\section{Closed points, divisors and skyscraper cohomology}
\label{pa-chap:RiemannRoch}

The valuation-theoretic organization of this chapter is inspired by Papaioannou
\cite{papaioannou-rr}; scheme-theoretic local facts are also compared with Hartshorne
\cite{hartshorne-ag}.

This chapter assembles the local and divisorial foundations of the Riemann--Roch ledger of a
smooth curve. On an integral scheme \(X\) smooth of relative dimension one over a field \(K\),
every point other than the generic point is closed, and the stalk of the structure sheaf there
is a discrete valuation ring; this yields, on the function field \(K(X)\), the order-of-vanishing
valuation \(\mathrm{ord}_x\) at each closed point \(x\). Each closed point also carries a residue
degree \([\kappa(x) : K]\), the \(K\)-dimension of its residue field. A Weil divisor is then a
finitely supported \(\Z\)-family on the closed points, and its degree is the residue-degree
weighted sum of its coefficients --- the correct notion over a non-closed base field. Finally, to
each \(K\)-module \(M\) and point \(x\) is attached the skyscraper sheaf of \(K\)-modules
concentrated at \(x\); its degree-zero cohomology recovers \(M\) and its degree-one cohomology
vanishes, the keystone input to the local contribution of a point in the cohomological
Riemann--Roch computation.

Throughout, \(K\) is a field, \(X\) is an integral scheme smooth of relative dimension one over
\(\Spec K\) --- the curve bundle of Chapter~\ref{pa-chap:Curves} --- and \(\eta\) denotes its generic
point. We write \(\struct{X,x}\) for the stalk at \(x\), \(\kappa(x)\) for its residue field,
and \(K(X)\) for the function field (the stalk at \(\eta\)).

\section{Discrete valuations at closed points}

\begin{theorem}[Discrete valuation ring at a closed point of an affine Dedekind chart]
  \label{pa-thm:affine_dedekind_dvr_stalk}
  \lean{AlgebraicGeometry.IsAffineOpen.isDiscreteValuationRing_stalk}
  \leanok
  \dcref{custom}

  Let \(X\) be an integral scheme, \(V \subseteq X\) an affine open whose ring of sections
  \(\Gamma(X, V)\) is a Dedekind domain, and \(x \in V\) a point which is not the generic point
  of \(X\). Then the stalk \(\struct{X,x}\) is a discrete valuation ring.
\end{theorem}
\begin{proof}
  \leanok
  Under the isomorphism \(V \cong \Spec \Gamma(X, V)\) the point \(x\) corresponds to a prime
  ideal \(\mathfrak q\) of \(\Gamma(X, V)\), and the stalk \(\struct{X,x}\) is the localization
  of \(\Gamma(X, V)\) at \(\mathfrak q\). The prime \(\mathfrak q\) is nonzero: if
  \(\mathfrak q = 0\) then \(x\) is the image of the generic point of \(\Spec \Gamma(X, V)\),
  which is the generic point of the nonempty open \(V\), hence the generic point of the
  irreducible scheme \(X\) --- contrary to hypothesis. The localization of a Dedekind domain at
  a nonzero prime ideal is a discrete valuation ring, which is the claim.
\end{proof}

\begin{theorem}[Discrete valuation ring at a closed point of a curve]
  \label{pa-thm:closedpoint_dvr_stalk}
  \lean{AlgebraicGeometry.isDiscreteValuationRing_stalk}
  \leanok
  \dcref{custom}

  \uses{pa-thm:affine_dedekind_dvr_stalk, pa-thm:relDim1_dedekind_charts}
  Let \(X\) be an integral scheme smooth of relative dimension one over a field \(K\), and let
  \(x \in X\) be a point other than the generic point. Then the stalk \(\struct{X,x}\) is a
  discrete valuation ring.
\end{theorem}
\begin{proof}
  \leanok
  By \ref{pa-thm:relDim1_dedekind_charts} the point \(x\) lies in an affine open \(V\) whose ring
  of sections \(\Gamma(X, V)\) is a Dedekind domain. Apply \ref{pa-thm:affine_dedekind_dvr_stalk}.
\end{proof}

\begin{theorem}[Dedekind stalk at a closed point of a curve]
  \label{pa-thm:closedpoint_dedekind_stalk}
  \lean{AlgebraicGeometry.isDedekindDomain_stalk}
  \leanok
  \dcref{custom}

  \uses{pa-thm:closedpoint_dvr_stalk}
  In the situation of \ref{pa-thm:closedpoint_dvr_stalk}, the stalk \(\struct{X,x}\) is a Dedekind
  domain.
\end{theorem}
\begin{proof}
  \leanok
  A discrete valuation ring is a local principal ideal domain, in particular a Dedekind domain;
  apply this to the stalk, which is a discrete valuation ring by \ref{pa-thm:closedpoint_dvr_stalk}.
\end{proof}

\begin{theorem}[Non-generic points of a curve are closed]
  \label{pa-thm:closedpoint_isClosed}
  \lean{AlgebraicGeometry.isClosed_singleton_of_ne_genericPoint}
  \leanok
  \dcref{custom}

  \uses{pa-thm:relDim1_specializes, pa-lem:height_one_closed_points}
  Let \(X\) be an integral scheme smooth of relative dimension one over a field \(K\), and let
  \(x \in X\) be a point other than the generic point. Then the singleton \(\{x\}\) is closed.
\end{theorem}
\begin{proof}
  \leanok
  By \ref{pa-thm:relDim1_specializes} the specialization order of \(X\) has height one: every
  generalization of any point is the generic point or the point itself. A scheme with
  irreducible underlying space and height-one specialization order has all its non-generic
  points closed (\ref{pa-lem:height_one_closed_points}); apply this to \(x\).
\end{proof}

\begin{definition}[The order valuation at a closed point]
  \label{pa-def:scheme_ord}
  \lean{AlgebraicGeometry.Scheme.ord}
  \leanok
  \source{papaioannou-algebraic-rr:page-0003}

  \uses{pa-thm:closedpoint_dvr_stalk, pa-thm:closedpoint_dedekind_stalk}
  \dcref{custom}
  For a closed point \(x\) of \(X\) (a point other than the generic point), the
  \emph{order valuation} \(\mathrm{ord}_x\) is the discrete valuation of the function field
  \(K(X)\) with values in \(\Z \cup \{+\infty\}\) determined by the discrete valuation ring
  \(\struct{X,x}\): it assigns to a nonzero rational function its order of vanishing at \(x\),
  and \(+\infty\) to \(0\).
\end{definition}
\begin{proof}
  \leanok
  The stalk \(\struct{X,x}\) is a discrete valuation ring (\ref{pa-thm:closedpoint_dvr_stalk}),
  hence a Dedekind domain (\ref{pa-thm:closedpoint_dedekind_stalk}) which is not a field, so its
  maximal ideal is a nonzero prime, i.e.\ a height-one prime of \(\struct{X,x}\). On an integral
  scheme the function field \(K(X)\) is the fraction field of every stalk, in particular of
  \(\struct{X,x}\); and a height-one prime of a Dedekind domain \(B\) determines a discrete
  valuation on \(\Frac B\) --- the \(\mathfrak p\)-adic valuation counting the multiplicity of
  \(\mathfrak p\) in the factorization of a fractional ideal. Taking \(B = \struct{X,x}\) and
  \(\mathfrak p\) its maximal ideal gives the order valuation on \(K(X)\).
\end{proof}

\section{The residue degree}

The curve \(X\) is a scheme over \(\Spec K\); as in Chapter~\ref{pa-chap:Cohomology} the base field
acts on the sections through the structure morphism, and the same discipline supplies the module
structure on each residue field.

\begin{definition}[The base-field algebra map on a residue field]
  \label{pa-def:residueOverAlgebraMap}
  \lean{AlgebraicGeometry.Scheme.residueOverAlgebraMap}
  \leanok
  \dcref{custom}

  \uses{pa-def:overAlgebraMap}
  For a scheme \(X\) over \(\Spec K\) and a point \(x \in X\), the canonical ring homomorphism
  \[
    K \longrightarrow \kappa(x)
  \]
  through the structure morphism: the composite of the global algebra map
  \(K \to \Gamma(X, \struct X)\) of \ref{pa-def:overAlgebraMap}, the germ
  \(\Gamma(X, \struct X) \to \struct{X,x}\) at \(x\), and the residue map
  \(\struct{X,x} \to \kappa(x)\). By restriction of scalars along it, the residue field
  \(\kappa(x)\) is a \(K\)-module (indeed a \(K\)-algebra); as with the module structure on the
  structure sheaf, this is used as an explicit map rather than a global instance.
\end{definition}
\begin{proof}
  \leanok
  Each of the three maps is a ring homomorphism, and \(K \to \Gamma(X, \struct X)\) is the base
  algebra map of \ref{pa-def:overAlgebraMap}; their composite is the stated ring homomorphism.
\end{proof}

\begin{definition}[The residue degree]
  \label{pa-def:residueDeg}
  \lean{AlgebraicGeometry.Scheme.residueDeg}
  \leanok
  \source{papaioannou-algebraic-rr:page-0004}

  \uses{pa-def:residueOverAlgebraMap}
  \dcref{custom}
  For a scheme \(X\) over \(\Spec K\) and a point \(x \in X\), the \emph{residue degree} at
  \(x\) is
  \[
    \deg_K(x) \;=\; [\kappa(x) : K] \;=\; \dim_K \kappa(x),
  \]
  the \(K\)-dimension of the residue field, taken with the \(K\)-module structure of
  \ref{pa-def:residueOverAlgebraMap}.
\end{definition}

\section{Weil divisors and their degree}

A Weil divisor on the curve records a finite prescription of orders at closed points. Its degree
must weight each point by its residue degree over the base field: this is the correct notion when
\(K\) is not algebraically closed. For example, on \(\PP^1_{\mathbb{R}}\) the rational function
\(t^2 + 1\) has a single closed zero \(x\), whose residue field is \(\mathbb{C}\), and a double
pole at \(\infty\), with residue field \(\mathbb{R}\); the unweighted point count \(1 - 2 = -1\)
is nonzero, whereas the residue-degree weighted degree
\(1 \cdot [\mathbb{C} : \mathbb{R}] - 2 \cdot [\mathbb{R} : \mathbb{R}] = 2 - 2 = 0\)
vanishes, as the degree of a principal divisor must.

\begin{definition}[Weil divisors on a curve]
  \label{pa-def:curveDivisor}
  \lean{AlgebraicGeometry.Scheme.CurveDivisor}
  \leanok
  \source{papaioannou-algebraic-rr:page-0006}

  \uses{pa-thm:closedpoint_isClosed}
  \dcref{custom}
  For an integral scheme \(X\) with generic point \(\eta\), a \emph{Weil divisor} is a finitely
  supported \(\Z\)-valued function on the non-generic points, i.e.\ an element of the free abelian
  group
  \[
    \Div(X) \;=\; \bigl(\{ x \in X : x \neq \eta \} \to_{\mathrm{fin}} \Z\bigr)
  \]
  on the set of non-generic points. It carries the pointwise abelian-group structure and the
  pointwise partial order, both independent of any base field. For the curve \(X\) the
  non-generic points are exactly the closed points (\ref{pa-thm:closedpoint_isClosed}), so a Weil
  divisor is a finite formal \(\Z\)-combination of closed points.
\end{definition}

\begin{definition}[The degree of a Weil divisor]
  \label{pa-def:curveDivisor_deg}
  \lean{AlgebraicGeometry.Scheme.CurveDivisor.deg}
  \leanok
  \source{papaioannou-algebraic-rr:page-0007}

  \uses{pa-def:curveDivisor, pa-def:residueDeg}
  \dcref{custom}
  For a scheme \(X\) over \(\Spec K\) and a Weil divisor \(D \in \Div(X)\), the \emph{degree} of
  \(D\), weighted by the residue degrees over the base field, is
  \[
    \deg_K D \;=\; \sum_{x} D_x \cdot [\kappa(x) : K] \;\in\; \Z ,
  \]
  the finite sum over the support of \(D\) of the coefficients \(D_x\) times the residue degrees
  \([\kappa(x) : K]\) of \ref{pa-def:residueDeg}.
\end{definition}

\begin{theorem}[The degree is a group homomorphism]
  \label{pa-thm:curveDivisor_deg_hom}
  \lean{AlgebraicGeometry.Scheme.CurveDivisor.deg_zero,
    AlgebraicGeometry.Scheme.CurveDivisor.deg_add,
    AlgebraicGeometry.Scheme.CurveDivisor.deg_neg}
  \leanok
  \source{papaioannou-algebraic-rr:page-0007}

  \uses{pa-def:curveDivisor_deg}
  \dcref{custom}
  The degree map \(\deg_K : \Div(X) \to \Z\) is a homomorphism of abelian groups: it satisfies
  \(\deg_K 0 = 0\), \(\deg_K(D + D') = \deg_K D + \deg_K D'\) and \(\deg_K(-D) = -\deg_K D\).
\end{theorem}
\begin{proof}
  \leanok
  The degree is the finitely supported sum \(D \mapsto \sum_x D_x \cdot [\kappa(x) : K]\). The
  zero divisor has empty support, so \(\deg_K 0 = 0\). For additivity, the summand
  \(n \mapsto n \cdot [\kappa(x) : K]\) is, at each point \(x\), additive in the coefficient
  \(n\) and vanishes at \(n = 0\); hence the finitely supported sum of the pointwise sums
  \((D + D')_x = D_x + D'_x\) splits as \(\sum_x D_x[\kappa(x):K] + \sum_x D'_x[\kappa(x):K]\),
  giving \(\deg_K(D + D') = \deg_K D + \deg_K D'\). Finally, applying additivity to \(D\) and
  \(-D\) and using \(\deg_K 0 = 0\) yields \(\deg_K(-D) = -\deg_K D\).
\end{proof}

\begin{lemma}[Degree of a single-point divisor]
  \label{pa-lem:curveDivisor_deg_single}
  \lean{AlgebraicGeometry.Scheme.CurveDivisor.deg_single}
  \leanok
  \dcref{custom}

  \uses{pa-def:curveDivisor_deg}
  For a closed point \(x\) and an integer \(n\), the degree of the single-point divisor
  \(n \cdot x\) is
  \[
    \deg_K(n \cdot x) \;=\; n \cdot [\kappa(x) : K] .
  \]
\end{lemma}
\begin{proof}
  \leanok
  The divisor \(n \cdot x\) is supported at the single point \(x\) with coefficient \(n\), so the
  defining sum of \ref{pa-def:curveDivisor_deg} has the single term \(n \cdot [\kappa(x) : K]\)
  (the summand vanishes at coefficient \(0\), so no other point contributes).
\end{proof}

\section{Skyscraper sheaves and the vanishing of \(H^1\)}

Fix a scheme \(X\) over \(\Spec K\) and work on its small Zariski site (the poset of opens with
open coverings), whose terminal object is the whole space \(\top = X\). For a point \(x \in X\)
and a \(K\)-module \(M\) the skyscraper sheaf concentrates \(M\) at \(x\).

\begin{definition}[The skyscraper sheaf of \(K\)-modules]
  \label{pa-def:skyModule}
  \lean{AlgebraicGeometry.skyModule}
  \leanok
  \dcref{custom}

  For \(x \in X\) and a \(K\)-module \(M\), the \emph{skyscraper sheaf} \(\mathrm{sky}_x(M)\) of
  \(K\)-modules on the small Zariski site of \(X\) is the sheaf whose sections over an open
  \(U\) are
  \[
    \mathrm{sky}_x(M)(U) \;=\;
      \begin{cases} M & x \in U, \\ 0 & x \notin U, \end{cases}
  \]
  with restriction between two opens containing \(x\) the identity of \(M\), and to an open
  avoiding \(x\) the zero map. Its restriction map between any two opens that both contain \(x\)
  is an isomorphism, and its sections over any open avoiding \(x\) are the zero module.
\end{definition}

\begin{definition}[Degree-zero cohomology of a skyscraper is its value]
  \label{pa-def:skyModuleGammaEquiv}
  \lean{AlgebraicGeometry.skyModuleGammaEquiv}
  \leanok
  \dcref{custom}

  \uses{pa-def:skyModule, pa-def:HModule_linearEquiv0}
  For \(x \in X\) and a \(K\)-module \(M\), there is a \(K\)-linear equivalence
  \[
    H^0\bigl(X, \mathrm{sky}_x(M)\bigr) \;\cong\; M .
  \]
\end{definition}
\begin{proof}
  \leanok
  Global sections live over the terminal open \(\top = X\), which contains \(x\), so
  \(\mathrm{sky}_x(M)(\top) = M\) by \ref{pa-def:skyModule}. Compose with the identification of
  degree-zero cohomology with global sections, \(H^0(F) \cong F(\top)\) of
  \ref{pa-def:HModule_linearEquiv0}, applied to \(F = \mathrm{sky}_x(M)\).
\end{proof}

\begin{theorem}[Global lifting along the cokernel of a skyscraper embedding]
  \label{pa-thm:sky_cokernel_app_top_surjective}
  \lean{AlgebraicGeometry.cokernelπ_app_top_surjective}
  \leanok
  \dcref{custom}

  \uses{pa-def:skyModule, pa-lem:sections_left_exact, pa-lem:app_injective_of_mono}
  Let \(x \in X\), let \(M\) be a \(K\)-module, and let \(\iota : \mathrm{sky}_x(M) \to G\) be a
  monomorphism of sheaves of \(K\)-modules. Then every global section of \(\coker \iota\) lifts
  to a global section of \(G\): the map \(G(\top) \to (\coker \iota)(\top)\) is surjective.
\end{theorem}
\begin{proof}
  \leanok
  Write \(\pi : G \to \coker \iota\) for the projection and let \(q \in (\coker \iota)(\top)\).
  The proof glues corrected local lifts, the correction being an \emph{elementwise} coboundary
  supplied by flasqueness of the skyscraper --- so no quasi-compactness or finite subcover is
  needed. Restriction maps of a sheaf \(F\) are written \(t|_W\).

  (a) \emph{Local lifts.} As an epimorphism of sheaves, \(\pi\) is locally surjective on
  sections: for each point \(p \in X\) there is an open \(W_p \ni p\) and a section
  \(s_p \in G(W_p)\) with \(\pi(s_p) = q|_{W_p}\). The opens \(W_p\) cover \(X\), since each
  contains its point \(p\).

  (b) \emph{The differences come from the skyscraper.} On \(W_p \cap W_{p'}\) the difference
  \(s_p|_{W_p \cap W_{p'}} - s_{p'}|_{W_p \cap W_{p'}}\) maps to
  \(q|_{W_p \cap W_{p'}} - q|_{W_p \cap W_{p'}} = 0\) under \(\pi\). Since \(\iota\) is a
  monomorphism it is the kernel of its cokernel on sections (\ref{pa-lem:sections_left_exact}), so
  there is a unique \(a_{pp'} \in \mathrm{sky}_x(M)(W_p \cap W_{p'})\) with
  \(\iota(a_{pp'}) = s_p - s_{p'}\) there.

  (c) \emph{Cocycle identity.} On a triple overlap \(W_p \cap W_{p'} \cap W_l\),
  \[
    \iota(a_{pl}) = s_p - s_l = (s_p - s_{p'}) + (s_{p'} - s_l) = \iota(a_{pp'} + a_{p'l}) ,
  \]
  and \(\iota\) is injective on sections (\ref{pa-lem:app_injective_of_mono}), so
  \(a_{pl} = a_{pp'} + a_{p'l}\) there.

  (d) \emph{Pivot on the support point.} Because \(\mathrm{sky}_x(M)\) is a skyscraper sheaf it
  is \emph{flasque}: every restriction map of \(\mathrm{sky}_x(M)\) is surjective. In particular
  the restriction \(\mathrm{sky}_x(M)(W_p) \to \mathrm{sky}_x(M)(W_p \cap W_x)\) is surjective,
  so choose \(b_p \in \mathrm{sky}_x(M)(W_p)\) with \(b_p|_{W_p \cap W_x} = a_{px}\), where the
  index \(x\)-th cover member \(W_x \ni x\) is the neighbourhood of the support point itself.

  (e) \emph{Restriction to \(W_x\) is injective.} For every open \(A\), the restriction
  \(\mathrm{sky}_x(M)(A) \to \mathrm{sky}_x(M)(A \cap W_x)\) is injective: if \(x \in A\) then
  both opens contain \(x\) and the restriction is an isomorphism (\ref{pa-def:skyModule}); if
  \(x \notin A\) then \(\mathrm{sky}_x(M)(A) = 0\) is trivial, so any map out of it is injective.

  (f) \emph{The cocycle is the coboundary of \(b\).} We claim
  \(a_{pp'} = b_p|_{W_p \cap W_{p'}} - b_{p'}|_{W_p \cap W_{p'}}\). By (e) it suffices to check
  this after restriction to \((W_p \cap W_{p'}) \cap W_x\). There the cocycle identity (c) with
  \(l = x\) gives \(a_{pp'} = a_{px} - a_{p'x}\), and \(a_{px} = b_p\), \(a_{p'x} = b_{p'}\)
  after restriction to \(\cap W_x\) by the choice in (d); hence
  \(a_{pp'} = b_p - b_{p'}\) restricted to \((W_p \cap W_{p'}) \cap W_x\), as required.

  (g) \emph{Corrected local sections glue.} Put \(t_p = s_p - \iota(b_p) \in G(W_p)\). On
  \(W_p \cap W_{p'}\),
  \[
    t_p - t_{p'} = (s_p - s_{p'}) - \iota(b_p - b_{p'}) = \iota(a_{pp'}) - \iota(a_{pp'}) = 0 ,
  \]
  using (b) and (f). So the \(t_p\) agree on overlaps and, by the sheaf condition, glue to a
  global section \(s \in G(\top)\) with \(s|_{W_p} = t_p\).

  (h) \emph{The glued section lifts \(q\).} On each \(W_p\), since \(\pi \circ \iota = 0\),
  \[
    \pi(s)|_{W_p} = \pi(t_p) = \pi(s_p) - \pi(\iota(b_p)) = q|_{W_p} .
  \]
  Sections of the sheaf \(\coker \iota\) that agree on the cover \(\{W_p\}\) are equal, so
  \(\pi(s) = q\). Thus \(s\) is the required lift.
\end{proof}

\begin{theorem}[Degree-one cohomology of a skyscraper vanishes]
  \label{pa-thm:skyModule_h1_vanishing}
  \lean{AlgebraicGeometry.skyModule_subsingleton_hModule_one}
  \leanok
  \dcref{custom}

  \uses{pa-def:skyModule, pa-def:HModule}
  For every point \(x \in X\) and every \(K\)-module \(M\),
  \[
    H^1\bigl(X, \mathrm{sky}_x(M)\bigr) \;=\; 0 .
  \]
\end{theorem}
\begin{proof}
  \leanok
  \uses{pa-lem:ext_one_vanishing_criterion, pa-thm:sky_cokernel_app_top_surjective, pa-def:constHomEquiv,
    pa-lem:constHomEquiv_naturality}
  The category of sheaves of \(K\)-modules on the small Zariski site of \(X\) is Grothendieck
  abelian, hence has enough injectives: embed \(\mathrm{sky}_x(M)\) into an injective sheaf by a
  monomorphism \(\iota : \mathrm{sky}_x(M) \to G\). By the \(\Ext^1\)-vanishing criterion
  \ref{pa-lem:ext_one_vanishing_criterion} (with \(L = K_J\) the constant sheaf), it suffices to
  lift every morphism \(K_J \to \coker \iota\) along \(\pi : G \to \coker \iota\). By
  \ref{pa-def:constHomEquiv} and its naturality \ref{pa-lem:constHomEquiv_naturality}, morphisms out
  of \(K_J\) are global sections, compatibly with postcomposition; so the required lifting is
  exactly the surjectivity of \(G(\top) \to (\coker \iota)(\top)\), which is
  \ref{pa-thm:sky_cokernel_app_top_surjective}.
\end{proof}

\section{Finiteness of the residue degree}

Over an arbitrary base field the residue field at a closed point of the curve is a finite extension
of \(K\); this finiteness is the arithmetic gate through which every point contributes a finite
local term to the ledger.

\begin{theorem}[Finiteness of the residue degree]
  \label{pa-thm:residueDeg_finite}
  \lean{AlgebraicGeometry.Scheme.residueDeg_finite}
  \leanok
  \source{papaioannou-algebraic-rr:page-0004}

  \uses{pa-thm:relDim1_dedekind_charts, pa-thm:affine_dedekind_dvr_stalk, pa-def:residueOverAlgebraMap}
  \dcref{custom}
  Let \(X\) be an integral scheme smooth of relative dimension one and locally of finite type over a
  field \(K\), and let \(x\) be a non-generic (hence closed) point. Then the residue field
  \(\kappa(x)\) is a finite \(K\)-module; equivalently the residue degree \([\kappa(x):K]\) is
  finite.
\end{theorem}
\begin{proof}
  \leanok
  By \ref{pa-thm:relDim1_dedekind_charts} the point \(x\) lies in an affine open \(V\) whose ring of
  sections \(B = \Gamma(X, V)\) is a Dedekind domain. Under \(V \cong \Spec B\) the point \(x\)
  corresponds to a prime \(\mathfrak q\) of \(B\), which is nonzero --- otherwise \(x\) would be the
  generic point, as in \ref{pa-thm:affine_dedekind_dvr_stalk} --- and hence maximal, so \(B/\mathfrak q\)
  is a field. The stalk \(\struct{X,x}\) is the localization \(B_{\mathfrak q}\), whose residue
  field is canonically \(B/\mathfrak q\); this identification is \(K\)-linear for the residue
  \(K\)-algebra structure of \ref{pa-def:residueOverAlgebraMap}, so \(\kappa(x) \cong B/\mathfrak q\)
  as \(K\)-modules. Now \(K \to B\) is of finite type (\(X\) is locally of finite type over \(K\),
  restricted to the chart \(V\)), hence so is the quotient \(K \to B/\mathfrak q\); a finite-type
  algebra over a field which is itself a field is a finite extension (Zariski's lemma, valid because
  a field is a Jacobson ring). Therefore \(B/\mathfrak q \cong \kappa(x)\) is finite over \(K\).
\end{proof}

\begin{theorem}[Positivity of the residue degree]
  \label{pa-thm:residueDeg_pos}
  \lean{AlgebraicGeometry.Scheme.residueDeg_pos}
  \leanok
  \dcref{custom}

  \uses{pa-thm:residueDeg_finite, pa-def:residueDeg}
  In the situation of \ref{pa-thm:residueDeg_finite}, the residue degree is positive:
  \(0 < [\kappa(x):K]\).
\end{theorem}
\begin{proof}
  \leanok
  By \ref{pa-thm:residueDeg_finite} the residue field \(\kappa(x)\) is a finite \(K\)-module; being a
  field extension of \(K\) it is nonzero and free over \(K\). A nonzero finite free module over a
  field has strictly positive rank, so \([\kappa(x):K] = \dim_K \kappa(x) > 0\).
\end{proof}

\section{Principal divisors}

A nonzero rational function on the curve has a principal divisor recording its zeros and poles with
multiplicities. The construction rests on the \(\Z\)-valued order at each closed point and on the
finiteness of the set of points where a given function has a zero or a pole.

\begin{definition}[The order homomorphism on function-field units]
  \label{pa-def:scheme_ordZ}
  \lean{AlgebraicGeometry.Scheme.ordZ}
  \leanok
  \source{papaioannou-algebraic-rr:page-0005}

  \uses{pa-def:scheme_ord}
  \dcref{custom}
  For a closed point \(x\), the \emph{order homomorphism}
  \[
    \mathrm{ord}_x : K(X)^\times \longrightarrow \Z
  \]
  sends a nonzero rational function to its order of vanishing at \(x\) (\ref{pa-def:scheme_ord}):
  positive at a zero, negative at a pole, with a uniformizer of the discrete valuation ring
  \(\struct{X,x}\) having order \(+1\). It is a homomorphism from the multiplicative group
  \(K(X)^\times\) to the additive group \(\Z\).
\end{definition}
\begin{proof}
  \leanok
  The order valuation of \ref{pa-def:scheme_ord} is a valuation of \(K(X)\), hence multiplicative on
  nonzero elements; restricting it to units and reading its value in \(\Z\) with the classical sign
  (a uniformizer, whose valuation generates the value group, is assigned \(+1\)) yields a group
  homomorphism \(K(X)^\times \to \Z\).
\end{proof}

\begin{lemma}[Trivial order where a section is invertible]
  \label{pa-lem:ord_eq_one_of_mem_basicOpen}
  \lean{AlgebraicGeometry.Scheme.ord_eq_one_of_mem_basicOpen}
  \leanok
  \dcref{custom}

  \uses{pa-def:scheme_ord, pa-thm:closedpoint_dvr_stalk}
  Let \(s \in \Gamma(X, U)\) represent the rational function \(g = \mathrm{germ}_\eta(s)\) at the
  generic point (with \(\eta \in U\)), and let \(x\) be a closed point of the basic open set where
  \(s\) is invertible, i.e.\ the germ of \(s\) at \(x\) is a unit of \(\struct{X,x}\). Then
  \(\mathrm{ord}_x g = 0\).
\end{lemma}
\begin{proof}
  \leanok
  Since \(\eta\) specializes to \(x\), the function \(g\) is the image of the germ \(\mathrm{germ}_x(s)\)
  under the localization map \(\struct{X,x} \to K(X)\). By hypothesis this germ is a unit of the
  discrete valuation ring \(\struct{X,x}\) (\ref{pa-thm:closedpoint_dvr_stalk}), and a unit of the
  valuation ring has order of vanishing \(0\); hence \(\mathrm{ord}_x g = 0\).
\end{proof}

\begin{theorem}[Finiteness of the order support]
  \label{pa-thm:ordZ_support_finite}
  \lean{AlgebraicGeometry.Scheme.ordZ_support_finite}
  \leanok
  \source{papaioannou-algebraic-rr:page-0006}

  \uses{pa-lem:ord_eq_one_of_mem_basicOpen, pa-thm:height_one_finite_closed, pa-thm:relDim1_specializes}
  \dcref{custom}
  For \(g \in K(X)^\times\), the set of closed points at which \(g\) has nonzero order,
  \(\{ x : \mathrm{ord}_x g \neq 0 \}\), is finite. (Here \(X\) is additionally assumed
  quasi-compact over \(K\), so that its underlying space is Noetherian.)
\end{theorem}
\begin{proof}
  \leanok
  Represent \(g\) by a section \(s\) on a nonempty open \(U \ni \eta\) (the generic point lies in
  every nonempty open). Then \(g\) is a unit at \(\eta\), so \(\eta\) lies in the basic open set
  \(D(s)\) where \(s\) is invertible. At every closed point of \(D(s)\) the order of \(g\) is \(0\)
  (\ref{pa-lem:ord_eq_one_of_mem_basicOpen}), so the locus where \(\mathrm{ord}_x g \neq 0\) is
  contained in the complement \(X \setminus D(s)\), a closed set not containing \(\eta\). On the
  curve --- an irreducible Noetherian space whose specialization order has height one
  (\ref{pa-thm:relDim1_specializes}) --- a proper closed set avoiding the generic point is finite
  (\ref{pa-thm:height_one_finite_closed}). Quasi-compactness over \(K\) supplies the Noetherian
  hypothesis: finite type alone gives only a locally Noetherian space.
\end{proof}

\begin{definition}[The principal divisor of a rational function]
  \label{pa-def:divOf}
  \lean{AlgebraicGeometry.Scheme.divOf}
  \leanok
  \source{papaioannou-algebraic-rr:page-0007}

  \uses{pa-def:curveDivisor, pa-def:scheme_ordZ, pa-thm:ordZ_support_finite}
  \dcref{custom}
  For \(g \in K(X)^\times\), the \emph{principal divisor}
  \[
    \mathrm{div}(g) = \sum_{x} \mathrm{ord}_x(g)\, x
  \]
  is the Weil divisor whose coefficient at a closed point \(x\) is the order of vanishing
  \(\mathrm{ord}_x(g)\): positive at zeros, negative at poles. It is finitely supported
  (\ref{pa-thm:ordZ_support_finite}), hence a well-defined element of \(\Div(X)\).
\end{definition}

\begin{theorem}[The principal divisor is a homomorphism]
  \label{pa-thm:divOf_hom}
  \lean{AlgebraicGeometry.Scheme.divOf_mul, AlgebraicGeometry.Scheme.divOf_one}
  \leanok
  \source{papaioannou-algebraic-rr:page-0007}

  \uses{pa-def:divOf, pa-def:scheme_ordZ}
  \dcref{custom}
  The principal divisor is a homomorphism from \(K(X)^\times\) to \(\Div(X)\): for all
  \(g, g' \in K(X)^\times\),
  \[
    \mathrm{div}(g g') = \mathrm{div}(g) + \mathrm{div}(g'), \qquad \mathrm{div}(1) = 0 .
  \]
\end{theorem}
\begin{proof}
  \leanok
  At each closed point \(x\) the order \(\mathrm{ord}_x\) is a homomorphism \(K(X)^\times \to \Z\)
  (\ref{pa-def:scheme_ordZ}), so \(\mathrm{ord}_x(g g') = \mathrm{ord}_x g + \mathrm{ord}_x g'\) and
  \(\mathrm{ord}_x(1) = 0\). Since the coefficients of \(\mathrm{div}(g)\) are these orders
  (\ref{pa-def:divOf}), the two identities hold coefficientwise, hence as divisors.
\end{proof}

\section{The divisor sheaf of a Weil divisor}

To a Weil divisor \(D\) is attached the sheaf \(\mathcal O(D)\) of rational functions with poles
bounded by \(D\). Its sections over an open \(U\) are the classical Riemann--Roch space of \(D\)
restricted to \(U\); the sheaf structure is available because on the irreducible curve a compatible
family of such functions shares a single underlying rational function.

\begin{definition}[The \(K\)-structure on the function field]
  \label{pa-def:functionFieldOverAlgebraMap}
  \lean{AlgebraicGeometry.Scheme.functionFieldOverAlgebraMap}
  \leanok
  \dcref{custom}

  \uses{pa-def:overAlgebraMap}
  For \(X\) over \(\Spec K\), the function field \(K(X)\) carries a canonical \(K\)-algebra structure
  through the structure morphism: the composite
  \[
    K \longrightarrow \Gamma(X, \struct X) \longrightarrow K(X)
  \]
  of the global algebra map \ref{pa-def:overAlgebraMap} with the germ at the generic point. By
  restriction of scalars along it, \(K(X)\) is a \(K\)-module (as with the residue field, this is
  used as an explicit map rather than a global instance).
\end{definition}

\begin{lemma}[Integral elements have no pole]
  \label{pa-lem:ord_algebraMap_stalk_le_one}
  \lean{AlgebraicGeometry.Scheme.ord_algebraMap_stalk_le_one}
  \leanok
  \dcref{custom}

  \uses{pa-thm:closedpoint_dvr_stalk, pa-def:scheme_ord}
  The image in \(K(X)\) of any element of the stalk \(\struct{X,x}\) at a closed point \(x\) has no
  pole at \(x\): its order of vanishing is \(\ge 0\).
\end{lemma}
\begin{proof}
  \leanok
  The stalk \(\struct{X,x}\) is a discrete valuation ring (\ref{pa-thm:closedpoint_dvr_stalk}), and it
  is exactly the valuation ring of \(\mathrm{ord}_x\) --- the order valuation is the adic valuation
  of its maximal ideal (\ref{pa-def:scheme_ord}). Every element of a valuation ring has nonnegative
  order, so the image of any stalk element has \(\mathrm{ord}_x \ge 0\).
\end{proof}

\begin{lemma}[Constants have no pole]
  \label{pa-lem:ord_functionFieldOverAlgebraMap_le_one}
  \lean{AlgebraicGeometry.Scheme.ord_functionFieldOverAlgebraMap_le_one}
  \leanok
  \dcref{custom}

  \uses{pa-lem:ord_algebraMap_stalk_le_one, pa-def:functionFieldOverAlgebraMap}
  The image in \(K(X)\) of a constant \(r \in K\) has no pole at any closed point \(x\):
  \(\mathrm{ord}_x(r) \ge 0\).
\end{lemma}
\begin{proof}
  \leanok
  The image of \(r\) is the germ at the generic point of the global section \(r \cdot 1\)
  (\ref{pa-def:functionFieldOverAlgebraMap}); its germ at \(x\) realises it as an element of the stalk
  \(\struct{X,x}\), so it has no pole at \(x\) by \ref{pa-lem:ord_algebraMap_stalk_le_one}.
\end{proof}

\begin{definition}[Sections of the divisor sheaf]
  \label{pa-def:divisorSections}
  \lean{AlgebraicGeometry.Scheme.divisorSections}
  \leanok
  \source{papaioannou-algebraic-rr:page-0007}

  \uses{pa-def:scheme_ord, pa-lem:ord_functionFieldOverAlgebraMap_le_one, pa-def:curveDivisor,
    pa-def:functionFieldOverAlgebraMap}
  \dcref{custom}
  For a Weil divisor \(D\) and a nonempty open \(U\), the sections of \(\mathcal O(D)\) over \(U\)
  are
  \[
    \mathcal O(D)(U) = \{ g \in K(X) : \text{at every closed point } x \in U,\ g \text{ has pole
      order} \le D_x \},
  \]
  the rational functions with poles bounded by \(D\) on \(U\); over the empty open the sections are
  the zero module (forced by the sheaf condition, since the empty sieve covers \(\bot\)). This is a
  \(K\)-submodule of \(K(X)\): it contains \(0\); it is closed under addition because a valuation is
  subadditive, \(\mathrm{ord}_x(g_1 + g_2) \ge \min(\mathrm{ord}_x g_1, \mathrm{ord}_x g_2)\), so
  the pole bound is preserved; and closed under \(K\)-scaling because constants have no poles
  (\ref{pa-lem:ord_functionFieldOverAlgebraMap_le_one}).
\end{definition}

\begin{lemma}[Nonempty opens of an irreducible space meet]
  \label{pa-lem:opens_inf_nonempty}
  \lean{AlgebraicGeometry.Scheme.opens_inf_nonempty}
  \leanok
  \dcref{custom}

  Any two nonempty open subsets of the irreducible space \(X\) have nonempty intersection.
\end{lemma}
\begin{proof}
  \leanok
  The underlying space of the integral scheme \(X\) is irreducible, so no two nonempty opens are
  disjoint.
\end{proof}

\begin{definition}[The divisor sheaf \(\mathcal O(D)\)]
  \label{pa-def:divisorSheaf}
  \lean{AlgebraicGeometry.Scheme.divisorSheaf}
  \leanok
  \dcref{custom}

  \uses{pa-def:divisorSections, pa-lem:opens_inf_nonempty}
  The presheaf \(U \mapsto \mathcal O(D)(U)\) of \ref{pa-def:divisorSections} is a sheaf of
  \(K\)-modules on the small Zariski site of \(X\), denoted \(\mathcal O(D)\).
\end{definition}
\begin{proof}
  \leanok
  Restriction between nonempty opens is the inclusion of subspaces of \(K(X)\), so a section is
  determined by its underlying rational function. Given a compatible family over a cover of \(U\):
  any two members with nonempty domain agree as rational functions, since their domains meet
  (\ref{pa-lem:opens_inf_nonempty}) and they restrict equally on the overlap; hence the family shares a
  single \(g \in K(X)\). Every closed point of \(U\) lies in some cover member, where \(g\) respects
  the bound \(D\); so \(g \in \mathcal O(D)(U)\) and restricts to the given sections. Over the empty
  open all sections are \(0\). Uniqueness is again the determinacy of a section by its underlying
  rational function.
\end{proof}

\begin{theorem}[Monotone inclusion of divisor sheaves]
  \label{pa-thm:divisorSheafLE}
  \lean{AlgebraicGeometry.Scheme.divisorSheafLE, AlgebraicGeometry.Scheme.divisorSheafLE_mono}
  \leanok
  \dcref{custom}

  \uses{pa-def:divisorSheaf, pa-def:curveDivisor}
  For \(D \le D'\) the identity on rational functions gives a monomorphism of sheaves
  \(\mathcal O(D) \hookrightarrow \mathcal O(D')\).
\end{theorem}
\begin{proof}
  \leanok
  For \(D \le D'\) one has \(D_x \le D'_x\) at every closed point, so a pole of order \(\le D_x\) is
  a fortiori a pole of order \(\le D'_x\); hence \(\mathcal O(D)(U) \subseteq \mathcal O(D')(U)\) as
  submodules of \(K(X)\) for every \(U\). These section-wise inclusions are injective and natural in
  \(U\), so they assemble into a morphism of sheaves \(\mathcal O(D) \hookrightarrow \mathcal O(D')\)
  which, being injective on every section, is a monomorphism.
\end{proof}

\begin{lemma}[No pole means integral]
  \label{pa-lem:exists_stalk_of_ord_le_one}
  \lean{AlgebraicGeometry.Scheme.exists_stalk_of_ord_le_one}
  \leanok
  \dcref{custom}

  \uses{pa-thm:closedpoint_dvr_stalk, pa-thm:closedpoint_dedekind_stalk, pa-def:scheme_ord}
  If a rational function \(g\) has no pole at a closed point \(x\) (order of vanishing \(\ge 0\)),
  then \(g\) is integral at \(x\): it lies in the image of the localization map
  \(\struct{X,x} \to K(X)\).
\end{lemma}
\begin{proof}
  \leanok
  The stalk \(\struct{X,x}\) is a discrete valuation ring (\ref{pa-thm:closedpoint_dvr_stalk}), a
  Dedekind domain (\ref{pa-thm:closedpoint_dedekind_stalk}) with a single nonzero prime --- its
  maximal ideal --- whose adic valuation is \(\mathrm{ord}_x\) (\ref{pa-def:scheme_ord}). Since
  \(K(X)\) is the fraction field of \(\struct{X,x}\) and there is only this one valuation to test,
  the single condition \(\mathrm{ord}_x g \ge 0\) places \(g\) in the valuation ring
  \(\struct{X,x}\).
\end{proof}

\section{Comparison of \(\mathcal O(0)\) with the structure sheaf}

The divisor sheaf of the zero divisor is the structure sheaf: a rational function with no pole
anywhere is a global regular section, and conversely. This identification transports the Euler
characteristic of \(\mathcal O(0)\) to that of \(\struct X\).

\begin{lemma}[The germ/stalk seam]
  \label{pa-lem:germ_generic_eq_algebraMap_germ}
  \lean{AlgebraicGeometry.Scheme.germ_generic_eq_algebraMap_germ}
  \leanok
  \dcref{custom}

  For a section \(s \in \Gamma(X, U)\) over an open containing both the generic point \(\eta\) and a
  point \(x\), the germ at \(\eta\) equals the image of the germ at \(x\) under the localization map
  \(\struct{X,x} \to K(X)\).
\end{lemma}
\begin{proof}
  \leanok
  The generic point specializes to \(x\), and the germ maps of the structure sheaf are compatible
  with the specialization map between stalks; applying this compatibility to \(s\) gives the stated
  identity.
\end{proof}

\begin{lemma}[Regular germs have poles bounded by zero]
  \label{pa-lem:germGenericLinear_mem}
  \lean{AlgebraicGeometry.Scheme.germGenericLinear_mem}
  \leanok
  \dcref{custom}

  \uses{pa-lem:germ_generic_eq_algebraMap_germ, pa-lem:ord_algebraMap_stalk_le_one, pa-def:divisorSections}
  For a section \(s \in \Gamma(X, U)\) over a nonempty open \(U\), its germ at the generic point
  \(\mathrm{germ}_\eta(s) \in K(X)\) has poles bounded by the zero divisor on \(U\), i.e.\ lies in
  \(\mathcal O(0)(U)\).
\end{lemma}
\begin{proof}
  \leanok
  At every closed point \(x \in U\), \(\mathrm{germ}_\eta(s)\) is the image of the germ
  \(\mathrm{germ}_x(s) \in \struct{X,x}\) (\ref{pa-lem:germ_generic_eq_algebraMap_germ}), hence has no
  pole at \(x\) (\ref{pa-lem:ord_algebraMap_stalk_le_one}). Thus it satisfies the zero-divisor bound at
  every point of \(U\).
\end{proof}

\begin{lemma}[Local regular representatives]
  \label{pa-lem:exists_local_section}
  \lean{AlgebraicGeometry.Scheme.exists_local_section}
  \leanok
  \dcref{custom}

  \uses{pa-lem:exists_stalk_of_ord_le_one, pa-lem:germ_generic_eq_algebraMap_germ, pa-def:divisorSections}
  A rational function \(g \in \mathcal O(0)(U)\) is, near every point \(x \in U\), the germ at
  \(\eta\) of a regular section over a neighbourhood of \(x\) inside \(U\).
\end{lemma}
\begin{proof}
  \leanok
  At \(x = \eta\), the function \(g\) is directly represented by a regular section on some
  neighbourhood. At a closed point \(x\), the zero-divisor bound says \(g\) has no pole at \(x\), so
  \(g\) is the image of some \(y \in \struct{X,x}\) (\ref{pa-lem:exists_stalk_of_ord_le_one}). Represent
  \(y\) by a regular section \(s\) on a neighbourhood \(W \ni x\) inside \(U\); then
  \(\mathrm{germ}_\eta(s)\) is the image of \(\mathrm{germ}_x(s) = y\), which is \(g\)
  (\ref{pa-lem:germ_generic_eq_algebraMap_germ}).
\end{proof}

\begin{lemma}[Global representability]
  \label{pa-lem:exists_section_germ_eq}
  \lean{AlgebraicGeometry.Scheme.exists_section_germ_eq}
  \leanok
  \dcref{custom}

  \uses{pa-lem:exists_local_section}
  Every \(g \in \mathcal O(0)(U)\) over a nonempty open \(U\) is the germ at \(\eta\) of a section
  \(s \in \Gamma(X, U)\).
\end{lemma}
\begin{proof}
  \leanok
  Cover \(U\) by the neighbourhoods \(W\) of \ref{pa-lem:exists_local_section} carrying regular
  sections \(s_W\) with \(\mathrm{germ}_\eta(s_W) = g\). On overlaps the \(s_W\) have equal germ at
  \(\eta\) (which lies in every nonempty overlap of the irreducible space), and since \(X\) is
  integral a section is determined by its germ at \(\eta\), so the \(s_W\) agree there. They
  therefore glue to a section \(s \in \Gamma(X, U)\) with \(\mathrm{germ}_\eta(s) = g\).
\end{proof}

\begin{theorem}[\(\mathcal O(0) \cong \struct X\)]
  \label{pa-thm:divisorSheafZeroIso}
  \lean{AlgebraicGeometry.Scheme.divisorSheafZeroIso}
  \leanok
  \dcref{custom}

  \uses{pa-lem:germGenericLinear_mem, pa-lem:exists_section_germ_eq, pa-def:divisorSheaf, pa-def:moduleKSheaf}
  The divisor sheaf of the zero divisor is the structure sheaf of \(X\) as a sheaf of \(K\)-modules:
  \[
    \mathcal O(0) \;\cong\; \struct X .
  \]
\end{theorem}
\begin{proof}
  \leanok
  The section-wise map \(\Gamma(X, U) \to \mathcal O(0)(U)\), \(s \mapsto \mathrm{germ}_\eta(s)\), is
  well defined (\ref{pa-lem:germGenericLinear_mem}) and \(K\)-linear. On a nonempty open it is
  injective, because \(X\) is integral so a section is determined by its germ at \(\eta\), and
  surjective by \ref{pa-lem:exists_section_germ_eq}; over the empty open both sides are the zero module.
  Being bijective on every open, these maps assemble into an isomorphism of sheaves of \(K\)-modules.
  This is the identification through which \(\chi(\mathcal O(0))\) equals \(\chi(\struct X)\).
\end{proof}

\section{The dévissage complex at a closed point}

Removing a single closed point from a divisor changes the divisor sheaf by a skyscraper at that
point. The dévissage short complex makes this precise: it relates \(\mathcal O(D-x)\), \(\mathcal
O(D)\) and a skyscraper sheaf built from the one-point jump module, the associated-graded piece of
the pole-order filtration at \(x\).

\begin{definition}[The one-point lattice]
  \label{pa-def:pointLattice}
  \lean{AlgebraicGeometry.pointLattice}
  \leanok
  \dcref{custom}

  \uses{pa-def:scheme_ord, pa-lem:ord_functionFieldOverAlgebraMap_le_one}
  For a closed point \(x\) and \(n \in \Z\), the \emph{one-point lattice} \(L_n\) is the
  \(K\)-submodule of \(K(X)\) of rational functions whose pole order at \(x\) is at most \(n\)
  (equivalently \(\mathrm{ord}_x g \ge -n\)). It is a submodule by the same argument as
  \ref{pa-def:divisorSections} --- the valuation is subadditive and \(K\)-multiples create no poles
  (\ref{pa-lem:ord_functionFieldOverAlgebraMap_le_one}) --- and \(L_m \subseteq L_n\) whenever
  \(m \le n\).
\end{definition}

\begin{definition}[The jump module]
  \label{pa-def:jumpModule}
  \lean{AlgebraicGeometry.jumpModule}
  \leanok
  \dcref{custom}

  \uses{pa-def:pointLattice, pa-def:curveDivisor}
  The \emph{jump module} at \(x\) is the quotient
  \[
    J \;=\; L_{D_x} \big/ L_{D_x - 1},
  \]
  the \(K\)-space of functions with pole order \(\le D_x\) at \(x\) modulo those with pole order
  \(\le D_x - 1\) --- the associated-graded piece of the pole-order filtration at the single point
  \(x\).
\end{definition}

\begin{definition}[The projection to the jump module]
  \label{pa-def:jumpProj}
  \lean{AlgebraicGeometry.jumpProj}
  \leanok
  \dcref{custom}

  \uses{pa-def:jumpModule, pa-def:divisorSections, pa-def:pointLattice}
  For an open \(U \ni x\), the \emph{jump projection} \(\mathrm{proj} : \mathcal O(D)(U) \to J\),
  \(g \mapsto [g]\). It is well defined and \(K\)-linear: a section of \(\mathcal O(D)\) over \(U\)
  has pole order \(\le D_x\) at the point \(x \in U\), so it lies in \(L_{D_x}\) and determines a
  class in \(J\); the class depends only on the underlying rational function.
\end{definition}

\begin{definition}[The divisor minus a point]
  \label{pa-def:devissageDivisor}
  \lean{AlgebraicGeometry.devissageDivisor}
  \leanok
  \dcref{custom}

  \uses{pa-def:curveDivisor}
  For a closed point \(x\), the divisor \(D - x\) is obtained by decreasing the coefficient of \(D\)
  at \(x\) by one and leaving every other coefficient unchanged. Thus \((D-x)_x = D_x - 1\) and
  \(D - x \le D\).
\end{definition}

\begin{definition}[The skyscraper projection \(\pi\)]
  \label{pa-def:devissagePi}
  \lean{AlgebraicGeometry.devissageπ}
  \leanok
  \dcref{custom}

  \uses{pa-def:jumpProj, pa-def:skyModule, pa-def:divisorSheaf}
  The morphism of sheaves \(\pi : \mathcal O(D) \to \mathrm{sky}_x(J)\) sending a section \(g\) over
  an open \(U \ni x\) to its jump class \([g] \in J\), and every section over an open avoiding \(x\)
  to \(0\). Naturality holds because the jump projection depends only on the underlying rational
  function, and the restriction of the skyscraper between two opens containing \(x\) is the identity
  of \(J\).
\end{definition}

\begin{lemma}[The dévissage complex vanishes]
  \label{pa-lem:devissage_comp_eq_zero}
  \lean{AlgebraicGeometry.devissage_comp_eq_zero}
  \leanok
  \dcref{custom}

  \uses{pa-def:devissagePi, pa-def:devissageDivisor, pa-thm:divisorSheafLE, pa-def:jumpProj}
  The composite \(\mathcal O(D-x) \hookrightarrow \mathcal O(D) \xrightarrow{\ \pi\ }
  \mathrm{sky}_x(J)\) is zero.
\end{lemma}
\begin{proof}
  \leanok
  Over an open \(U \ni x\), a section of \(\mathcal O(D-x)\) has pole order \(\le (D-x)_x = D_x - 1\)
  at \(x\), so it lies in \(L_{D_x - 1}\) and its jump class in \(J = L_{D_x}/L_{D_x - 1}\) vanishes.
  Over an open avoiding \(x\) the target \(\mathrm{sky}_x(J)\) has zero sections. Hence the composite
  is zero on every open.
\end{proof}

\begin{definition}[The dévissage short complex]
  \label{pa-def:devissageSES}
  \lean{AlgebraicGeometry.devissageSES, AlgebraicGeometry.devissageSES_mono_f}
  \leanok
  \dcref{custom}

  \uses{pa-lem:devissage_comp_eq_zero, pa-thm:divisorSheafLE, pa-def:devissagePi}
  The \emph{dévissage short complex}
  \[
    0 \longrightarrow \mathcal O(D-x) \longrightarrow \mathcal O(D)
      \xrightarrow{\ \pi\ } \mathrm{sky}_x(J) \longrightarrow 0
  \]
  of sheaves of \(K\)-modules: the left map is the inclusion \ref{pa-thm:divisorSheafLE} for
  \(D - x \le D\) and the right map is \(\pi\); their composite vanishes
  (\ref{pa-lem:devissage_comp_eq_zero}). Its left leg is a monomorphism.
\end{definition}
\begin{proof}
  \leanok
  The left leg is the inclusion of \ref{pa-thm:divisorSheafLE}, which is a monomorphism of sheaves.
\end{proof}

\section{Short exactness of the dévissage sequence}

The dévissage short complex is short exact. The monomorphism is the inclusion of
\ref{pa-def:devissageSES}; the epimorphism is local surjectivity of \(\pi\), which at the support point
excises a finite bad locus to lift a jump class; middle exactness is checked on sections and
transported to the sheaf level by reflection of limits.

\begin{lemma}[Local lifting of a jump class]
  \label{pa-lem:devissage_local_lift}
  \lean{AlgebraicGeometry.devissage_local_lift}
  \leanok
  \dcref{custom}

  \uses{pa-def:jumpModule, pa-thm:ordZ_support_finite, pa-def:divOf, pa-thm:closedpoint_isClosed,
    pa-def:pointLattice, pa-def:divisorSections, pa-def:jumpProj}
  Let \(x \in U\) and let \(t\) be a section of \(\mathrm{sky}_x(J)\) over \(U\), i.e.\ a class in
  \(J\). There is an open \(W\) with \(x \in W \subseteq U\) and a section of \(\mathcal O(D)\) over
  \(W\) whose jump class is \(t|_W\).
\end{lemma}
\begin{proof}
  \leanok
  Choose a representative \(g \in L_{D_x}\) of the class \(t\) (the quotient
  \(L_{D_x} \twoheadrightarrow J\) is surjective). If \(t = 0\), the zero section over \(U\) itself
  has jump class \(0\). Otherwise \(g \neq 0\), so \(g \in K(X)^\times\). The closed points \(z\)
  where \(g\) violates the \(D\)-bound (pole order at \(z\) exceeding \(D_z\)) form a finite set:
  outside the support of \(\mathrm{div}(g)\) and the support of \(D\), the function \(g\) has neither
  zero nor pole and \(D_z = 0\), so the bound holds; both supports are finite
  (\ref{pa-thm:ordZ_support_finite} and finiteness of a divisor's support). This bad set is a finite
  union of closed points (\ref{pa-thm:closedpoint_isClosed}), hence closed, and does not contain \(x\),
  where \(g \in L_{D_x}\) already respects the bound. Let \(W = U \setminus (\text{bad set})\); then
  \(x \in W \subseteq U\) and \(g\) satisfies the \(D\)-bound at every closed point of \(W\), so
  \(g \in \mathcal O(D)(W)\), and by construction its jump class is \(t|_W\).
\end{proof}

\begin{theorem}[The projection is an epimorphism]
  \label{pa-thm:devissageSES_epi_g}
  \lean{AlgebraicGeometry.devissageSES_epi_g}
  \leanok
  \dcref{custom}

  \uses{pa-lem:devissage_local_lift, pa-def:devissagePi, pa-def:skyModule, pa-def:devissageSES}
  The dévissage projection \(\pi : \mathcal O(D) \to \mathrm{sky}_x(J)\) is an epimorphism of
  sheaves.
\end{theorem}
\begin{proof}
  \leanok
  It suffices that \(\pi\) is locally surjective. Over an open avoiding \(x\) the target skyscraper
  sections are zero, so the zero section lifts any section there; over a neighbourhood of \(x\)
  every jump section lifts by \ref{pa-lem:devissage_local_lift}. A locally surjective morphism of
  sheaves of \(K\)-modules is an epimorphism.
\end{proof}

\begin{lemma}[Sectionwise middle exactness]
  \label{pa-lem:devissageSES_map_exact}
  \lean{AlgebraicGeometry.devissageSES_map_exact}
  \leanok
  \dcref{custom}

  \uses{pa-def:jumpProj, pa-def:pointLattice, pa-def:devissageDivisor, pa-thm:divisorSheafLE, pa-def:devissagePi,
    pa-def:divisorSections}
  Over every open \(W\), the sequence of sections
  \(\mathcal O(D-x)(W) \to \mathcal O(D)(W) \xrightarrow{\pi_W} \mathrm{sky}_x(J)(W)\) is exact in
  \(K\)-modules.
\end{lemma}
\begin{proof}
  \leanok
  Let \(s\) be a section of \(\mathcal O(D)\) over \(W\) with \(\pi_W(s) = 0\). If \(x \in W\): the
  vanishing of \(\pi_W(s)\) means the jump class of \(s\) is zero, i.e.\ \(s\) has pole order
  \(\le D_x - 1\) at \(x\); together with the \(D\)-bound at the other points --- where \(D\) and
  \(D-x\) have equal coefficients --- this says \(s \in \mathcal O(D-x)(W)\). If \(x \notin W\): the
  divisors \(D\) and \(D-x\) agree on \(W\), so \(\mathcal O(D)(W) = \mathcal O(D-x)(W)\) and \(s\)
  already lies in the subspace. In either case \(s\) is the image of a \((D-x)\)-section under the
  inclusion, which is the kernel of \(\pi_W\).
\end{proof}

\begin{theorem}[Middle exactness of the dévissage complex]
  \label{pa-thm:devissageSES_exact}
  \lean{AlgebraicGeometry.devissageSES_exact}
  \leanok
  \dcref{custom}

  \uses{pa-lem:devissageSES_map_exact, pa-def:devissageSES, pa-def:devissagePi}
  The dévissage short complex is exact at \(\mathcal O(D)\).
\end{theorem}
\begin{proof}
  \leanok
  Over each open \(W\) the sections complex is exact (\ref{pa-lem:devissageSES_map_exact}) and its left
  leg is injective (\ref{pa-def:devissageSES}), so the left map realises \(\mathcal O(D-x)(W)\) as the
  kernel of \(\pi_W\) in \(K\)-modules. The forgetful functor from sheaves to presheaves creates
  limits, and evaluation at the opens jointly reflects them; hence these section-wise kernel forks
  assemble into a kernel fork of sheaves, exhibiting the left map as the kernel of \(\pi\). This is
  exactness in the middle.
\end{proof}

\begin{theorem}[The dévissage sequence is short exact]
  \label{pa-thm:devissageSES_shortExact}
  \lean{AlgebraicGeometry.devissageSES_shortExact}
  \leanok
  \dcref{custom}

  \uses{pa-def:devissageSES, pa-thm:devissageSES_epi_g, pa-thm:devissageSES_exact}
  The dévissage sequence
  \[
    0 \longrightarrow \mathcal O(D-x) \longrightarrow \mathcal O(D)
      \xrightarrow{\ \pi\ } \mathrm{sky}_x(J) \longrightarrow 0
  \]
  is a short exact sequence of sheaves of \(K\)-modules.
\end{theorem}
\begin{proof}
  \leanok
  The left leg is a monomorphism (\ref{pa-def:devissageSES}), the right leg \(\pi\) is an epimorphism
  (\ref{pa-thm:devissageSES_epi_g}), and the complex is exact in the middle
  (\ref{pa-thm:devissageSES_exact}); these are exactly the three certificates of a short exact
  sequence.
\end{proof}

\section{The jump module and the residue field}

The jump module at a closed point is the residue field, viewed as a \(K\)-space; its dimension is
therefore the residue degree. This is the local input identifying the degree-zero cohomology of the
skyscraper term with the residue degree of the point.

\begin{lemma}[A uniformizer at a closed point]
  \label{pa-lem:uniformizer}
  \lean{AlgebraicGeometry.uniformizer, AlgebraicGeometry.ord_uniformizer}
  \leanok
  \dcref{custom}

  \uses{pa-thm:closedpoint_dvr_stalk, pa-def:scheme_ord}
  There is a rational function \(t \in K(X)\) with \(\mathrm{ord}_x t = 1\): a uniformizer of the
  discrete valuation ring \(\struct{X,x}\) (\ref{pa-thm:closedpoint_dvr_stalk}), a function with a
  simple zero at \(x\).
\end{lemma}
\begin{proof}
  \leanok
  The stalk \(\struct{X,x}\) is a discrete valuation ring, so its maximal ideal has a generator
  \(t\), a uniformizer; the adic valuation \(\mathrm{ord}_x\) (\ref{pa-def:scheme_ord}) sends the
  generator to \(1\), so \(\mathrm{ord}_x t = 1\).
\end{proof}

\begin{theorem}[The shift isomorphism]
  \label{pa-thm:shiftMap}
  \lean{AlgebraicGeometry.shiftMap, AlgebraicGeometry.shiftMap_surjective}
  \leanok
  \dcref{custom}

  \uses{pa-lem:uniformizer, pa-def:pointLattice}
  Let \(t\) be a uniformizer at \(x\). Multiplication by \(t^a\) is a \(K\)-linear isomorphism
  \(L_a \xrightarrow{\ \sim\ } L_0\), carrying \(L_{a-1}\) onto \(L_{-1}\).
\end{theorem}
\begin{proof}
  \leanok
  Since \(\mathrm{ord}_x(t^a g) = a + \mathrm{ord}_x g\), a function \(g\) has pole order \(\le a\)
  at \(x\) iff \(t^a g\) has pole order \(\le 0\); so multiplication by \(t^a\) maps \(L_a\) into
  \(L_0\) (and \(L_{a-1}\) into \(L_{-1}\)), with two-sided inverse multiplication by \(t^{-a}\),
  hence is a bijection. It is \(K\)-linear because the \(K\)-action on \(K(X)\) is multiplication in
  the field, which commutes with multiplication by \(t^a\).
\end{proof}

\begin{theorem}[The base map to the residue field]
  \label{pa-thm:baseHom}
  \lean{AlgebraicGeometry.baseHom, AlgebraicGeometry.baseHom_surjective,
    AlgebraicGeometry.baseHom_ker}
  \leanok
  \dcref{custom}

  \uses{pa-lem:exists_stalk_of_ord_le_one, pa-def:pointLattice, pa-def:residueOverAlgebraMap,
    pa-thm:closedpoint_dvr_stalk}
  The residue map induces a \(K\)-linear surjection \(L_0 \twoheadrightarrow \kappa(x)\) with kernel
  \(L_{-1}\).
\end{theorem}
\begin{proof}
  \leanok
  A function \(g \in L_0\) has no pole at \(x\), hence is the image of a unique element of the stalk
  \(\struct{X,x}\) (\ref{pa-lem:exists_stalk_of_ord_le_one}); send \(g\) to the residue of that element
  in \(\kappa(x)\). This is \(K\)-linear for the residue \(K\)-structure of
  \ref{pa-def:residueOverAlgebraMap} (using injectivity of \(\struct{X,x} \to K(X)\)) and surjective,
  because the residue map \(\struct{X,x} \to \kappa(x)\) is (\ref{pa-thm:closedpoint_dvr_stalk}). Its
  kernel is exactly \(L_{-1}\): the residue vanishes iff the stalk preimage lies in the maximal
  ideal, i.e.\ has strictly positive order of vanishing, i.e.\ \(g\) has pole order \(\le -1\) --- a
  zero --- at \(x\).
\end{proof}

\begin{theorem}[The jump module is the residue field]
  \label{pa-thm:jumpEquivResidueField}
  \lean{AlgebraicGeometry.jumpEquivResidueField}
  \leanok
  \dcref{custom}

  \uses{pa-thm:shiftMap, pa-thm:baseHom, pa-def:jumpModule}
  There is a \(K\)-linear isomorphism
  \[
    J \;=\; L_{D_x}\big/L_{D_x - 1} \;\cong\; \kappa(x).
  \]
\end{theorem}
\begin{proof}
  \leanok
  Compose the shift isomorphism \(L_{D_x} \xrightarrow{\sim} L_0\) (\ref{pa-thm:shiftMap},
  multiplication by \(t^{D_x}\)) with the residue surjection \(L_0 \twoheadrightarrow \kappa(x)\)
  (\ref{pa-thm:baseHom}). The composite \(L_{D_x} \to \kappa(x)\) is surjective, and its kernel is the
  preimage of \(L_{-1}\) under the shift, namely \(L_{D_x - 1}\) (the shift carries \(L_{D_x-1}\)
  onto \(L_{-1}\)). By the first isomorphism theorem, \(J = L_{D_x}/L_{D_x-1} \cong \kappa(x)\).
\end{proof}

\begin{theorem}[Dimension of the jump module]
  \label{pa-thm:finrank_jumpModule}
  \lean{AlgebraicGeometry.finrank_jumpModule}
  \leanok
  \dcref{custom}

  \uses{pa-thm:jumpEquivResidueField, pa-def:residueDeg}
  The \(K\)-dimension of the jump module is the residue degree:
  \[
    \dim_K J \;=\; [\kappa(x):K] \;=\; \deg_K(x).
  \]
\end{theorem}
\begin{proof}
  \leanok
  Immediate from the isomorphism \(J \cong \kappa(x)\) (\ref{pa-thm:jumpEquivResidueField}) and the
  definition of the residue degree as \(\dim_K \kappa(x)\) (\ref{pa-def:residueDeg}).
\end{proof}

\begin{corollary}[The jump module is finite]
  \label{pa-cor:moduleFinite_jumpModule}
  \lean{AlgebraicGeometry.moduleFinite_jumpModule}
  \leanok
  \dcref{custom}

  \uses{pa-thm:jumpEquivResidueField, pa-thm:residueDeg_finite}
  The jump module \(J\) is a finite \(K\)-module.
\end{corollary}
\begin{proof}
  \leanok
  By \ref{pa-thm:jumpEquivResidueField}, \(J \cong \kappa(x)\), which is finite over \(K\)
  (\ref{pa-thm:residueDeg_finite}).
\end{proof}

\section{The multiplication isomorphism}

Multiplication by a nonzero rational function \(g\) identifies \(\mathcal O(D)\) with
\(\mathcal O(D - \mathrm{div}\, g)\). This is the engine forcing the degree of a principal divisor
to vanish: the two sheaves being isomorphic, their Euler characteristics agree.

\begin{theorem}[The bound-shift]
  \label{pa-thm:mem_boundedSections_mul_iff}
  \lean{AlgebraicGeometry.Scheme.mem_boundedSections_mul_iff}
  \leanok
  \dcref{custom}

  \uses{pa-def:divOf, pa-def:divisorSections, pa-def:scheme_ord}
  Let \(g \in K(X)^\times\) and \(h \in K(X)\). For any open \(U\), the product \(g h\) has poles
  bounded by \(D - \mathrm{div}(g)\) on \(U\) if and only if \(h\) has poles bounded by \(D\) on
  \(U\).
\end{theorem}
\begin{proof}
  \leanok
  At each closed point \(z \in U\), the valuation is multiplicative, \(\mathrm{ord}_z(gh) =
  \mathrm{ord}_z g + \mathrm{ord}_z h\), and \((D - \mathrm{div}\, g)_z = D_z - \mathrm{ord}_z g\)
  since \(\mathrm{div}(g)_z = \mathrm{ord}_z g\) (\ref{pa-def:divOf}). Hence the bound
  \(\mathrm{ord}_z(gh) \ge -(D - \mathrm{div}\, g)_z\) reads
  \(\mathrm{ord}_z g + \mathrm{ord}_z h \ge -D_z + \mathrm{ord}_z g\), i.e.\ \(\mathrm{ord}_z h \ge
  -D_z\): the finite value \(\mathrm{ord}_z g\) cancels, leaving exactly the \(D\)-bound for \(h\).
  As this equivalence holds at every closed point of \(U\), it holds for the whole section
  (\ref{pa-def:divisorSections}).
\end{proof}

\begin{theorem}[Multiplication isomorphism of divisor sheaves]
  \label{pa-thm:mulEquivDivisorSheaf}
  \lean{AlgebraicGeometry.Scheme.mulEquivDivisorSheaf}
  \leanok
  \source{papaioannou-algebraic-rr:page-0007}

  \uses{pa-thm:mem_boundedSections_mul_iff, pa-def:divisorSheaf, pa-def:divOf}
  \dcref{custom}
  For \(g \in K(X)^\times\) and any Weil divisor \(D\), multiplication by \(g\) is an isomorphism of
  sheaves of \(K\)-modules
  \[
    \mathcal O(D) \;\cong\; \mathcal O\bigl(D - \mathrm{div}\, g\bigr).
  \]
\end{theorem}
\begin{proof}
  \leanok
  Multiplication by the unit \(g\) is a \(K\)-linear automorphism of \(K(X)\), with inverse
  multiplication by \(g^{-1}\). By the bound-shift \ref{pa-thm:mem_boundedSections_mul_iff} it
  restricts, on every nonempty open \(U\), to a bijection \(\mathcal O(D)(U) \to
  \mathcal O(D - \mathrm{div}\, g)(U)\) (the inverse being restriction of multiplication by
  \(g^{-1}\)), and both sides are the zero module over the empty open. Being bijective on every
  open, these maps assemble into an isomorphism of sheaves. Together with the additivity of
  \(\mathrm{div}\) this is the engine of \(\deg(\mathrm{div}\, g) = 0\): it identifies
  \(\mathcal O(D)\) with \(\mathcal O(D - \mathrm{div}\, g)\), forcing their Euler characteristics to
  agree.
\end{proof}

\section{The cohomological ledger:
  \texorpdfstring{\(h^0\), \(h^1\) and \(\chi\)}{h0, h1 and chi}}

The Euler characteristic of the Riemann--Roch ledger is a numerical invariant of a sheaf of
modules. On an arbitrary small site \((\mathcal C, J)\) with coefficients in a commutative ring
\(R\), any sheaf \(F\) of \(R\)-modules has cohomology \(H^n(F) = \Ext^n(R_J, F)\)
(Definition~\ref{pa-def:HModule}); this section attaches to \(F\) its dimension counts in degrees
\(0\) and \(1\) and their alternating combination, and records their invariance under sheaf
isomorphism. Nothing here is scheme-specific --- the curve enters only later, where all higher
cohomology vanishes and the truncation below is the honest Euler characteristic.

\begin{definition}[Transport of cohomology along a sheaf isomorphism]
  \label{pa-def:HModule_mapEquiv}
  \lean{CategoryTheory.Sheaf.HModule.mapEquiv}
  \leanok
  \dcref{custom}

  \uses{pa-def:HModule_map}
  Let \(F, G\) be sheaves of \(R\)-modules on \((\mathcal C, J)\) and \(e : F \cong G\) an
  isomorphism of sheaves. For every \(n\), \(e\) induces an \(R\)-linear equivalence
  \[
    H^n(F) \;\cong\; H^n(G),
  \]
  with forward map the functorial \(H^n(e)\) of the isomorphism and inverse the functorial
  \(H^n(e^{-1})\) of its inverse.
\end{definition}
\begin{proof}
  \leanok
  Take the forward map to be \(H^n(e)\) and the backward map \(H^n(e^{-1})\)
  (Definition~\ref{pa-def:HModule_map}). By functoriality of \(H^n\), their composites are
  \(H^n(e^{-1}) \circ H^n(e) = H^n(e^{-1} \circ e) = H^n(\id_F) = \id\) and symmetrically
  \(H^n(e) \circ H^n(e^{-1}) = \id\); so the two maps are mutually inverse \(R\)-linear
  bijections, an \(R\)-linear equivalence.
\end{proof}

\begin{definition}[The ledger invariants \(h^0\), \(h^1\) and \(\chi\)]
  \label{pa-def:h0_h1_chi}
  \lean{CategoryTheory.Sheaf.h0, CategoryTheory.Sheaf.h1, CategoryTheory.Sheaf.chi}
  \leanok
  \dcref{custom}

  \uses{pa-def:HModule}
  For a sheaf \(F\) of \(R\)-modules on \((\mathcal C, J)\), set
  \[
    h^0(F) = \dim_R H^0(F), \qquad h^1(F) = \dim_R H^1(F) \quad\in \mathbb N,
  \]
  the \(R\)-dimensions of the degree-zero and degree-one cohomology, and
  \[
    \chi(F) \;=\; h^0(F) - h^1(F) \quad\in \Z,
  \]
  the \emph{(truncated) Euler characteristic}. The dimensions are read by \(\finrank\), which
  returns \(0\) on a module that is not finite-dimensional; no finiteness hypothesis is imposed
  here, the finiteness bookkeeping being discharged downstream. On a curve, where all higher
  cohomology vanishes, \(\chi\) is the honest Euler characteristic.
\end{definition}

\begin{theorem}[Isomorphism-invariance of the ledger]
  \label{pa-thm:chi_congr}
  \lean{CategoryTheory.Sheaf.h0_congr, CategoryTheory.Sheaf.h1_congr,
    CategoryTheory.Sheaf.chi_congr}
  \leanok
  \dcref{custom}

  \uses{pa-def:h0_h1_chi, pa-def:HModule_mapEquiv}
  Isomorphic sheaves have equal ledger invariants: if \(F \cong G\) then
  \(h^0(F) = h^0(G)\), \(h^1(F) = h^1(G)\) and \(\chi(F) = \chi(G)\).
\end{theorem}
\begin{proof}
  \leanok
  An isomorphism \(F \cong G\) induces an \(R\)-linear equivalence
  \(H^n(F) \cong H^n(G)\) in each degree (Definition~\ref{pa-def:HModule_mapEquiv}); linearly
  equivalent modules have equal dimension, so \(h^0(F) = h^0(G)\) and \(h^1(F) = h^1(G)\).
  Subtracting gives \(\chi(F) = \chi(G)\).
\end{proof}

\begin{theorem}[Degenerate ledger when \(H^1\) vanishes]
  \label{pa-thm:chi_eq_h0}
  \lean{CategoryTheory.Sheaf.h1_eq_zero, CategoryTheory.Sheaf.chi_eq_h0}
  \leanok
  \dcref{custom}

  \uses{pa-def:h0_h1_chi}
  If \(R\) is nontrivial and \(H^1(F) = 0\), then \(h^1(F) = 0\) and
  \(\chi(F) = h^0(F)\).
\end{theorem}
\begin{proof}
  \leanok
  A zero module has dimension \(0\) over a nontrivial ring, so \(h^1(F) = \dim_R H^1(F) = 0\);
  substituting into \(\chi(F) = h^0(F) - h^1(F)\) gives \(\chi(F) = h^0(F)\).
\end{proof}

\section{The six-term cohomology slice of a short exact sequence}

A short exact sequence \(0 \to X_1 \to X_2 \to X_3 \to 0\) of sheaves of \(R\)-modules gives a
long exact cohomology sequence; its degree-\(\le 1\) slice
\[
  0 \to H^0 X_1 \to H^0 X_2 \to H^0 X_3 \xrightarrow{\ \delta\ }
    H^1 X_1 \to H^1 X_2 \to H^1 X_3
\]
is the exactness input to both the finiteness bookkeeping and the additivity of \(\chi\). This
section packages the slice --- its connecting map, its exactness, and its two ends --- and, over a
field, the finiteness transfer and the additivity of \(\chi\) that ride on it. The additivity rests
on the vanishing of the alternating sum of dimensions along the five-term exact sequence
(Theorem~\ref{pa-thm:finrank_alt_sum_exact5}).

\begin{theorem}[Extension of finite dimensionality along an exact pair]
  \label{pa-thm:finiteDimensional_of_exact}
  \lean{FiniteDimensional.of_exact}
  \leanok
  \dcref{custom}

  Let \(K\) be a division ring and let \(V_1 \xrightarrow{f} V_2 \xrightarrow{g} V_3\) be a pair
  of \(K\)-linear maps, exact at \(V_2\) (\(\ker g = \operatorname{im} f\)), with \(V_1\) and
  \(V_3\) finite-dimensional. Then \(V_2\) is finite-dimensional.
\end{theorem}
\begin{proof}
  \leanok
  Corestrict \(g\) to its range, \(g' : V_2 \to \operatorname{im} g\); this is surjective, and the
  pair \((f, g')\) has the same kernel of second map as \((f, g)\), so it is still exact at
  \(V_2\). The range \(\operatorname{im} g\) is a subspace of the finite-dimensional \(V_3\), hence
  finite-dimensional. Thus \(V_2\) sits in an exact sequence
  \(V_1 \to V_2 \to \operatorname{im} g \to 0\) with surjective right map and finite-dimensional
  ends: its kernel \(\operatorname{im} f\) is a quotient of the finite \(V_1\) and its coimage is
  the finite \(\operatorname{im} g\), so \(V_2\) is finite-dimensional.
\end{proof}

\begin{definition}[The connecting homomorphism]
  \label{pa-def:HModule_delta}
  \lean{CategoryTheory.Sheaf.HModule.delta}
  \leanok
  \dcref{custom}

  \uses{pa-def:HModule, pa-def:constModuleSheaf}
  For a short exact sequence \(0 \to X_1 \to X_2 \to X_3 \to 0\) of sheaves of \(R\)-modules and
  \(n_0 + 1 = n_1\), the \emph{connecting homomorphism}
  \[
    \delta : H^{n_0}(X_3) \longrightarrow H^{n_1}(X_1)
  \]
  is postcomposition of a class in \(\Ext^{n_0}(R_J, X_3)\) with the extension class
  \(\Ext^1(X_3, X_1)\) of the sequence.
\end{definition}

\begin{theorem}[Exactness of the cohomology slice]
  \label{pa-thm:HModule_slice_exact}
  \lean{CategoryTheory.Sheaf.HModule.exact_map_f_map_g,
    CategoryTheory.Sheaf.HModule.exact_map_g_delta,
    CategoryTheory.Sheaf.HModule.exact_delta_map_f}
  \leanok
  \dcref{custom}

  \uses{pa-def:HModule_map, pa-def:HModule_delta}
  For a short exact sequence \(0 \to X_1 \to X_2 \to X_3 \to 0\) of sheaves of \(R\)-modules and
  \(n_0 + 1 = n_1\), the cohomology sequence is exact at each of its three interior joints:
  \[
    H^{n}(X_1) \to H^{n}(X_2) \to H^{n}(X_3), \qquad
    H^{n_0}(X_2) \to H^{n_0}(X_3) \xrightarrow{\delta} H^{n_1}(X_1), \qquad
    H^{n_0}(X_3) \xrightarrow{\delta} H^{n_1}(X_1) \to H^{n_1}(X_2).
  \]
\end{theorem}
\begin{proof}
  \leanok
  These are the three interior slots of the long exact sequence of the covariant functor
  \(\Ext^\bullet(R_J, -)\) applied to the short exact sequence, read elementwise. The maps
  \(H^n(X_i \to X_j)\) are the functorial maps of Definition~\ref{pa-def:HModule_map} and the
  connecting map is \(\delta\) of Definition~\ref{pa-def:HModule_delta}; the identification of the
  covariant \(\Ext\)-sequence's maps with these is by definition of postcomposition on both
  sides.
\end{proof}

\begin{theorem}[Left end of the slice: injectivity at \(H^0\)]
  \label{pa-thm:HModule_injective_map_f_zero}
  \lean{CategoryTheory.Sheaf.HModule.injective_map_f_zero}
  \leanok
  \dcref{custom}

  \uses{pa-def:HModule_map}
  For a short exact sequence \(0 \to X_1 \xrightarrow{f} X_2 \to X_3 \to 0\) of sheaves of
  \(R\)-modules, the map \(H^0(f) : H^0(X_1) \to H^0(X_2)\) is injective.
\end{theorem}
\begin{proof}
  \leanok
  The map \(f : X_1 \to X_2\) is a monomorphism, and postcomposition with a monomorphism is
  injective on \(\Ext^0(R_J, -) = \Hom(R_J, -)\); this is exactly injectivity of \(H^0(f)\).
\end{proof}

\begin{theorem}[Right end of the slice: surjectivity into a vanishing term]
  \label{pa-thm:HModule_surjective_map_f}
  \lean{CategoryTheory.Sheaf.HModule.surjective_map_f}
  \leanok
  \dcref{custom}

  \uses{pa-thm:HModule_slice_exact}
  For a short exact sequence \(0 \to X_1 \xrightarrow{f} X_2 \to X_3 \to 0\) of sheaves of
  \(R\)-modules and a degree \(n\) with \(H^n(X_3) = 0\), the map
  \(H^n(f) : H^n(X_1) \to H^n(X_2)\) is surjective.
\end{theorem}
\begin{proof}
  \leanok
  Exactness at \(H^n(X_2)\) (Theorem~\ref{pa-thm:HModule_slice_exact}) identifies the image of
  \(H^n(f)\) with the kernel of \(H^n(X_2) \to H^n(X_3)\). Since \(H^n(X_3) = 0\), that map is
  zero and its kernel is all of \(H^n(X_2)\); hence \(H^n(f)\) is surjective.
\end{proof}

\begin{theorem}[Finiteness transfer along the slice]
  \label{pa-thm:HModule_finiteness_transfer}
  \lean{CategoryTheory.Sheaf.HModule.moduleFinite_middle,
    CategoryTheory.Sheaf.HModule.moduleFinite_left_zero,
    CategoryTheory.Sheaf.HModule.moduleFinite_left_succ}
  \leanok
  \dcref{custom}

  \uses{pa-thm:finiteDimensional_of_exact, pa-thm:HModule_slice_exact,
    pa-thm:HModule_injective_map_f_zero, pa-def:HModule_delta}
  Let \(R\) be a field and \(0 \to X_1 \to X_2 \to X_3 \to 0\) a short exact sequence of sheaves
  of \(R\)-modules, with \(n_0 + 1 = n_1\). Then:
  \begin{enumerate}
    \item if \(H^n(X_1)\) and \(H^n(X_3)\) are finite-dimensional, so is \(H^n(X_2)\);
    \item if \(H^0(X_2)\) is finite-dimensional, so is \(H^0(X_1)\);
    \item if \(H^{n_0}(X_3)\) and \(H^{n_1}(X_2)\) are finite-dimensional, so is \(H^{n_1}(X_1)\).
  \end{enumerate}
\end{theorem}
\begin{proof}
  \leanok
  (1) The sequence \(H^n(X_1) \to H^n(X_2) \to H^n(X_3)\) is exact at \(H^n(X_2)\)
  (Theorem~\ref{pa-thm:HModule_slice_exact}) with finite-dimensional ends, so \(H^n(X_2)\) is
  finite-dimensional by the extension lemma \ref{pa-thm:finiteDimensional_of_exact}. (2) The map
  \(H^0(X_1) \to H^0(X_2)\) is injective (Theorem~\ref{pa-thm:HModule_injective_map_f_zero}), so
  \(H^0(X_1)\) embeds in the finite-dimensional \(H^0(X_2)\) and is finite-dimensional. (3) The
  sequence \(H^{n_0}(X_3) \xrightarrow{\delta} H^{n_1}(X_1) \to H^{n_1}(X_2)\) is exact at
  \(H^{n_1}(X_1)\) (Theorem~\ref{pa-thm:HModule_slice_exact}, with \(\delta\) of
  Definition~\ref{pa-def:HModule_delta}) and has finite-dimensional ends, so \(H^{n_1}(X_1)\) is
  finite-dimensional by \ref{pa-thm:finiteDimensional_of_exact}.
\end{proof}

\begin{theorem}[Additivity of \(\chi\) along a short exact sequence]
  \label{pa-thm:chi_eq_add_of_shortExact}
  \lean{CategoryTheory.Sheaf.chi_eq_add_of_shortExact}
  \leanok
  \dcref{custom}

  \uses{pa-thm:finrank_alt_sum_exact5, pa-thm:HModule_injective_map_f_zero, pa-thm:HModule_slice_exact,
    pa-thm:HModule_surjective_map_f, pa-def:HModule_delta, pa-thm:chi_eq_h0, pa-def:h0_h1_chi}
  Let \(R\) be a field and \(0 \to X_1 \to X_2 \to X_3 \to 0\) a short exact sequence of sheaves
  of \(R\)-modules with \(H^1(X_3) = 0\) and all of \(H^0(X_1), H^0(X_2), H^0(X_3), H^1(X_1),
  H^1(X_2)\) finite-dimensional. Then
  \[
    \chi(X_2) \;=\; \chi(X_1) + \chi(X_3).
  \]
\end{theorem}
\begin{proof}
  \leanok
  Since \(H^1(X_3) = 0\), the degree-\(\le 1\) slice truncates to the five-term sequence
  \[
    0 \to H^0 X_1 \to H^0 X_2 \to H^0 X_3 \xrightarrow{\delta}
      H^1 X_1 \to H^1 X_2 \to 0,
  \]
  exact throughout: the left map is injective (Theorem~\ref{pa-thm:HModule_injective_map_f_zero}),
  the three interior joints are exact (Theorem~\ref{pa-thm:HModule_slice_exact}), and the right map
  is surjective because \(H^1 X_3 = 0\) (Theorem~\ref{pa-thm:HModule_surjective_map_f}). All five
  terms being finite-dimensional, the alternating sum of their dimensions vanishes
  (Theorem~\ref{pa-thm:finrank_alt_sum_exact5}):
  \[
    h^0 X_1 - h^0 X_2 + h^0 X_3 - h^1 X_1 + h^1 X_2 \;=\; 0 .
  \]
  Moreover \(h^1(X_3) = 0\) (Theorem~\ref{pa-thm:chi_eq_h0}). Rearranging the vanishing alternating
  sum,
  \[
    (h^0 X_2 - h^1 X_2) = (h^0 X_1 - h^1 X_1) + (h^0 X_3 - h^1 X_3),
  \]
  which is \(\chi(X_2) = \chi(X_1) + \chi(X_3)\).
\end{proof}

\section{Finiteness of the cohomology of divisor sheaves}

We return to the curve bundle \(X\): integral, smooth of relative dimension one and quasi-compact
over the field \(K\). Assume the structure sheaf, as a sheaf of \(K\)-modules, has
finite-dimensional \(H^0\) and \(H^1\) --- the two properness inputs, taken here as standing
hypotheses. This section proves that then \emph{every} divisor sheaf \(\mathcal O(D)\) has
finite-dimensional \(H^0\) and \(H^1\). The engine is dévissage: a one-point induction that
transfers finiteness through the six-term slice of the dévissage short exact sequence
\(0 \to \mathcal O(D-x) \to \mathcal O(D) \to \mathrm{sky}_x(J) \to 0\), whose skyscraper term
has finite cohomology.

\begin{definition}[The one-point divisor]
  \label{pa-def:curveDivisor_single}
  \lean{AlgebraicGeometry.Scheme.CurveDivisor.single}
  \leanok
  \dcref{custom}

  \uses{pa-def:curveDivisor}
  For a closed point \(x\) (a point other than the generic point \(\eta\)) and \(n \in \Z\), the
  \emph{one-point divisor} \(n \cdot x\) is the Weil divisor supported at \(x\) with coefficient
  \(n\).
\end{definition}

\begin{lemma}[Arithmetic of the one-point divisor]
  \label{pa-lem:curveDivisor_single_arith}
  \lean{AlgebraicGeometry.Scheme.CurveDivisor.single_zero,
    AlgebraicGeometry.Scheme.CurveDivisor.single_add,
    AlgebraicGeometry.Scheme.CurveDivisor.nsmul_single_one,
    AlgebraicGeometry.Scheme.CurveDivisor.deg_single'}
  \leanok
  \dcref{custom}

  \uses{pa-def:curveDivisor_single, pa-lem:curveDivisor_deg_single, pa-def:residueDeg}
  The one-point divisor at a fixed closed point \(x\) is additive in its multiplicity and
  compatible with the group operations and the degree:
  \[
    0 \cdot x = 0, \qquad m \cdot x + n \cdot x = (m+n)\cdot x, \qquad
    n \cdot (1 \cdot x) = n \cdot x \ (n \in \mathbb N), \qquad
    \deg_K(n \cdot x) = n \cdot [\kappa(x):K].
  \]
\end{lemma}
\begin{proof}
  \leanok
  The one-point divisor \(n \cdot x\) is the finitely supported function sending \(x\) to \(n\)
  and every other point to \(0\); such functions are additive in the value and vanish at value
  \(0\), giving the first two identities, and the third follows by induction on \(n\) from the
  second. The degree identity is the single-point degree formula
  (Lemma~\ref{pa-lem:curveDivisor_deg_single}) with the residue degree
  \([\kappa(x):K]\) of Definition~\ref{pa-def:residueDeg}.
\end{proof}

\begin{lemma}[The dévissage twist is subtraction of a point]
  \label{pa-lem:devissageDivisor_sub}
  \lean{AlgebraicGeometry.devissageDivisor_eq_sub,
    AlgebraicGeometry.devissageDivisor_add_single}
  \leanok
  \dcref{custom}

  \uses{pa-def:devissageDivisor, pa-def:curveDivisor_single}
  The dévissage twist \(D - x\) (Definition~\ref{pa-def:devissageDivisor}) is the subtraction of the
  one-point divisor \(1 \cdot x\):
  \[
    D - x = D - 1 \cdot x, \qquad (D - x) + 1\cdot x = D .
  \]
\end{lemma}
\begin{proof}
  \leanok
  The twist decreases the coefficient of \(D\) at \(x\) by one and leaves the others unchanged,
  which is exactly subtracting the one-point divisor \(1 \cdot x\); adding \(1 \cdot x\) back
  restores the original coefficient, giving \(D\).
\end{proof}

\begin{lemma}[Degree drop along the twist]
  \label{pa-lem:deg_devissageDivisor}
  \lean{AlgebraicGeometry.deg_devissageDivisor}
  \leanok
  \dcref{custom}

  \uses{pa-lem:devissageDivisor_sub, pa-thm:curveDivisor_deg_hom, pa-lem:curveDivisor_single_arith}
  The degree drops by the residue degree along the dévissage twist:
  \[
    \deg_K(D - x) \;=\; \deg_K D - [\kappa(x):K].
  \]
\end{lemma}
\begin{proof}
  \leanok
  The degree is a homomorphism (Theorem~\ref{pa-thm:curveDivisor_deg_hom}), so
  \(\deg_K\bigl((D - x) + 1\cdot x\bigr) = \deg_K(D - x) + \deg_K(1 \cdot x)\); the left side is
  \(\deg_K D\) since \((D - x) + 1 \cdot x = D\) (Lemma~\ref{pa-lem:devissageDivisor_sub}), and
  \(\deg_K(1 \cdot x) = [\kappa(x):K]\) (Lemma~\ref{pa-lem:curveDivisor_single_arith}). Subtracting
  gives the claim.
\end{proof}

\begin{theorem}[Dévissage induction on Weil divisors]
  \label{pa-thm:induction_devissage}
  \lean{AlgebraicGeometry.Scheme.CurveDivisor.induction_devissage}
  \leanok
  \dcref{custom}

  \uses{pa-def:curveDivisor, pa-def:devissageDivisor, pa-lem:devissageDivisor_sub}
  Let \(P\) be a predicate on Weil divisors of \(X\) that holds for the zero divisor and, for
  every closed point \(x\), transfers in both directions along the one-point dévissage:
  \(P(D) \iff P(D - x)\) for all \(D\). Then \(P(D)\) holds for every \(D\).
\end{theorem}
\begin{proof}
  \leanok
  A Weil divisor is a finitely supported \(\Z\)-family on the closed points, so it suffices to
  prove the following single-bump induction: a predicate on such families that holds at \(0\) and
  satisfies \(P(F + 1 \cdot x) \iff P(F)\) for every point \(x\) and family \(F\) holds
  everywhere. First, from the \(\pm 1\)-bump equivalence one gets \(P(F + n \cdot x) \iff P(F)\)
  for every \(n \in \Z\), by induction on \(n\): the case \(n = 0\) is trivial, and the step
  from \(n\) to \(n \pm 1\) composes one more bump equivalence, using
  \(F + (n\pm 1)\cdot x = (F + n\cdot x) \pm 1\cdot x\). Now argue by the finitely-supported
  induction principle: \(P(0)\) holds by hypothesis, and adding a single generator \(b \cdot x\)
  to a family \(F\) preserves \(P\) by the equivalence just established. Hence \(P\) holds for
  every family, i.e.\ every Weil divisor. The hypothesis \(P(D) \iff P(D - x)\) is the bump
  equivalence in the form \(P\bigl((D - x) + 1\cdot x\bigr) \iff P(D - x)\), using
  \((D - x) + 1\cdot x = D\) (Lemma~\ref{pa-lem:devissageDivisor_sub}).
\end{proof}

\begin{theorem}[Cohomology of a skyscraper is finite]
  \label{pa-thm:moduleFinite_skyModule}
  \lean{AlgebraicGeometry.moduleFinite_hModule_skyModule_zero,
    AlgebraicGeometry.moduleFinite_hModule_skyModule_one}
  \leanok
  \dcref{custom}

  \uses{pa-def:skyModuleGammaEquiv, pa-thm:skyModule_h1_vanishing}
  For a point \(x \in X\) and a finite-dimensional \(K\)-module \(M\), the cohomology of the
  skyscraper \(\mathrm{sky}_x(M)\) is finite-dimensional in degrees \(0\) and \(1\):
  \(H^0\) is \(M\) and \(H^1\) vanishes.
\end{theorem}
\begin{proof}
  \leanok
  In degree \(0\), \(H^0(\mathrm{sky}_x(M)) \cong M\) (Definition~\ref{pa-def:skyModuleGammaEquiv}),
  which is finite-dimensional by hypothesis. In degree \(1\),
  \(H^1(\mathrm{sky}_x(M)) = 0\) (Theorem~\ref{pa-thm:skyModule_h1_vanishing}), and the zero module
  is finite-dimensional.
\end{proof}

\begin{theorem}[Finiteness of the cohomology of every divisor sheaf]
  \label{pa-thm:moduleFinite_divisorSheaf}
  \lean{AlgebraicGeometry.moduleFinite_hModule_divisorSheaf_zero,
    AlgebraicGeometry.moduleFinite_hModule_divisorSheaf_one}
  \leanok
  \dcref{custom}

  \uses{pa-thm:induction_devissage, pa-thm:divisorSheafZeroIso, pa-thm:devissageSES_shortExact,
    pa-thm:moduleFinite_skyModule, pa-cor:moduleFinite_jumpModule, pa-thm:HModule_finiteness_transfer,
    pa-def:jumpModule}
  On the curve bundle \(X\), quasi-compact over \(K\), assume \(H^0(X, \struct{X/K})\) and
  \(H^1(X, \struct{X/K})\) are finite-dimensional over \(K\). Then for every Weil divisor \(D\),
  both \(H^0(\mathcal O(D))\) and \(H^1(\mathcal O(D))\) are finite-dimensional.
\end{theorem}
\begin{proof}
  \leanok
  Apply the dévissage induction (Theorem~\ref{pa-thm:induction_devissage}) to the predicate
  ``\(H^0(\mathcal O(D))\) and \(H^1(\mathcal O(D))\) are both finite-dimensional''.

  \emph{Base case.} For \(D = 0\), the identification \(\mathcal O(0) \cong \struct{X/K}\)
  (Theorem~\ref{pa-thm:divisorSheafZeroIso}) transports finite-dimensionality of \(H^0\) and \(H^1\)
  from the structure sheaf, which holds by the standing hypothesis.

  \emph{Step.} Fix a closed point \(x\). The dévissage short exact sequence
  \(0 \to \mathcal O(D-x) \to \mathcal O(D) \to \mathrm{sky}_x(J) \to 0\)
  (Theorem~\ref{pa-thm:devissageSES_shortExact}) has third term the skyscraper on the jump module
  \(J\) (Definition~\ref{pa-def:jumpModule}), which is finite-dimensional
  (Corollary~\ref{pa-cor:moduleFinite_jumpModule}); hence \(H^0(\mathrm{sky}_x(J))\) is
  finite-dimensional and \(H^1(\mathrm{sky}_x(J)) = 0\)
  (Theorem~\ref{pa-thm:moduleFinite_skyModule}). We must show the finiteness of the cohomology of
  \(\mathcal O(D)\) is equivalent to that of \(\mathcal O(D-x)\); both directions use the
  finiteness transfer of Theorem~\ref{pa-thm:HModule_finiteness_transfer} for this sequence, whose
  terms \(X_1, X_2, X_3\) are \(\mathcal O(D-x), \mathcal O(D), \mathrm{sky}_x(J)\).
  \begin{itemize}
    \item \emph{From \(D-x\) to \(D\).} If \(H^0, H^1\) of \(\mathcal O(D-x)\) are
      finite-dimensional, then so are those of \(\mathcal O(D)\), by the middle-term transfer
      (1) applied in degrees \(0\) and \(1\): the two outer terms \(\mathcal O(D-x)\) and
      \(\mathrm{sky}_x(J)\) have finite cohomology.
    \item \emph{From \(D\) to \(D-x\).} If \(H^0, H^1\) of \(\mathcal O(D)\) are
      finite-dimensional, then \(H^0(\mathcal O(D-x))\) is finite-dimensional as it embeds in
      \(H^0(\mathcal O(D))\) (left-term transfer (2)), and \(H^1(\mathcal O(D-x))\) is
      finite-dimensional as an extension controlled by the finite \(H^0(\mathrm{sky}_x(J))\) and
      \(H^1(\mathcal O(D))\) (left-term transfer (3)).
  \end{itemize}
  Both directions established, the induction concludes.
\end{proof}

\section{The \texorpdfstring{\(\chi\)}{chi}-ledger and the Riemann inequality}

The Euler characteristic of a divisor sheaf is now computed. Under the standing finiteness
hypotheses of the previous section, subtracting a closed point drops \(\chi\) by the residue
degree; iterating from \(\mathcal O(0) \cong \struct{X/K}\) closes the ledger into
\(\chi(\mathcal O(D)) = \chi(\struct{X/K}) + \deg_K D\). Two consequences follow: principal
divisors have degree zero, and \(h^0\) grows without bound along multiples of any closed point.

\begin{theorem}[The Euler characteristic of a skyscraper]
  \label{pa-thm:chi_skyModule}
  \lean{AlgebraicGeometry.h0_skyModule, AlgebraicGeometry.chi_skyModule}
  \leanok
  \dcref{custom}

  \uses{pa-def:skyModuleGammaEquiv, pa-thm:skyModule_h1_vanishing, pa-def:h0_h1_chi, pa-thm:chi_eq_h0}
  For a point \(x \in X\) and a \(K\)-module \(M\) (with \(K\) a field), the degree-zero
  dimension and the Euler characteristic of the skyscraper are both the dimension of the value
  module:
  \[
    h^0(\mathrm{sky}_x(M)) = \dim_K M, \qquad \chi(\mathrm{sky}_x(M)) = \dim_K M .
  \]
\end{theorem}
\begin{proof}
  \leanok
  Since \(H^0(\mathrm{sky}_x(M)) \cong M\) (Definition~\ref{pa-def:skyModuleGammaEquiv}),
  \(h^0(\mathrm{sky}_x(M)) = \dim_K M\). The degree-one cohomology vanishes
  (Theorem~\ref{pa-thm:skyModule_h1_vanishing}), so
  \(\chi(\mathrm{sky}_x(M)) = h^0(\mathrm{sky}_x(M))\) (Theorem~\ref{pa-thm:chi_eq_h0}), which is
  \(\dim_K M\).
\end{proof}

\begin{theorem}[The dévissage step of the \(\chi\)-ledger]
  \label{pa-thm:chi_step}
  \lean{AlgebraicGeometry.chi_step}
  \leanok
  \dcref{custom}

  \uses{pa-thm:chi_eq_add_of_shortExact, pa-thm:devissageSES_shortExact, pa-thm:chi_skyModule,
    pa-thm:finrank_jumpModule, pa-thm:moduleFinite_divisorSheaf, pa-thm:moduleFinite_skyModule,
    pa-cor:moduleFinite_jumpModule, pa-def:devissageDivisor}
  Under the standing finiteness hypotheses, subtracting a closed point \(x\) from a Weil divisor
  \(D\) drops the Euler characteristic by the residue degree:
  \[
    \chi(\mathcal O(D)) \;=\; \chi(\mathcal O(D - x)) + [\kappa(x):K].
  \]
\end{theorem}
\begin{proof}
  \leanok
  Apply the additivity of \(\chi\) (Theorem~\ref{pa-thm:chi_eq_add_of_shortExact}) to the dévissage
  short exact sequence \(0 \to \mathcal O(D-x) \to \mathcal O(D) \to \mathrm{sky}_x(J) \to 0\)
  (Theorem~\ref{pa-thm:devissageSES_shortExact}). Its hypotheses hold: the skyscraper term has
  \(H^1 = 0\) (Theorem~\ref{pa-thm:moduleFinite_skyModule}), and all five relevant cohomology
  modules are finite-dimensional --- those of \(\mathcal O(D-x)\) and \(\mathcal O(D)\) by
  Theorem~\ref{pa-thm:moduleFinite_divisorSheaf}, and \(H^0(\mathrm{sky}_x(J))\) because \(J\) is
  finite (Corollary~\ref{pa-cor:moduleFinite_jumpModule}). Hence
  \(\chi(\mathcal O(D)) = \chi(\mathcal O(D-x)) + \chi(\mathrm{sky}_x(J))\). Finally
  \(\chi(\mathrm{sky}_x(J)) = \dim_K J\) (Theorem~\ref{pa-thm:chi_skyModule}) and
  \(\dim_K J = [\kappa(x):K]\) (Theorem~\ref{pa-thm:finrank_jumpModule}), giving the stated identity.
\end{proof}

\begin{theorem}[The \(\chi\)-ledger, closed]
  \label{pa-thm:chi_divisorSheaf}
  \lean{AlgebraicGeometry.chi_divisorSheaf}
  \leanok
  \dcref{custom}

  \uses{pa-thm:induction_devissage, pa-thm:divisorSheafZeroIso, pa-thm:chi_congr, pa-thm:chi_step,
    pa-lem:deg_devissageDivisor, pa-thm:curveDivisor_deg_hom, pa-lem:devissageDivisor_sub}
  Under the standing finiteness hypotheses, for every Weil divisor \(D\),
  \[
    \chi(\mathcal O(D)) \;=\; \chi(\struct{X/K}) + \deg_K D .
  \]
\end{theorem}
\begin{proof}
  \leanok
  Dévissage induction (Theorem~\ref{pa-thm:induction_devissage}) on the predicate
  ``\(\chi(\mathcal O(D)) = \chi(\struct{X/K}) + \deg_K D\)''. For \(D = 0\):
  \(\chi(\mathcal O(0)) = \chi(\struct{X/K})\) by isomorphism-invariance
  (Theorem~\ref{pa-thm:chi_congr}) along \(\mathcal O(0) \cong \struct{X/K}\)
  (Theorem~\ref{pa-thm:divisorSheafZeroIso}), and \(\deg_K 0 = 0\)
  (Theorem~\ref{pa-thm:curveDivisor_deg_hom}). For the step at a closed point \(x\), combine the
  ledger step \(\chi(\mathcal O(D)) = \chi(\mathcal O(D-x)) + [\kappa(x):K]\)
  (Theorem~\ref{pa-thm:chi_step}) with the degree drop
  \(\deg_K(D-x) = \deg_K D - [\kappa(x):K]\) (Lemma~\ref{pa-lem:deg_devissageDivisor}): the
  statement for \(D-x\) and the statement for \(D\) differ on both sides by \([\kappa(x):K]\), so
  each implies the other (using \(D - x = D - 1\cdot x\),
  Lemma~\ref{pa-lem:devissageDivisor_sub}).
\end{proof}

\begin{theorem}[Principal divisors have degree zero]
  \label{pa-thm:deg_divOf}
  \lean{AlgebraicGeometry.deg_divOf}
  \leanok
  \dcref{custom}

  \uses{pa-thm:mulEquivDivisorSheaf, pa-thm:chi_congr, pa-thm:chi_divisorSheaf, pa-thm:curveDivisor_deg_hom}
  For a nonzero rational function \(g \in K(X)^\times\),
  \[
    \deg_K(\mathrm{div}\, g) \;=\; 0 .
  \]
\end{theorem}
\begin{proof}
  \leanok
  Multiplication by \(g\) is an isomorphism of sheaves
  \(\mathcal O(0) \cong \mathcal O(0 - \mathrm{div}\, g)\)
  (Theorem~\ref{pa-thm:mulEquivDivisorSheaf}), so the two sheaves have equal Euler characteristic
  (Theorem~\ref{pa-thm:chi_congr}). Evaluating both with the closed ledger
  (Theorem~\ref{pa-thm:chi_divisorSheaf}):
  \[
    \chi(\struct{X/K}) + \deg_K 0 \;=\; \chi(\struct{X/K}) + \deg_K(0 - \mathrm{div}\, g).
  \]
  Cancelling \(\chi(\struct{X/K})\) and using \(\deg_K 0 = 0\) and
  \(\deg_K(-\mathrm{div}\, g) = -\deg_K(\mathrm{div}\, g)\)
  (Theorem~\ref{pa-thm:curveDivisor_deg_hom}) leaves \(-\deg_K(\mathrm{div}\, g) = 0\), i.e.\
  \(\deg_K(\mathrm{div}\, g) = 0\).
\end{proof}

\begin{theorem}[The Riemann inequality]
  \label{pa-thm:riemann_inequality}
  \lean{AlgebraicGeometry.riemann_inequality}
  \leanok
  \dcref{custom}

  \uses{pa-thm:chi_divisorSheaf, pa-def:h0_h1_chi}
  Under the standing finiteness hypotheses, for every Weil divisor \(D\),
  \[
    \deg_K D + \chi(\struct{X/K}) \;\le\; h^0(\mathcal O(D)).
  \]
\end{theorem}
\begin{proof}
  \leanok
  By the closed ledger (Theorem~\ref{pa-thm:chi_divisorSheaf}),
  \(\chi(\mathcal O(D)) = \chi(\struct{X/K}) + \deg_K D\). On the other hand
  \(\chi(\mathcal O(D)) = h^0(\mathcal O(D)) - h^1(\mathcal O(D)) \le h^0(\mathcal O(D))\), since
  \(h^1 \ge 0\). Combining the two gives \(\deg_K D + \chi(\struct{X/K}) \le h^0(\mathcal O(D))\).
\end{proof}

\begin{theorem}[Unbounded growth of \(h^0\) along multiples of a point]
  \label{pa-thm:h0_nsmul_point_unbounded}
  \lean{AlgebraicGeometry.h0_nsmul_point_unbounded}
  \leanok
  \dcref{custom}

  \uses{pa-thm:chi_divisorSheaf, pa-lem:curveDivisor_single_arith, pa-thm:residueDeg_pos, pa-def:h0_h1_chi}
  Under the standing finiteness hypotheses, for every closed point \(x\) and every bound \(N\)
  there is \(n \in \mathbb N\) with
  \[
    N \;\le\; h^0\bigl(\mathcal O(n \cdot x)\bigr).
  \]
  No rational point is needed: it suffices that \([\kappa(x):K] \ge 1\).
\end{theorem}
\begin{proof}
  \leanok
  For each \(n\), the closed ledger (Theorem~\ref{pa-thm:chi_divisorSheaf}) and the degree of a
  one-point divisor (Lemma~\ref{pa-lem:curveDivisor_single_arith}) give
  \(\chi(\mathcal O(n\cdot x)) = \chi(\struct{X/K}) + n\,[\kappa(x):K]\). The residue degree is
  positive, \([\kappa(x):K] \ge 1\) (Theorem~\ref{pa-thm:residueDeg_pos}), so
  \(\chi(\mathcal O(n\cdot x)) \ge \chi(\struct{X/K}) + n\). Since
  \(\chi(\mathcal O(n\cdot x)) \le h^0(\mathcal O(n\cdot x))\) (as \(h^1 \ge 0\)) and
  \(\chi(\struct{X/K}) = h^0(\struct{X/K}) - h^1(\struct{X/K}) \ge -h^1(\struct{X/K})\), we obtain
  \[
    h^0(\mathcal O(n\cdot x)) \;\ge\; n - h^1(\struct{X/K}).
  \]
  Taking \(n = N + h^1(\struct{X/K})\) yields \(h^0(\mathcal O(n\cdot x)) \ge N\).
\end{proof}

\section{Riemann--Roch on the challenge curve}

Finally, instantiate the ledger on the bundled challenge curve \(C\) over the field \(k\): smooth
of relative dimension one, proper, and geometrically irreducible. Properness supplies both
missing ingredients --- the finiteness of \(H^1(C, \struct C)\), which is the genus, and the
identification \(H^0(C, \struct C) \cong k\), which makes \(h^0(\struct C) = 1\). The conditional
ledger then reads as Riemann--Roch-lite: \(\chi(\mathcal O(D)) = 1 - g(C) + \deg_k D\).

\begin{definition}[The \(k\)-linear structure map on global sections]
  \label{pa-def:overAlgebraMapTopLinear}
  \lean{AlgebraicGeometry.Scheme.overAlgebraMapTopLinear}
  \leanok
  \dcref{custom}

  \uses{pa-def:overAlgebraMap, pa-def:moduleKSheaf}
  For a scheme \(X\) over \(\Spec k\), the structure map \(k \to \Gamma(X, \top)\) of
  Definition~\ref{pa-def:overAlgebraMap} is \(k\)-linear for the \(k\)-module structure it induces
  on \(\Gamma(X, \top)\) (Definition~\ref{pa-def:moduleKSheaf}).
\end{definition}
\begin{proof}
  \leanok
  The \(k\)-module structure on \(\Gamma(X, \top)\) is restriction of scalars along the very same
  ring map \(k \to \Gamma(X, \top)\), so additivity and \(k\)-homogeneity of that map are the
  defining module axioms.
\end{proof}

\begin{theorem}[The \(k\)-linear \(\Gamma(X, \top) \cong k\) seam]
  \label{pa-thm:overAlgebraMapTopEquiv}
  \lean{AlgebraicGeometry.Scheme.bijective_overAlgebraMap_top,
    AlgebraicGeometry.Scheme.overAlgebraMapTopEquiv}
  \leanok
  \dcref{custom}

  \uses{pa-def:overAlgebraMapTopLinear}
  Let \(X\) be a scheme over \(\Spec k\) whose structure morphism is bijective on global sections
  (the induced map \(\Gamma(\Spec k, \struct{\Spec k}) \to \Gamma(X, \struct X)\) is bijective).
  Then the structure map \(k \to \Gamma(X, \top)\) is a \(k\)-linear equivalence
  \[
    k \;\cong\; \Gamma(X, \top).
  \]
\end{theorem}
\begin{proof}
  \leanok
  The structure map \(k \to \Gamma(X, \top)\) factors as the composite of the inverse of the
  canonical isomorphism \(\Gamma(\Spec k, \struct{\Spec k}) \cong k\), the map on global sections
  induced by the structure morphism, and the identity restriction \(\Gamma(X, \top) \to
  \Gamma(X, \top)\). The first factor is bijective (an isomorphism), the last is the identity,
  and the middle is bijective by hypothesis; hence the composite is bijective. A bijective
  \(k\)-linear map (Definition~\ref{pa-def:overAlgebraMapTopLinear}) is a \(k\)-linear equivalence.
\end{proof}

\begin{theorem}[\(h^0 = 1\) when global sections are constant]
  \label{pa-thm:h0_moduleKSheaf_bijective}
  \lean{AlgebraicGeometry.Scheme.h0_moduleKSheaf_of_bijective,
    AlgebraicGeometry.moduleFinite_hModule_zero_of_bijective}
  \leanok
  \dcref{custom}

  \uses{pa-thm:overAlgebraMapTopEquiv, pa-def:moduleKSheafHZero, pa-def:h0_h1_chi}
  Let \(X\) be a scheme over \(\Spec k\) whose structure morphism is bijective on global
  sections. Then \(H^0(X, \struct{X/k})\) is one-dimensional over \(k\); in particular it is
  finite-dimensional and \(h^0(\struct{X/k}) = 1\).
\end{theorem}
\begin{proof}
  \leanok
  Degree-zero cohomology is global sections, \(H^0(X, \struct{X/k}) \cong \Gamma(X, \top)\)
  (Definition~\ref{pa-def:moduleKSheafHZero}), and the seam identifies \(\Gamma(X, \top) \cong k\)
  (Theorem~\ref{pa-thm:overAlgebraMapTopEquiv}). Thus \(H^0(X, \struct{X/k}) \cong k\) is
  one-dimensional, so it is finite-dimensional and \(h^0(\struct{X/k}) = \dim_k k = 1\).
\end{proof}

\begin{theorem}[Global sections of the challenge curve]
  \label{pa-thm:h0_moduleKSheaf_curve}
  \lean{AlgebraicGeometry.moduleFinite_hModule_zero, AlgebraicGeometry.h0_moduleKSheaf}
  \leanok
  \dcref{custom}

  \uses{pa-thm:h0_moduleKSheaf_bijective, pa-thm:sections_k}
  For the challenge curve \(C\) --- proper, smooth of relative dimension one and geometrically
  irreducible over \(k\) --- the cohomology \(H^0(C, \struct C)\) is a one-dimensional (hence
  finite-dimensional) \(k\)-vector space, and \(h^0(\struct C) = 1\).
\end{theorem}
\begin{proof}
  \leanok
  Being smooth, \(C\) is geometrically reduced, and together with geometric irreducibility this
  makes \(C\) geometrically integral; being also proper, its structure morphism is bijective on
  global sections (Theorem~\ref{pa-thm:sections_k}). Apply
  Theorem~\ref{pa-thm:h0_moduleKSheaf_bijective}.
\end{proof}

\begin{theorem}[Euler characteristic of the structure sheaf of the curve]
  \label{pa-thm:chi_moduleKSheaf_curve}
  \lean{AlgebraicGeometry.chi_moduleKSheaf}
  \leanok
  \dcref{custom}

  \uses{pa-thm:h0_moduleKSheaf_curve, pa-def:genus, pa-def:h0_h1_chi}
  For the challenge curve \(C\),
  \[
    \chi(\struct C) \;=\; 1 - g(C),
  \]
  where \(g(C) = \dim_k H^1(C, \struct C)\) is the genus.
\end{theorem}
\begin{proof}
  \leanok
  By Theorem~\ref{pa-thm:h0_moduleKSheaf_curve}, \(h^0(\struct C) = 1\). By definition of the genus
  (Definition~\ref{pa-def:genus}), \(h^1(\struct C) = \dim_k H^1(C, \struct C) = g(C)\). Therefore
  \(\chi(\struct C) = h^0(\struct C) - h^1(\struct C) = 1 - g(C)\).
\end{proof}

\begin{theorem}[Riemann--Roch-lite for the challenge curve]
  \label{pa-thm:chi_divisorSheaf_curve}
  \lean{AlgebraicGeometry.chi_divisorSheaf_curve}
  \leanok
  \dcref{custom}

  \uses{pa-thm:chi_divisorSheaf, pa-thm:chi_moduleKSheaf_curve, pa-thm:h0_moduleKSheaf_curve,
    pa-thm:moduleFinite_hModule_one}
  For every Weil divisor \(D\) on the challenge curve \(C\),
  \[
    \chi(\mathcal O(D)) \;=\; 1 - g(C) + \deg_k D .
  \]
\end{theorem}
\begin{proof}
  \leanok
  The two standing finiteness hypotheses are met on the curve: \(H^0(C, \struct C)\) is
  finite-dimensional (Theorem~\ref{pa-thm:h0_moduleKSheaf_curve}) and \(H^1(C, \struct C)\) is
  finite-dimensional (Theorem~\ref{pa-thm:moduleFinite_hModule_one}). So the closed ledger applies
  (Theorem~\ref{pa-thm:chi_divisorSheaf}): \(\chi(\mathcal O(D)) = \chi(\struct C) + \deg_k D\).
  Substituting \(\chi(\struct C) = 1 - g(C)\) (Theorem~\ref{pa-thm:chi_moduleKSheaf_curve}) gives
  the claim.
\end{proof}

\begin{theorem}[The Riemann inequality for the challenge curve]
  \label{pa-thm:riemann_inequality_curve}
  \lean{AlgebraicGeometry.riemann_inequality_curve}
  \leanok
  \source{papaioannou-algebraic-rr:page-0007}

  \uses{pa-thm:riemann_inequality, pa-thm:chi_moduleKSheaf_curve, pa-thm:h0_moduleKSheaf_curve,
    pa-thm:moduleFinite_hModule_one}
  \dcref{custom}

  For every Weil divisor \(D\) on the challenge curve \(C\),
  \[
    \deg_k D + 1 - g(C) \;\le\; h^0(\mathcal O(D)).
  \]
\end{theorem}
\begin{proof}
  \leanok
  The standing finiteness hypotheses hold on \(C\) (Theorem~\ref{pa-thm:h0_moduleKSheaf_curve} and
  Theorem~\ref{pa-thm:moduleFinite_hModule_one}), so the general inequality
  \(\deg_k D + \chi(\struct C) \le h^0(\mathcal O(D))\) applies
  (Theorem~\ref{pa-thm:riemann_inequality}). Substituting \(\chi(\struct C) = 1 - g(C)\)
  (Theorem~\ref{pa-thm:chi_moduleKSheaf_curve}) gives the stated form. This is the classical Riemann
  inequality \(\dim L(D) \ge \deg D + 1 - g\), read off the genus as the supremum
  \(g = \sup_D\{\deg D - \dim L(D) + 1\}\).
\end{proof}

\begin{theorem}[Unbounded \(h^0\) along a closed point of the curve]
  \label{pa-thm:h0_nsmul_point_unbounded_curve}
  \lean{AlgebraicGeometry.h0_nsmul_point_unbounded_curve}
  \leanok
  \dcref{custom}

  \uses{pa-thm:h0_nsmul_point_unbounded, pa-thm:h0_moduleKSheaf_curve, pa-thm:moduleFinite_hModule_one}
  For every closed point \(x\) of the challenge curve \(C\) and every bound \(N\), there is
  \(n \in \mathbb N\) with \(N \le h^0(\mathcal O(n \cdot x))\). No rational point is needed.
\end{theorem}
\begin{proof}
  \leanok
  The standing finiteness hypotheses hold on \(C\) (Theorem~\ref{pa-thm:h0_moduleKSheaf_curve} and
  Theorem~\ref{pa-thm:moduleFinite_hModule_one}), so the general unboundedness
  (Theorem~\ref{pa-thm:h0_nsmul_point_unbounded}) applies verbatim.
\end{proof}

\section{The chart colength dictionary}
\label{pa-sec:chart_colength}

This section proves the dictionary between the algebra of an affine chart of the curve and
the divisor data at its closed points. Let \(V\) be an affine open of \(X\) containing the
generic point \(\eta\), and write \(B = \Gamma(X, V)\), a \(K\)-algebra through the
structure morphism (\ref{pa-def:overAlgebraMap}); all dimensions over \(K\) below are taken
for that structure. Under the isomorphism \(V \cong \Spec B\), a point \(x \in V\)
corresponds to a prime ideal \(\mathfrak p_x \subseteq B\), and conversely a prime
\(\mathfrak q\) of \(B\) determines a point \(x_{\mathfrak q} \in V\); the first four
lemmas record this correspondence between closed points of \(V\) and nonzero primes of
\(B\). The section ring \(B\) is then shown to be a Dedekind domain for \emph{every}
chart, and the colength \(\dim_K B/(f)\) of a nonzero section \(f\) is computed in
divisor terms: the primes dividing \((f)\) are the primes of the closed points where
\(f\) vanishes, the multiplicity of \(\mathfrak p_x\) in \((f)\) is the order of
vanishing at \(x\) of the rational function determined by \(f\), and the residue
\(B/\mathfrak p_x\) is the residue field \(\kappa(x)\). The engine behind the colength
formula is the localized Dedekind colength of
Section~\ref{pa-sec:localized_colength}.

\begin{lemma}[The point of a prime lies in the chart]
  \label{pa-lem:fromSpec_base_mem}
  \lean{AlgebraicGeometry.IsAffineOpen.fromSpec_base_mem}
  \leanok
  \dcref{custom}

  Let \(X\) be a scheme and \(V\) an affine open. For every prime \(\mathfrak q\) of
  \(B = \Gamma(X, V)\), the point \(x_{\mathfrak q}\) lies in \(V\).
\end{lemma}
\begin{proof}
  \leanok
  The canonical morphism \(\Spec B \to X\) is an open immersion with image \(V\), and
  \(x_{\mathfrak q}\) is by definition the image of the point \(\mathfrak q\) under it.
\end{proof}

\begin{lemma}[The correspondences are mutually inverse]
  \label{pa-lem:primeIdealOf_fromSpec_base}
  \lean{AlgebraicGeometry.IsAffineOpen.primeIdealOf_fromSpec_base}
  \leanok
  \dcref{custom}

  \uses{pa-lem:fromSpec_base_mem}
  Let \(X\) be a scheme and \(V\) an affine open. For every prime \(\mathfrak q\) of
  \(B\),
  \[
    \mathfrak p_{x_{\mathfrak q}} \;=\; \mathfrak q .
  \]
\end{lemma}
\begin{proof}
  \leanok
  The assignment \(x \mapsto \mathfrak p_x\) is defined so that the point of \(V\)
  determined by \(\mathfrak p_x\) is \(x\) again; since the morphism
  \(\Spec B \to X\) is an open immersion, it is injective on points, and the two
  primes \(\mathfrak p_{x_{\mathfrak q}}\) and \(\mathfrak q\), having the same image
  \(x_{\mathfrak q}\), coincide.
\end{proof}

\begin{lemma}[Nonzero prime at a non-generic point]
  \label{pa-lem:primeIdealOf_ne_bot}
  \lean{AlgebraicGeometry.IsAffineOpen.primeIdealOf_ne_bot}
  \leanok
  \dcref{custom}

  Let \(X\) be an integral scheme and \(V\) an affine open. For a point \(x \in V\)
  which is not the generic point of \(X\), the prime \(\mathfrak p_x\) is nonzero.
\end{lemma}
\begin{proof}
  \leanok
  Suppose \(\mathfrak p_x = 0\). Since \(X\) is integral and \(V\) is a nonempty open,
  \(B\) is a domain, so the zero ideal is the generic point of \(\Spec B\). An open
  immersion between irreducible schemes with dense image carries the generic point to
  the generic point, and \(V\) is dense in the irreducible \(X\); hence \(x\), the
  image of the generic point of \(\Spec B\), is the generic point of \(X\) ---
  contrary to hypothesis.
\end{proof}

\begin{lemma}[The point of a nonzero prime is non-generic]
  \label{pa-lem:fromSpec_base_ne_genericPoint}
  \lean{AlgebraicGeometry.IsAffineOpen.fromSpec_base_ne_genericPoint}
  \leanok
  \dcref{custom}

  \uses{pa-lem:fromSpec_base_mem}
  Let \(X\) be an integral scheme and \(V\) an affine open. For a prime
  \(\mathfrak q \neq 0\) of \(B\), the point \(x_{\mathfrak q}\) is not the generic
  point of \(X\).
\end{lemma}
\begin{proof}
  \leanok
  If \(x_{\mathfrak q}\) were the generic point of \(X\), then --- as in
  \ref{pa-lem:primeIdealOf_ne_bot}, the open immersion \(\Spec B \to X\) matches generic
  points, and is injective --- the prime \(\mathfrak q\) would be the generic point of
  \(\Spec B\), which is the zero ideal of the domain \(B\); contradiction.
\end{proof}

\begin{theorem}[Finite type of the chart sections]
  \label{pa-thm:overAlgebraMap_finiteType}
  \lean{AlgebraicGeometry.Scheme.overAlgebraMap_finiteType}
  \leanok
  \dcref{custom}

  \uses{pa-def:overAlgebraMap}
  Let \(K\) be a commutative ring and \(X\) a scheme locally of finite type over
  \(\Spec K\). For every affine open \(V\) of \(X\), the algebra map
  \(K \to \Gamma(X, V)\) of \ref{pa-def:overAlgebraMap} is of finite type.
\end{theorem}
\begin{proof}
  \leanok
  The algebra map is the composite of the identification \(K \cong \Gamma(\Spec K)\)
  with the map on sections \(\Gamma(\Spec K) \to \Gamma(X, V)\) induced by the
  structure morphism between the affine opens \(V\) and \(\Spec K\). For a morphism
  locally of finite type, the induced map between the section rings of an affine open
  of the source contained in the preimage of an affine open of the target is a ring
  map of finite type; apply this to \(V\) and \(\Spec K\).
\end{proof}

\begin{theorem}[Dedekind sections on every affine chart]
  \label{pa-thm:isDedekindDomain_section}
  \lean{AlgebraicGeometry.isDedekindDomain_section}
  \leanok
  \dcref{custom}

  Let \(X\) be an integral scheme, smooth of relative dimension one and locally of
  finite type over the field \(K\), and let \(V\) be an affine open containing the
  generic point (equivalently, nonempty). Then \(B = \Gamma(X, V)\) is a Dedekind
  domain. (This upgrades \ref{pa-thm:relDim1_dedekind_charts}, which produces \emph{some}
  Dedekind chart around each point, to \emph{every} affine chart.)
\end{theorem}
\begin{proof}
  \leanok
  \uses{pa-thm:overAlgebraMap_finiteType, pa-thm:closedpoint_dvr_stalk,
    pa-lem:fromSpec_base_mem, pa-lem:primeIdealOf_fromSpec_base,
    pa-lem:fromSpec_base_ne_genericPoint}
  \(B\) is a domain: \(V\) is a nonempty open of the integral scheme \(X\). \(B\) is
  Noetherian: the algebra map \(K \to B\) is of finite type
  (\ref{pa-thm:overAlgebraMap_finiteType}), and a finite-type algebra over a field is
  Noetherian by Hilbert's basis theorem. It remains to see that the localization of
  \(B\) at each nonzero prime \(\mathfrak q\) is a discrete valuation ring; a
  Noetherian domain, not a field, with this property is Dedekind. The point
  \(x := x_{\mathfrak q}\) lies in \(V\) (\ref{pa-lem:fromSpec_base_mem}) and is not the
  generic point (\ref{pa-lem:fromSpec_base_ne_genericPoint}), so the stalk
  \(\struct{X,x}\) is a discrete valuation ring (\ref{pa-thm:closedpoint_dvr_stalk}).
  The stalk at a point of an affine open is the localization of the section ring at
  the corresponding prime, and \(\mathfrak p_x = \mathfrak q\)
  (\ref{pa-lem:primeIdealOf_fromSpec_base}); hence \(B_{\mathfrak q} \cong
  \struct{X,x}\) is a discrete valuation ring.
\end{proof}

\begin{theorem}[Finiteness of the chart residue]
  \label{pa-thm:moduleFinite_quotient_primeIdealOf}
  \lean{AlgebraicGeometry.moduleFinite_quotient_primeIdealOf}
  \leanok
  \dcref{custom}

  \uses{pa-def:overAlgebraMap}
  Let \(X\) be as in \ref{pa-thm:isDedekindDomain_section}, \(V\) an affine open, and
  \(x \in V\) a point other than the generic point. Then the residue
  \(B/\mathfrak p_x\) is a finite \(K\)-module.
\end{theorem}
\begin{proof}
  \leanok
  \uses{pa-thm:closedpoint_isClosed, pa-thm:overAlgebraMap_finiteType}
  The singleton \(\{x\}\) is closed in \(X\) (\ref{pa-thm:closedpoint_isClosed}), so
  \(x\) is a closed point of the chart \(V \cong \Spec B\) and \(\mathfrak p_x\) is a
  maximal ideal; \(B/\mathfrak p_x\) is a field. The composite
  \(K \to B \to B/\mathfrak p_x\) is of finite type
  (\ref{pa-thm:overAlgebraMap_finiteType} followed by a surjection), and a field which is
  a finite-type algebra over a field is a finite extension (Zariski's lemma). Hence
  \(B/\mathfrak p_x\) is a finite \(K\)-module.
\end{proof}

\begin{theorem}[The residue leg: the chart residue is the residue field]
  \label{pa-thm:finrank_quotient_primeIdealOf}
  \lean{AlgebraicGeometry.finrank_quotient_primeIdealOf}
  \leanok
  \dcref{custom}

  \uses{pa-def:residueDeg}
  Let \(X\) be an integral scheme smooth of relative dimension one over the field
  \(K\), \(V\) an affine open, and \(x \in V\) a point other than the generic point.
  Then the residue of the chart at \(\mathfrak p_x\) is the residue field of \(x\),
  \(K\)-linearly; in particular
  \[
    \dim_K\bigl(B/\mathfrak p_x\bigr) \;=\; [\kappa(x) : K] \;=\; \deg_K(x) .
  \]
\end{theorem}
\begin{proof}
  \leanok
  \uses{pa-thm:closedpoint_isClosed, pa-def:residueOverAlgebraMap}
  As in \ref{pa-thm:moduleFinite_quotient_primeIdealOf}, \(\mathfrak p_x\) is maximal
  (\ref{pa-thm:closedpoint_isClosed}). The stalk \(\struct{X,x}\) is the localization of
  \(B\) at \(\mathfrak p_x\), and the residue field of the localization of a ring at a
  maximal ideal is the quotient by that ideal: \(\kappa(x) =
  \struct{X,x}/\mathfrak m_x \cong B/\mathfrak p_x\), by the map induced by
  \(B \to \struct{X,x}\). This isomorphism is \(K\)-linear: the \(K\)-structure of
  \(\kappa(x)\) is the composite of the algebra map into the sections, the germ, and
  the residue map (\ref{pa-def:residueOverAlgebraMap}), and the germ of the algebra map
  through \(V\) agrees with the germ of the global algebra map. Taking dimensions
  gives \(\dim_K(B/\mathfrak p_x) = \dim_K \kappa(x) = \deg_K(x)\)
  (\ref{pa-def:residueDeg}).
\end{proof}

\begin{theorem}[The multiplicity leg: order of vanishing is the multiplicity]
  \label{pa-thm:toAdd_ordZ_eq_count_factors}
  \lean{AlgebraicGeometry.toAdd_ordZ_eq_count_factors}
  \leanok
  \dcref{custom}

  \uses{pa-def:scheme_ordZ}
  Let \(X\) be an integral scheme smooth of relative dimension one over the field
  \(K\), \(V\) an affine open containing the generic point \(\eta\), and assume
  \(B = \Gamma(X, V)\) is a Dedekind domain (\ref{pa-thm:isDedekindDomain_section}). Let
  \(x \in V\) be a point other than \(\eta\), let \(f \in B\) be nonzero, and let
  \(g \in K(X)^\times\) be a unit whose value is the germ of \(f\) at \(\eta\). Then
  the order of vanishing of \(g\) at \(x\) is the multiplicity of \(\mathfrak p_x\) in
  the prime factorization of \((f) \subseteq B\):
  \[
    \mathrm{ord}_x(g) \;=\; \mathrm{ord}_{\mathfrak p_x}(f) .
  \]
\end{theorem}
\begin{proof}
  \leanok
  \uses{pa-lem:primeIdealOf_ne_bot, pa-thm:closedpoint_dvr_stalk, pa-def:scheme_ord,
    pa-lem:germ_generic_eq_algebraMap_germ, pa-lem:adic_valuation_count,
    pa-thm:count_factors_span_localization}
  The prime \(\mathfrak p_x\) is nonzero (\ref{pa-lem:primeIdealOf_ne_bot}), and the
  stalk \(\struct{X,x}\) is the localization of \(B\) at \(\mathfrak p_x\) --- a
  discrete valuation ring (\ref{pa-thm:closedpoint_dvr_stalk}) with maximal ideal
  \(\mathfrak m_x\). The germ of \(f\) at \(x\) is its image under
  \(B \to \struct{X,x}\); it is nonzero, because on an integral scheme the germ map is
  injective on sections and \(f \neq 0\). By the germ/stalk seam
  (\ref{pa-lem:germ_generic_eq_algebraMap_germ}), \(g\) --- the germ of \(f\) at
  \(\eta\) --- is the image of the germ of \(f\) at \(x\) under
  \(\struct{X,x} \to K(X)\).

  The order valuation \(\mathrm{ord}_x\) is the \(\mathfrak m_x\)-adic valuation of
  the discrete valuation ring \(\struct{X,x}\) (\ref{pa-def:scheme_ord}), so its value on
  \(g\) is the \(\mathfrak m_x\)-adic valuation of the germ of \(f\) at \(x\), which
  is the multiplicity of \(\mathfrak m_x\) in the factorization of the principal
  ideal generated by that germ (\ref{pa-lem:adic_valuation_count}, applied in the
  Dedekind domain \(\struct{X,x}\)). Finally, multiplicities are preserved by
  localization at the prime (\ref{pa-thm:count_factors_span_localization}, with
  \(S = \struct{X,x}\), a Dedekind local ring which is the localization of \(B\) at
  \(\mathfrak p_x\)):
  \(\mathrm{ord}_{\mathfrak m_x}\bigl(f \cdot \struct{X,x}\bigr)
  = \mathrm{ord}_{\mathfrak p_x}(f)\). Chaining the identities gives
  \(\mathrm{ord}_x(g) = \mathrm{ord}_{\mathfrak p_x}(f)\).
\end{proof}

\begin{theorem}[Finiteness of the chart colength module]
  \label{pa-thm:moduleFinite_quotient_span_section}
  \lean{AlgebraicGeometry.moduleFinite_quotient_span_section}
  \leanok
  \dcref{custom}

  \uses{pa-def:overAlgebraMap}
  Let \(X\) be an integral scheme, smooth of relative dimension one and locally of
  finite type over the field \(K\), and \(V\) an affine open containing \(\eta\).
  For every nonzero \(f \in B = \Gamma(X, V)\), the quotient \(B/(f)\) is a finite
  \(K\)-module.
\end{theorem}
\begin{proof}
  \leanok
  \uses{pa-thm:isDedekindDomain_section, pa-lem:prime_factor_span, pa-lem:fromSpec_base_mem,
    pa-lem:primeIdealOf_fromSpec_base, pa-lem:fromSpec_base_ne_genericPoint,
    pa-thm:moduleFinite_quotient_primeIdealOf, pa-thm:moduleFinite_quotient_span}
  \(B\) is a Dedekind domain (\ref{pa-thm:isDedekindDomain_section}). Each prime factor
  \(P\) of \((f)\) is a nonzero prime (\ref{pa-lem:prime_factor_span}); its point
  \(x_P\) lies in \(V\) (\ref{pa-lem:fromSpec_base_mem}), is not the generic point
  (\ref{pa-lem:fromSpec_base_ne_genericPoint}), and satisfies
  \(\mathfrak p_{x_P} = P\) (\ref{pa-lem:primeIdealOf_fromSpec_base}); so the residue
  \(B/P\) is a finite \(K\)-module by
  \ref{pa-thm:moduleFinite_quotient_primeIdealOf}. The finiteness of the Dedekind
  colength module (\ref{pa-thm:moduleFinite_quotient_span}) concludes.
\end{proof}

\begin{theorem}[The chart colength dictionary]
  \label{pa-thm:finrank_quotient_span_section}
  \lean{AlgebraicGeometry.finrank_quotient_span_section}
  \leanok
  \dcref{custom}

  \uses{pa-def:scheme_ordZ, pa-def:residueDeg}
  Let \(X\) be an integral scheme, smooth of relative dimension one and locally of
  finite type over the field \(K\), and \(V\) an affine open containing \(\eta\). Let
  \(f \in B = \Gamma(X, V)\), let \(g \in K(X)^\times\) be a unit whose value is the
  germ of \(f\) at \(\eta\), and let \(S\) be a finite set of non-generic points of
  \(X\), contained in \(V\), such that at every non-generic point of \(V\) outside
  \(S\) the germ of \(f\) is a unit of the stalk. Then, in \(\Z\),
  \[
    \dim_K\bigl(B/(f)\bigr) \;=\;
      \sum_{x \in S} \mathrm{ord}_x(g) \cdot \deg_K(x) .
  \]
\end{theorem}
\begin{proof}
  \leanok
  \uses{pa-thm:isDedekindDomain_section, pa-thm:finrank_quotient_span_eq_sum_count,
    pa-thm:toAdd_ordZ_eq_count_factors, pa-thm:finrank_quotient_primeIdealOf,
    pa-thm:moduleFinite_quotient_primeIdealOf, pa-lem:prime_factor_span,
    pa-lem:primeIdealOf_ne_bot, pa-lem:fromSpec_base_mem, pa-lem:primeIdealOf_fromSpec_base,
    pa-lem:fromSpec_base_ne_genericPoint, pa-lem:ord_eq_one_of_mem_basicOpen,
    pa-lem:mem_pow_iff_le_count}
  \(B\) is a Dedekind domain (\ref{pa-thm:isDedekindDomain_section}), and \(f \neq 0\):
  its germ at \(\eta\) is the value of the unit \(g\). As in
  \ref{pa-thm:moduleFinite_quotient_span_section}, every prime factor of \((f)\) has
  finite residue over \(K\), so the localized Dedekind colength formula
  (\ref{pa-thm:finrank_quotient_span_eq_sum_count}) gives
  \[
    \dim_K\bigl(B/(f)\bigr) \;=\;
      \sum_{P \mid (f)} \mathrm{ord}_P(f) \cdot \dim_K(B/P) .
  \]
  It remains to identify the right side with the sum over \(S\). Note first that a
  section is a non-unit at a point \(z \in V\) exactly when it lies in
  \(\mathfrak p_z\): the stalk at \(z\) is the localization of \(B\) at
  \(\mathfrak p_z\), where the image of an element of \(B\) is a unit precisely when
  the element avoids the prime.

  \emph{Terms of the \(S\)-sum outside the factor set vanish.} If \(x \in S\) and the
  germ of \(f\) at \(x\) is a unit, then \(x\) lies in the basic open of \(V\) where
  \(f\) is invertible, so \(\mathrm{ord}_x(g) = 0\)
  (\ref{pa-lem:ord_eq_one_of_mem_basicOpen}) and the term vanishes. Otherwise
  \(f \in \mathfrak p_x\); then \(\mathfrak p_x\) is a prime factor of \((f)\): it is
  nonzero (\ref{pa-lem:primeIdealOf_ne_bot}), and membership \(f \in \mathfrak p_x =
  \mathfrak p_x^1\) forces multiplicity \(\ge 1\)
  (\ref{pa-lem:mem_pow_iff_le_count}).

  \emph{The assignment \(x \mapsto \mathfrak p_x\) matches the two index sets.} It is
  injective on \(V\), the inverse being \(\mathfrak q \mapsto x_{\mathfrak q}\)
  (\ref{pa-lem:primeIdealOf_fromSpec_base}). It is surjective onto the prime factors:
  a prime factor \(P\) of \((f)\) contains \(f\) (\ref{pa-lem:prime_factor_span}), its
  point \(x_P\) lies in \(V\) (\ref{pa-lem:fromSpec_base_mem}), is non-generic
  (\ref{pa-lem:fromSpec_base_ne_genericPoint}) and satisfies \(\mathfrak p_{x_P} = P\)
  (\ref{pa-lem:primeIdealOf_fromSpec_base}); and \(x_P \in S\), for otherwise the
  outside-\(S\) hypothesis would make the germ of \(f\) at \(x_P\) a unit, while
  \(f \in P = \mathfrak p_{x_P}\).

  \emph{The terms match.} For \(x \in S\) whose prime is a factor of \((f)\), the
  multiplicity leg gives \(\mathrm{ord}_x(g) = \mathrm{ord}_{\mathfrak p_x}(f)\)
  (\ref{pa-thm:toAdd_ordZ_eq_count_factors}) and the residue leg gives
  \(\dim_K(B/\mathfrak p_x) = \deg_K(x)\)
  (\ref{pa-thm:finrank_quotient_primeIdealOf}). Hence the sum over the prime factors of
  \((f)\) coincides, term by term, with the sum over \(S\) of the nonzero terms, and
  the formula follows.
\end{proof}

\section{The degree homomorphism on the \v Cech Picard group}

The degree of a Weil divisor is invariant under linear equivalence, so it descends from
\(\Div(X)\) to the \v Cech Picard group \(\Pic(X) \cong \Div(X)/\mathrm{Princ}(X)\)
(\ref{pa-thm:curveDivisor_picClass_eq_iff}). This section records the degree homomorphism
\(\deg_{\Pic} : \Pic(X) \to \Z\) on classes and its interface: the normalization \(\deg_{\Pic}
\struct X(D) = \deg_K D\), additivity, and the connection to the \(\chi\)-ledger
\(\chi(\mathcal O(D)) = \chi(\struct{X/K}) + \deg_{\Pic}\struct X(D)\). The same linear-equivalence
argument makes the Euler characteristic of a class well-defined, this time with no finiteness
hypothesis at all.

\begin{lemma}[Linearly equivalent divisors have equal degree]
  \label{pa-lem:deg_eq_deg_of_picClass_eq}
  \lean{AlgebraicGeometry.deg_eq_deg_of_picClass_eq}
  \leanok
  \dcref{custom}

  \uses{pa-def:curveDivisor_deg, pa-thm:curveDivisor_picClass_eq_iff, pa-thm:deg_divOf,
    pa-thm:curveDivisor_deg_hom}
  If two Weil divisors have the same \v Cech Picard class, \(\struct X(D) = \struct X(D')\), then
  \(\deg_K D = \deg_K D'\).
\end{lemma}
\begin{proof}
  \leanok
  Equal classes differ by a principal divisor (\ref{pa-thm:curveDivisor_picClass_eq_iff}): there is
  \(g \in K(X)^\times\) with \(D - D' = \mathrm{div}\, g\). Principal divisors have degree zero
  (Theorem~\ref{pa-thm:deg_divOf}), so \(\deg_K(D - D') = 0\); and \(\deg_K(D - D') = \deg_K D -
  \deg_K D'\) by the degree homomorphism (Theorem~\ref{pa-thm:curveDivisor_deg_hom}). Hence
  \(\deg_K D = \deg_K D'\).
\end{proof}

\begin{theorem}[The Euler characteristic of a class is well-defined]
  \label{pa-thm:chi_divisorSheaf_eq_of_picClass_eq}
  \lean{AlgebraicGeometry.chi_divisorSheaf_eq_of_picClass_eq}
  \leanok
  \dcref{custom}

  \uses{pa-def:divisorSheaf, pa-def:h0_h1_chi, pa-thm:curveDivisor_picClass_eq_iff,
    pa-thm:mulEquivDivisorSheaf, pa-thm:chi_congr}
  If two Weil divisors have the same \v Cech Picard class, \(\struct X(D) = \struct X(D')\), then
  their divisor sheaves have the same Euler characteristic, \(\chi(\mathcal O(D)) = \chi(\mathcal
  O(D'))\). No finiteness hypothesis is needed.
\end{theorem}
\begin{proof}
  \leanok
  Equal classes differ by a principal divisor (\ref{pa-thm:curveDivisor_picClass_eq_iff}): \(D -
  \mathrm{div}\, g = D'\) for some \(g \in K(X)^\times\). Multiplication by \(g\) is an isomorphism
  of divisor sheaves \(\mathcal O(D) \cong \mathcal O(D - \mathrm{div}\, g) = \mathcal O(D')\)
  (Theorem~\ref{pa-thm:mulEquivDivisorSheaf}), and isomorphic sheaves have equal ledger invariants
  (Theorem~\ref{pa-thm:chi_congr}), so \(\chi(\mathcal O(D)) = \chi(\mathcal O(D'))\). The argument
  uses only the isomorphism and its transport, never the finiteness of any cohomology.
\end{proof}

\begin{definition}[The degree homomorphism on the Picard group]
  \label{pa-def:classDeg}
  \lean{AlgebraicGeometry.classDeg, AlgebraicGeometry.classDeg_picClass}
  \leanok
  \dcref{custom}

  \uses{pa-lem:deg_eq_deg_of_picClass_eq, pa-def:curveDivisor_picClass, pa-def:curveDivisor_deg}
  The degree descends to a group homomorphism
  \[
    \deg_{\Pic} : \Pic(X) \to \Z, \qquad \deg_{\Pic}\struct X(D) = \deg_K D
    \quad (D \in \Div(X)),
  \]
  the \emph{degree of a Picard class}: the degree of any Weil divisor realizing the class. The
  displayed normalization \(\deg_{\Pic}\struct X(D) = \deg_K D\) holds for every Weil divisor
  \(D\), and is the defining property.
\end{definition}
\begin{proof}
  \leanok
  Every class \(L \in \Pic(X)\) is realized by a Weil divisor, \(L = \struct X(D)\)
  (\ref{pa-def:curveDivisor_picClass}); set \(\deg_{\Pic} L = \deg_K D\). This is independent of the
  chosen realization: if \(\struct X(D) = \struct X(D')\) then \(\deg_K D = \deg_K D'\)
  (Lemma~\ref{pa-lem:deg_eq_deg_of_picClass_eq}). Applying this to a divisor's own class gives the
  normalization \(\deg_{\Pic}\struct X(D) = \deg_K D\). It is a homomorphism: the trivial class is
  \(\struct X(0)\), of degree \(\deg_K 0 = 0\); and for classes realized by \(D, D'\) the product
  is \(\struct X(D)\cdot\struct X(D') = \struct X(D + D')\), whose degree is \(\deg_K(D + D') =
  \deg_K D + \deg_K D'\) (the class map and the degree are both homomorphisms,
  \ref{pa-def:curveDivisor_deg}).
\end{proof}

\begin{theorem}[Additivity of the class degree]
  \label{pa-thm:classDeg_mul}
  \lean{AlgebraicGeometry.classDeg_mul}
  \leanok
  \dcref{custom}

  \uses{pa-def:classDeg, pa-def:curveDivisor_picClass, pa-thm:curveDivisor_picClass_hom,
    pa-thm:curveDivisor_deg_hom}
  The degree of a product of Picard classes is the sum of the degrees: for \(L, L' \in \Pic(X)\),
  \[
    \deg_{\Pic}(L \cdot L') = \deg_{\Pic} L + \deg_{\Pic} L' .
  \]
\end{theorem}
\begin{proof}
  \leanok
  Realize the classes by Weil divisors, \(L = \struct X(D)\) and \(L' = \struct X(D')\)
  (\ref{pa-def:curveDivisor_picClass}). Their product is \(L \cdot L' = \struct X(D + D')\) by
  additivity of the class map (Theorem~\ref{pa-thm:curveDivisor_picClass_hom}), so by the
  normalization (\ref{pa-def:classDeg}) and additivity of the divisor degree
  (Theorem~\ref{pa-thm:curveDivisor_deg_hom}),
  \[
    \deg_{\Pic}(L \cdot L') = \deg_K(D + D') = \deg_K D + \deg_K D' = \deg_{\Pic} L + \deg_{\Pic}
    L'.
  \]
\end{proof}

\begin{theorem}[The \(\chi\)-ledger in class form]
  \label{pa-thm:chi_divisorSheaf_classDeg}
  \lean{AlgebraicGeometry.chi_divisorSheaf_classDeg}
  \leanok
  \dcref{custom}

  \uses{pa-def:classDeg, pa-thm:chi_divisorSheaf}
  Under the standing finiteness hypotheses, for every Weil divisor \(D\),
  \[
    \chi(\mathcal O(D)) = \chi(\struct{X/K}) + \deg_{\Pic}\struct X(D) .
  \]
\end{theorem}
\begin{proof}
  \leanok
  By the normalization \(\deg_{\Pic}\struct X(D) = \deg_K D\) (\ref{pa-def:classDeg}), the claim is the
  closed ledger \(\chi(\mathcal O(D)) = \chi(\struct{X/K}) + \deg_K D\)
  (Theorem~\ref{pa-thm:chi_divisorSheaf}).
\end{proof}

\begin{theorem}[The point-product class map has the same degree]
  \label{pa-thm:classDeg_divisorClass}
  \lean{AlgebraicGeometry.classDeg_divisorClass}
  \leanok
  \dcref{custom}

  \uses{pa-def:classDeg, pa-def:divisorClass, pa-thm:divisorClass_eq_picClass}
  For every Weil divisor \(D\), the degree of the point-product class equals the degree of \(D\):
  \[
    \deg_{\Pic}\bigl(\mathrm{divisorClass}(D)\bigr) = \deg_K D .
  \]
\end{theorem}
\begin{proof}
  \leanok
  The point-product class map and the anchored class map coincide,
  \(\mathrm{divisorClass}(D) = \struct X(D)\) (Theorem~\ref{pa-thm:divisorClass_eq_picClass}), so the
  degree of \(\mathrm{divisorClass}(D)\) is the degree of \(\struct X(D)\), which is \(\deg_K D\) by
  the normalization (\ref{pa-def:classDeg}).
\end{proof}

\begin{theorem}[Riemann--Roch-lite in class form on the challenge curve]
  \label{pa-thm:chi_divisorSheaf_classDeg_curve}
  \lean{AlgebraicGeometry.chi_divisorSheaf_classDeg_curve}
  \leanok
  \dcref{custom}

  \uses{pa-def:classDeg, pa-thm:chi_divisorSheaf_curve}
  For every Weil divisor \(D\) on the challenge curve \(C\),
  \[
    \chi(\mathcal O(D)) = 1 - g(C) + \deg_{\Pic}\struct X(D) .
  \]
  The standing finiteness hypotheses are discharged on the curve, so there is no side condition.
\end{theorem}
\begin{proof}
  \leanok
  On the curve the closed ledger reads \(\chi(\mathcal O(D)) = 1 - g(C) + \deg_k D\)
  (Theorem~\ref{pa-thm:chi_divisorSheaf_curve}). Substituting the normalization \(\deg_{\Pic}\struct
  X(D) = \deg_k D\) (\ref{pa-def:classDeg}) gives the claim.
\end{proof}

\section{Cohomology of a Picard class is well-defined}
\label{pa-sec:classCohomology}

Throughout this section \(X\) is a curve bundle over a field \(K\): integral, smooth of relative
dimension one, and quasi-compact over \(\Spec K\). The Euler characteristic of a divisor sheaf
is a \v Cech Picard class invariant (\ref{pa-thm:chi_divisorSheaf_eq_of_picClass_eq}); the same
holds \emph{degree by degree} for \(h^0\) and \(h^1\), and for the vanishing and
finite-dimensionality of the individual cohomology modules. All of these transport along one
isomorphism of divisor sheaves, so the arguments are one line each on the ledger API.

The common input is: if two Weil divisors \(D, D'\) have equal class, \([\struct X(D)] =
[\struct X(D')]\), then the extraction lemma \ref{pa-thm:curveDivisor_picClass_eq_iff} writes \(D -
D' = \operatorname{div} g\) for a rational function \(g\), and multiplication by \(g\)
(\ref{pa-thm:mulEquivDivisorSheaf}) realises an isomorphism of sheaves of \(K\)-modules \(\struct
X(D) \cong \struct X(D')\). Every statement below is transport along this isomorphism.

\begin{theorem}[\(h^0\) and \(h^1\) are Picard-class invariants]
  \label{pa-thm:h0h1_divisorSheaf_class}
  \lean{AlgebraicGeometry.h0_divisorSheaf_eq_of_picClass_eq,
    AlgebraicGeometry.h1_divisorSheaf_eq_of_picClass_eq}
  \leanok
  \dcref{custom}

  \uses{pa-def:divisorSheaf, pa-def:h0_h1_chi, pa-thm:curveDivisor_picClass_eq_iff,
    pa-thm:mulEquivDivisorSheaf, pa-thm:chi_congr}
  Let \(D, D'\) be Weil divisors on \(X\) with equal class, \([\struct X(D)] = [\struct
  X(D')]\). Then
  \[
    h^0\bigl(\struct X(D)\bigr) = h^0\bigl(\struct X(D')\bigr), \qquad
    h^1\bigl(\struct X(D)\bigr) = h^1\bigl(\struct X(D')\bigr).
  \]
  No finiteness hypothesis is needed. This is the witness-independence that makes a per-class
  quantifier ``for all \(D\) with \([\struct X(D)] = \lambda\)'' well-posed.
\end{theorem}
\begin{proof}
  \leanok
  The class isomorphism \(\struct X(D) \cong \struct X(D')\) described above induces, in each
  fixed degree \(n\), an isomorphism of cohomology modules \(H^n(\struct X(D)) \cong
  H^n(\struct X(D'))\); passing to dimensions (\(h^0\)-, \(h^1\)-congruence,
  \ref{pa-thm:chi_congr}) gives the two equalities.
\end{proof}

\begin{theorem}[Vanishing of \(H^1\) is a Picard-class invariant]
  \label{pa-thm:subsingleton_hModule_one_class}
  \lean{AlgebraicGeometry.subsingleton_hModule_one_of_picClass_eq}
  \leanok
  \dcref{custom}

  \uses{pa-def:divisorSheaf, pa-def:HModule, pa-thm:curveDivisor_picClass_eq_iff,
    pa-thm:mulEquivDivisorSheaf, pa-def:HModule_mapEquiv}
  Let \(D, D'\) be Weil divisors on \(X\) with \([\struct X(D)] = [\struct X(D')]\). If the
  degree-one cohomology \(H^1(\struct X(D))\) vanishes (in the strong sense of being a
  subsingleton), then so does \(H^1(\struct X(D'))\).
\end{theorem}
\begin{proof}
  \leanok
  The class isomorphism \(\struct X(D) \cong \struct X(D')\) induces a \(K\)-linear equivalence
  \(H^1(\struct X(D)) \simeq H^1(\struct X(D'))\) (\ref{pa-def:HModule_mapEquiv}), and being a
  subsingleton transports across an equivalence. This is the form the consumers need: it
  survives isomorphisms and does not require finite-dimensionality, unlike an equality
  \(h^1 = 0\).
\end{proof}

\begin{theorem}[Finite-dimensionality of cohomology is a Picard-class invariant]
  \label{pa-thm:moduleFinite_hModule_divisorSheaf_class}
  \lean{AlgebraicGeometry.moduleFinite_hModule_divisorSheaf_zero_of_picClass_eq,
    AlgebraicGeometry.moduleFinite_hModule_divisorSheaf_one_of_picClass_eq}
  \leanok
  \dcref{custom}

  \uses{pa-def:divisorSheaf, pa-def:HModule, pa-thm:curveDivisor_picClass_eq_iff,
    pa-thm:mulEquivDivisorSheaf, pa-def:HModule_mapEquiv}
  Let \(D, D'\) be Weil divisors on \(X\) with \([\struct X(D)] = [\struct X(D')]\). If
  \(H^0(\struct X(D))\) (respectively \(H^1(\struct X(D))\)) is a finite-dimensional
  \(K\)-module, then so is \(H^0(\struct X(D'))\) (respectively \(H^1(\struct X(D'))\)).
\end{theorem}
\begin{proof}
  \leanok
  The class isomorphism \(\struct X(D) \cong \struct X(D')\) induces, in each degree, a
  \(K\)-linear equivalence of cohomology modules (\ref{pa-def:HModule_mapEquiv}); finite
  dimensionality transports across a linear equivalence.
\end{proof}

\section{The degree homomorphism over an extension field}

Everything above applies to the challenge curve \(C\) over its base field \(k\). We now record
that the same degree apparatus is available over \emph{every} field extension \(K/k\), on the
base-changed curve \(C_K := (C \otimes \Spec K)_{\mathrm{left}}\), regarded as a \(K\)-scheme via
the second projection \(\mathrm{pr}_2 \colon C_K \to \Spec K\). The point is that the geometric
half of the standing hypotheses --- proper, smooth of relative dimension one, geometrically
irreducible --- transports along base change, from which integrality and the finiteness of
\(H^0\) and \(H^1\) re-derive over \(K\); so the degree homomorphism \(\deg_{\Pic}\) is defined on
\(\Pic(C_K)\) with no further input.

\begin{theorem}[The base-changed curve is a smooth proper curve over \(K\)]
  \label{pa-thm:baseChangeCurve_bundle}
  \lean{AlgebraicGeometry.instOverBaseChange, AlgebraicGeometry.baseChange_over_eq_snd_left,
    AlgebraicGeometry.instSmoothOfRelativeDimensionSndLeft, AlgebraicGeometry.instIsProperSndLeft,
    AlgebraicGeometry.instGeometricallyIrreducibleSndLeft,
    AlgebraicGeometry.instSmoothOfRelativeDimensionBaseChange,
    AlgebraicGeometry.instIsProperBaseChange, AlgebraicGeometry.instQuasiCompactBaseChange}
  \leanok
  \dcref{custom}
  \uses{pa-def:relCurve}

  Let \(K/k\) be a field extension. The base-changed curve \(C_K\), as a scheme over \(\Spec K\)
  via the second projection, is smooth of relative dimension one, proper, quasi-compact, and
  geometrically irreducible over \(K\).
\end{theorem}
\begin{proof}
  \leanok
  The second projection \(C_K \to \Spec K\) is the base change of the structure morphism \(C \to
  \Spec k\) along \(\Spec K \to \Spec k\) (the curve product exhibits \(C_K\) as a fibre product,
  as for the relative curve of Definition~\ref{pa-def:relCurve}). Smoothness of relative dimension
  one, properness, and geometric irreducibility are each stable under base change, so all three
  transport from \(C/k\) to \(C_K/K\); properness gives quasi-compactness.
\end{proof}

\begin{theorem}[The base-changed curve is integral]
  \label{pa-thm:baseChangeCurve_integral}
  \lean{AlgebraicGeometry.instIsIntegralBaseChange}
  \leanok
  \dcref{custom}
  \uses{pa-thm:baseChangeCurve_bundle, pa-prop:curve_integral}

  For every field extension \(K/k\), the base-changed curve \(C_K\) is an integral scheme.
\end{theorem}
\begin{proof}
  \leanok
  Being smooth, \(C_K\) is geometrically reduced over \(K\); it is geometrically irreducible over
  \(K\) by Theorem~\ref{pa-thm:baseChangeCurve_bundle}. A geometrically reduced and geometrically
  irreducible scheme over a field is integral (Proposition~\ref{pa-prop:curve_integral}). This is
  the substance of the base change: geometric irreducibility and geometric reducedness of \(C/k\)
  yield integrality of \(C \otimes_k K\) for \emph{every} field extension \(K\), a property the
  bare fibre product \(C \otimes_k K\) need not have without them.
\end{proof}

\begin{theorem}[Finiteness of the cohomology of the base-changed curve]
  \label{pa-thm:baseChangeCurve_cohomFinite}
  \lean{AlgebraicGeometry.instModuleFiniteHModuleZeroBaseChange,
    AlgebraicGeometry.instModuleFiniteHModuleOneBaseChange}
  \leanok
  \dcref{custom}
  \uses{pa-thm:baseChangeCurve_bundle, pa-thm:h0_moduleKSheaf_curve, pa-thm:moduleFinite_hModule_one}

  For every field extension \(K/k\), both \(H^0(C_K, \struct{C_K/K})\) and \(H^1(C_K,
  \struct{C_K/K})\) are finite-dimensional \(K\)-vector spaces.
\end{theorem}
\begin{proof}
  \leanok
  By Theorem~\ref{pa-thm:baseChangeCurve_bundle}, \(C_K\) is a smooth curve over \(K\) --- proper,
  smooth of relative dimension one, geometrically irreducible. The curve-level finiteness results
  therefore apply over \(K\): \(H^0\) is finite-dimensional
  (Theorem~\ref{pa-thm:h0_moduleKSheaf_curve}) and \(H^1\) is finite-dimensional
  (Theorem~\ref{pa-thm:moduleFinite_hModule_one}).
\end{proof}

\begin{remark}[The degree homomorphism over every extension field]
  \label{pa-rem:classDeg_baseChange}
  \dcref{custom}
  \uses{pa-thm:baseChangeCurve_bundle, pa-thm:baseChangeCurve_integral,
    pa-thm:baseChangeCurve_cohomFinite, pa-def:classDeg}
  Theorems~\ref{pa-thm:baseChangeCurve_bundle}, \ref{pa-thm:baseChangeCurve_integral} and
  \ref{pa-thm:baseChangeCurve_cohomFinite} discharge every hypothesis of the degree stack for \(C_K\)
  over \(K\). Hence the degree homomorphism \(\deg_{\Pic} \colon \Pic(C_K) \to \Z\)
  (Definition~\ref{pa-def:classDeg}) is defined on the base-changed curve, for every field extension
  \(K/k\).
\end{remark}

\section{Invariance of the degree under the base-field transition}
\label{pa-sec:degree_base_field_invariance}

Let \(k\) be the base field of the challenge curve \(C\) (proper, smooth of relative
dimension one, geometrically irreducible), and let \(\varphi : K_1 \to K_2\) be a
\(k\)-algebra homomorphism between field extensions of \(k\), with base-field transition
\(\pi : C_{K_2} \to C_{K_1}\)
(Chapter~\ref{pa-chap:Cohomology}, Section~\ref{pa-sec:base_field_transition}). Both ends carry
the full degree apparatus (\ref{pa-rem:classDeg_baseChange}). This section proves that the
degree of a \v Cech Picard class is invariant under the transition,
\[
  \deg_{\Pic}\bigl(\pi^{*}L\bigr) \;=\; \deg_{\Pic}(L)
  \qquad \bigl(L \in \Pic(C_{K_1})\bigr),
\]
and descends the degree to the relative Picard group over a field. The invariance is the
algebraic form of the classical fact that the degree of a divisor does not change under
extension of the base field: a closed point \(x'\) of \(C_{K_1}\) pulls back to the fibre
divisor \(\pi^{*}(x') = \sum_{x'' \mapsto x'} e_{x''} \cdot x''\), and the multiplicities
and residue degrees rebalance so that
\(\sum e_{x''}\,[\kappa(x'') : K_2] = [\kappa(x') : K_1]\) --- the \emph{fiber-degree
identity} proved below by comparing the colength of one chart algebra at both ends of
the transition.

\begin{lemma}[The degree of a power of a class]
  \label{pa-lem:classDeg_zpow}
  \lean{AlgebraicGeometry.classDeg_zpow}
  \leanok
  \dcref{custom}

  \uses{pa-def:classDeg, pa-thm:classDeg_mul}
  Under the standing hypotheses of \ref{pa-def:classDeg}, for every class \(L\) and every
  integer \(n \in \Z\),
  \[
    \deg_{\Pic}\bigl(L^{\,n}\bigr) \;=\; n \cdot \deg_{\Pic}(L).
  \]
\end{lemma}
\begin{proof}
  \leanok
  The trivial class has degree zero: \(1 = \struct X(0)\) and \(\deg_K 0 = 0\)
  (\ref{pa-def:classDeg}). By additivity (\ref{pa-thm:classDeg_mul}),
  \(\deg_{\Pic}(L) + \deg_{\Pic}(L^{-1}) = \deg_{\Pic}(1) = 0\), so inversion negates
  the degree. Induct on \(n\) in both directions from \(n = 0\), using additivity at
  each step.
\end{proof}

\begin{theorem}[The fiber-degree identity]
  \label{pa-thm:sum_ordZ_residueDeg_baseFieldTransition}
  \lean{AlgebraicGeometry.sum_ordZ_residueDeg_baseFieldTransition}
  \leanok
  \dcref{custom}

  \uses{pa-def:scheme_ordZ, pa-def:residueDeg, pa-lem:baseFieldTransition_genericPoint,
    pa-def:functionFieldMap}
  Let \(V\) be an affine chart of \(C_{K_1}\) containing the generic point, let
  \(t \in \Gamma(C_{K_1}, V)\) be a section whose germ at \(\eta\) is the value of a
  unit \(g \in K(C_{K_1})^{\times}\) and which is a unit at every point of \(V\) other
  than a fixed non-generic point \(x' \in V\), and let \(S\) be a finite set of
  non-generic points of \(C_{K_2}\), contained in \(\pi^{-1}V\), outside which the
  pulled-back section \(\pi^{\sharp}t\) is a unit on \(\pi^{-1}V\). Then, in \(\Z\),
  \[
    \sum_{x'' \in S} \mathrm{ord}_{x''}\bigl(\pi^{\sharp}g\bigr) \cdot
      \bigl[\kappa(x'') : K_2\bigr]
    \;=\; \mathrm{ord}_{x'}(g) \cdot \bigl[\kappa(x') : K_1\bigr].
  \]
  (Mathematically, the left side is the sum over the fibre of \(x'\): the hypotheses
  force \(S\) to exhaust the non-unit locus of \(\pi^{\sharp}t\) on the chart, which is
  contained in \(\pi^{-1}(x')\) since unit germs pull back to unit germs; points of
  \(S\) outside that locus contribute order zero. The statement is parameterised by
  such a finite \(S\) exactly so that the consumer may take for \(S\) the support of a
  divisor already in hand.)
\end{theorem}
\begin{proof}
  \leanok
  \uses{pa-thm:finrank_quotient_span_section, pa-def:transitionSectionsBaseChange,
    pa-lem:transitionSectionsBaseChange_tmul, pa-lem:baseFieldTransition_flat_affine_surjective,
    pa-lem:functionFieldMap_germ, pa-thm:baseChangeCurve_integral, pa-thm:baseChangeCurve_bundle}
  Write \(B_1 = \Gamma(C_{K_1}, V)\) and
  \(B_2 = \Gamma(C_{K_2}, \pi^{-1}V)\). The chart \(\pi^{-1}V\) is affine (the
  transition is an affine morphism,
  \ref{pa-lem:baseFieldTransition_flat_affine_surjective}) and contains the generic point
  of \(C_{K_2}\) (the transition matches generic points,
  \ref{pa-lem:baseFieldTransition_genericPoint}); both curves carry the standing bundle
  hypotheses over their fields (\ref{pa-thm:baseChangeCurve_bundle},
  \ref{pa-thm:baseChangeCurve_integral}). By germ naturality
  (\ref{pa-lem:functionFieldMap_germ}), the germ of \(\pi^{\sharp}t\) at the generic point
  of \(C_{K_2}\) is the value of the unit \(\pi^{\sharp}g\).

  \emph{The two colength dictionaries.} The chart colength dictionary
  (\ref{pa-thm:finrank_quotient_span_section}) applies at both ends. Over \(K_2\), with
  the finite set \(S\) (whose members lie in \(\pi^{-1}V\), and outside which
  \(\pi^{\sharp}t\) is a unit, by hypothesis):
  \[
    \dim_{K_2}\bigl(B_2/(\pi^{\sharp}t)\bigr)
      \;=\; \sum_{x'' \in S} \mathrm{ord}_{x''}\bigl(\pi^{\sharp}g\bigr)
        \cdot \bigl[\kappa(x'') : K_2\bigr].
  \]
  Over \(K_1\), with the singleton \(S = \{x'\}\) (the section \(t\) is a unit at every
  other point of \(V\), by hypothesis):
  \[
    \dim_{K_1}\bigl(B_1/(t)\bigr)
      \;=\; \mathrm{ord}_{x'}(g) \cdot \bigl[\kappa(x') : K_1\bigr].
  \]

  \emph{The colength bridge.} Regard \(K_2\) as a \(K_1\)-algebra through \(\varphi\).
  The sections base change (\ref{pa-def:transitionSectionsBaseChange}) is a ring
  isomorphism \(B_1 \otimes_{K_1} K_2 \cong B_2\) carrying \(t \otimes 1\) to
  \(\pi^{\sharp}t\) (\ref{pa-lem:transitionSectionsBaseChange_tmul}) and commuting with
  the \(K_2\)-structures, hence carrying the principal ideal \((t \otimes 1)\) onto
  \((\pi^{\sharp}t)\) and inducing a \(K_2\)-linear isomorphism of the quotients. By
  right-exactness of the tensor product, tensoring the exact sequence
  \(B_1 \xrightarrow{\,\cdot t\,} B_1 \to B_1/(t) \to 0\) with \(K_2\) gives
  \(\bigl(B_1/(t)\bigr) \otimes_{K_1} K_2 \cong
  \bigl(B_1 \otimes_{K_1} K_2\bigr)/(t \otimes 1)\). Chaining,
  \[
    \bigl(B_1/(t)\bigr) \otimes_{K_1} K_2 \;\cong\; B_2/(\pi^{\sharp}t)
    \quad (K_2\text{-linearly}),
  \]
  and since extension of scalars along a field extension preserves dimension,
  \[
    \dim_{K_2}\bigl(B_2/(\pi^{\sharp}t)\bigr)
      \;=\; \dim_{K_2}\Bigl(\bigl(B_1/(t)\bigr) \otimes_{K_1} K_2\Bigr)
      \;=\; \dim_{K_1}\bigl(B_1/(t)\bigr).
  \]
  Chaining the three displays gives the identity.
\end{proof}

\begin{lemma}[Pulled local equations are regular on integral schemes]
  \label{pa-lem:pullbackEqn_regular}
  \lean{AlgebraicGeometry.Scheme.LocalEquations.pullbackEqn_germ_mem_nonZeroDivisors}
  \leanok
  \dcref{custom}

  \uses{pa-def:localEquations_pullback, pa-def:localEquations, pa-lem:functionFieldMap_germ}
  Let \(f : Y \to X\) be a morphism of \emph{integral} schemes with
  \(f(\eta_Y) = \eta_X\), and let \(E\) be a local-equation system on \(X\). Then the
  pulled-back equations of \(E\) (\ref{pa-def:localEquations_pullback}) are automatically
  regular: their germs at every point are nonzerodivisors. No flatness of \(f\) is
  needed --- integrality of both schemes discharges the regularity hypothesis of the
  pullback construction. (For the base-field transition this is strictly more general
  than the flat-pullback route: flatness of \(\pi\) is never consumed.)
\end{lemma}
\begin{proof}
  \leanok
  Fix a pulled equation \(f^{\sharp}(e)\), \(e\) an equation of \(E\) on a member
  \(U\), and a point \(z\) of \(f^{-1}U\). The stalks of an integral scheme are
  domains, so it suffices to show the germ of \(f^{\sharp}(e)\) at \(z\) is nonzero.
  Suppose it were zero. On an integral scheme, a section of a nonempty open is
  determined by its germ at any one point (the restriction maps between nonempty opens
  and the germ maps are injective, all sections embedding in the function field), so
  \(f^{\sharp}(e) = 0\), and in particular its germ at \(\eta_Y\) vanishes. But by
  germ naturality (\ref{pa-lem:functionFieldMap_germ}) that germ is the image under the
  \emph{injective} function-field map of the germ of \(e\) at \(\eta_X\), which is a
  nonzerodivisor of the field \(K(X)\) by regularity of \(E\)
  (\ref{pa-def:localEquations}), i.e.\ nonzero. Contradiction.
\end{proof}

\begin{definition}[The pulled tracked point system]
  \label{pa-def:pointTransitionEquations}
  \lean{AlgebraicGeometry.pointTransitionEquations,
    AlgebraicGeometry.pointTransitionEquations_eqn,
    AlgebraicGeometry.pointTransitionUnit, AlgebraicGeometry.coe_pointTransitionUnit}
  \leanok
  \dcref{custom}

  \uses{pa-def:pointEquations, pa-lem:pointUniformizerData, pa-def:localEquations_pullback,
    pa-lem:pullbackEqn_regular, pa-lem:baseFieldTransition_genericPoint, pa-def:functionFieldMap,
    pa-lem:uniformizer, pa-thm:baseChangeCurve_integral}
  Let \(x'\) be a non-generic point of \(C_{K_1}\), with tracked point equations at
  \(x'\) (\ref{pa-def:pointEquations}, from the tracked uniformizer data of
  \ref{pa-lem:pointUniformizerData}). The \emph{pulled tracked point system} is the
  pullback of this local-equation system along the transition \(\pi\)
  (\ref{pa-def:localEquations_pullback}), with regularity discharged by integrality of
  both curves (\ref{pa-lem:pullbackEqn_regular}, \ref{pa-thm:baseChangeCurve_integral},
  generic points matching by \ref{pa-lem:baseFieldTransition_genericPoint}). Its
  \emph{trivialising unit} is the pullback
  \(\pi^{\sharp}(t) \in K(C_{K_2})^{\times}\) of the chosen uniformizer \(t\) at
  \(x'\) (\ref{pa-lem:uniformizer}, \ref{pa-def:functionFieldMap}).
\end{definition}

\begin{lemma}[The trivialising elements of the pulled point presentation]
  \label{pa-lem:pointTransitionEquations_elem}
  \lean{AlgebraicGeometry.pointTransitionEquations_presentation_elem_of_ne,
    AlgebraicGeometry.pointTransitionEquations_presentation_elem_of_eq}
  \leanok
  \dcref{custom}

  \uses{pa-def:pointTransitionEquations, pa-lem:localEquations_presentation,
    pa-def:pointEquations, pa-lem:functionFieldMap_germ}
  The meromorphic presentation (\ref{pa-lem:localEquations_presentation}) of the pulled
  tracked point system has trivialising elements
  \[
    f_z \;=\;
    \begin{cases}
      \pi^{\sharp}(t), & \pi(z) = x', \\
      1, & \pi(z) \neq x',
    \end{cases}
  \]
  at each point \(z\) of \(C_{K_2}\).
\end{lemma}
\begin{proof}
  \leanok
  The trivialising element at \(z\) is the germ at \(\eta_{C_{K_2}}\) of the pulled
  equation indexed by \(z\), which is the \(\pi\)-pullback of the equation of the
  point system at \(\pi(z)\) (\ref{pa-def:pointTransitionEquations}). If
  \(\pi(z) \neq x'\), that equation is the constant \(1\)
  (\ref{pa-def:pointEquations}), and \(1\) pulls back to \(1\), with germ \(1\). If
  \(\pi(z) = x'\), the equation is the tracked section \(s\), whose germ at
  \(\eta_{C_{K_1}}\) is the chosen uniformizer \(t\) (\ref{pa-def:pointEquations});
  by germ naturality (\ref{pa-lem:functionFieldMap_germ}) the germ at
  \(\eta_{C_{K_2}}\) of \(\pi^{\sharp}s\) is \(\pi^{\sharp}(t)\).
\end{proof}

\begin{theorem}[E-iv: the degree is invariant under the base-field transition]
  \label{pa-thm:classDeg_cechPicMap_baseFieldTransition}
  \lean{AlgebraicGeometry.classDeg_cechPicMap_baseFieldTransition}
  \leanok
  \dcref{custom}

  \uses{pa-def:classDeg, pa-thm:baseChangeCurve_bundle, pa-thm:baseChangeCurve_integral,
    pa-thm:baseChangeCurve_cohomFinite, pa-def:overSpecMap}
  For every \(k\)-algebra homomorphism \(\varphi : K_1 \to K_2\) of field extensions of
  \(k\) and every \v Cech Picard class \(L\) on \(C_{K_1}\),
  \[
    \deg_{\Pic}\bigl(\pi^{*}L\bigr) \;=\; \deg_{\Pic}(L),
  \]
  where \(\pi^{*} : \Pic(C_{K_1}) \to \Pic(C_{K_2})\) is pullback of \v Cech Picard
  classes along the base-field transition. No finiteness or separability of the
  extensions is assumed.
\end{theorem}
\begin{proof}
  \leanok
  \uses{pa-def:curveDivisor_picClass, pa-thm:curveDivisor_picClass_hom, pa-lem:classDeg_zpow,
    pa-thm:curveDivisor_picClass_single, pa-thm:localEquations_picClass_pullback,
    pa-def:pointTransitionEquations, pa-lem:pointTransitionEquations_elem,
    pa-lem:localEquations_presentation, pa-def:presentationDivisor, pa-def:curveDivisor_single,
    pa-lem:curveDivisor_single_arith, pa-thm:sum_ordZ_residueDeg_baseFieldTransition,
    pa-lem:uniformizer, pa-lem:pointUniformizerData, pa-thm:ordZ_germGenericUnits,
    pa-lem:baseFieldTransition_genericPoint, pa-def:scheme_ordZ, pa-def:residueDeg}
  \emph{Reduction to a one-point divisor.} The class map
  \(\Div(C_{K_1}) \to \Pic(C_{K_1})\) is surjective (\ref{pa-def:curveDivisor_picClass}),
  so \(L = \struct{C_{K_1}}(D)\) for a Weil divisor \(D\), a finite \(\Z\)-combination of
  one-point divisors. Both sides of the claimed identity are additive in \(L\): the
  class map and the pullback \(\pi^{*}\) are homomorphisms
  (\ref{pa-thm:curveDivisor_picClass_hom}) and \(\deg_{\Pic}\) is additive with
  \(\deg_{\Pic}(L^n) = n\deg_{\Pic}(L)\) (\ref{pa-lem:classDeg_zpow}). It therefore
  suffices to prove the identity for \(L = \struct{C_{K_1}}(1 \cdot x')\), \(x'\) a
  non-generic point.

  \emph{The pulled point class as a divisor class.} The one-point class is the class
  of the tracked point equations at \(x'\)
  (\ref{pa-thm:curveDivisor_picClass_single}). Pulling back a local-equation system
  computes the pullback of its class (\ref{pa-thm:localEquations_picClass_pullback},
  with regularity discharged as in \ref{pa-def:pointTransitionEquations}), so
  \(\pi^{*}L\) is the Picard class of the pulled tracked point system; and the class
  of any local-equation system is the class of the divisor of its presentation
  (\ref{pa-lem:localEquations_presentation}, \ref{pa-def:curveDivisor_picClass}). Writing
  \(D''\) for the presentation divisor of the pulled system
  (\ref{pa-def:presentationDivisor}), the normalization of \ref{pa-def:classDeg} turns both
  sides into divisor degrees:
  \[
    \deg_{\Pic}\bigl(\pi^{*}L\bigr) = \deg_{K_2}\bigl(D''\bigr), \qquad
    \deg_{\Pic}(L) = \deg_{K_1}\bigl(1 \cdot x'\bigr) = \bigl[\kappa(x') : K_1\bigr]
  \]
  (\ref{pa-lem:curveDivisor_single_arith}).

  \emph{The degree of the pulled point divisor.} The coefficient of \(D''\) at a point
  \(z\) is the order at \(z\) of the trivialising element \(f_z\)
  (\ref{pa-def:presentationDivisor}). Off the fibre of \(x'\) the element is \(1\)
  (\ref{pa-lem:pointTransitionEquations_elem}), of order zero; so \(D''\) is supported in
  the finite set \(S := \mathrm{supp}(D'') \subseteq \pi^{-1}(x')\), and
  \[
    \deg_{K_2}(D'') \;=\; \sum_{z \in S}
      \mathrm{ord}_z\bigl(\pi^{\sharp}t\bigr) \cdot \bigl[\kappa(z) : K_2\bigr],
  \]
  the coefficient on the fibre being the order of \(\pi^{\sharp}(t)\)
  (\ref{pa-lem:pointTransitionEquations_elem}).

  \emph{Closing with the fiber-degree identity.} Choose an affine chart
  \(V \ni x'\) inside the tracked neighbourhood of the uniformizer data (affine opens
  form a basis), containing the generic point, and restrict the tracked section
  \(s\) to it; the restriction still has germ \(t\) at \(\eta\) and is still a unit at
  every point of \(V\) other than \(x'\) (\ref{pa-lem:pointUniformizerData}). The
  hypotheses of the fiber-degree identity
  (\ref{pa-thm:sum_ordZ_residueDeg_baseFieldTransition}) hold with this chart, the unit
  \(g = t\), and the finite set \(S\): its members lie in
  \(\pi^{-1}V\) (they lie in the fibre of \(x' \in V\)), and at a point \(z\) of
  \(\pi^{-1}V\) outside \(S\) the pulled section is a unit --- off the fibre because
  \(s\) is a unit at \(\pi(z) \neq x'\) and unit germs pull back to unit germs; on the
  fibre because a non-unit germ at \(z\) would give
  \(\mathrm{ord}_z(\pi^{\sharp}t) \neq 0\) (a rational function that is the generic
  germ of a section non-invertible at \(z\) has nontrivial order there --- the
  contrapositive of \ref{pa-thm:ordZ_germGenericUnits}, through the discrete valuation of
  the stalk at \(z\)), putting \(z\) in the support \(S\). The identity gives
  \[
    \sum_{z \in S} \mathrm{ord}_z\bigl(\pi^{\sharp}t\bigr) \cdot
      \bigl[\kappa(z) : K_2\bigr]
    \;=\; \mathrm{ord}_{x'}(t) \cdot \bigl[\kappa(x') : K_1\bigr]
    \;=\; \bigl[\kappa(x') : K_1\bigr],
  \]
  the uniformizer having order one (\ref{pa-lem:uniformizer}). Chaining the displays
  proves the one-point case, and with it the theorem.
\end{proof}

\begin{lemma}[The degree vanishes on classes from the base field]
  \label{pa-lem:classDeg_picFromBase}
  \lean{AlgebraicGeometry.classDeg_eq_zero_of_mem_picFromBase}
  \leanok
  \dcref{custom}

  \uses{pa-def:classDeg}
  Let \(K/k\) be a field extension. Every \v Cech Picard class on \(C_K\) pulled back
  from \(\Spec K\) along the second projection is trivial; in particular its degree is
  zero.
\end{lemma}
\begin{proof}
  \leanok
  The underlying space of \(\Spec K\) is a single point, and the \v Cech Picard group
  of a scheme with one-point underlying space is trivial: every member of a pointed
  cover is an open containing the point, hence the whole space, so a unit cocycle
  \((g_{xy})\) lives on the cover all of whose members and overlaps are the total
  space; fixing a base point \(x_0\) and setting \(b_x := g_{x x_0}\), the cocycle
  identity gives \(g_{xy} = g_{x x_0}\, g_{x_0 y} = b_x\, b_y^{-1}\), a coboundary. So
  the pulled-back class is the pullback of the trivial class, which is trivial, of
  degree \(\deg_{\Pic}(1) = 0\) (\ref{pa-def:classDeg}).
\end{proof}

\begin{definition}[The degree on the relative Picard group over a field]
  \label{pa-def:relPicDeg}
  \lean{AlgebraicGeometry.relPicDeg, AlgebraicGeometry.relPicDeg_relPicMk}
  \leanok
  \dcref{custom}

  \uses{pa-lem:classDeg_picFromBase, pa-def:classDeg}
  For a field extension \(K/k\), recall the relative Picard group
  \(\mathrm{relPic}(C, \Spec K) = \Pic(C_K) \big/ (\text{classes from } \Spec K)\)
  (the quotient recalled at the head of Chapter~\ref{pa-chap:PicardEtale}). The degree
  homomorphism kills the classes from the base
  (\ref{pa-lem:classDeg_picFromBase}), so it descends along the quotient projection to
  the \emph{relative degree}
  \[
    \deg_{\mathrm{rel}} : \mathrm{relPic}(C, \Spec K) \longrightarrow \Z,
    \qquad \deg_{\mathrm{rel}}\bigl([L]\bigr) = \deg_{\Pic}(L),
  \]
  a homomorphism of groups; the displayed anchor holds for every
  \(L \in \Pic(C_K)\).
\end{definition}
\begin{proof}
  \leanok
  A group homomorphism trivial on a subgroup factors uniquely through the quotient by
  it; the anchor is the factorization identity.
\end{proof}

\begin{theorem}[E-iv descends to the relative Picard group]
  \label{pa-thm:relPicDeg_relPicAlgMap}
  \lean{AlgebraicGeometry.relPicDeg_relPicAlgMap}
  \leanok
  \dcref{custom}

  \uses{pa-def:relPicDeg, pa-def:relPicAlgMap, pa-thm:classDeg_cechPicMap_baseFieldTransition}
  For every \(k\)-algebra homomorphism \(\varphi : K_1 \to K_2\) of field extensions
  of \(k\) and every relative Picard class \(x \in \mathrm{relPic}(C, \Spec K_1)\),
  \[
    \deg_{\mathrm{rel}}\bigl(\varphi^{*}x\bigr) \;=\; \deg_{\mathrm{rel}}(x),
  \]
  where \(\varphi^{*}\) is restriction of relative Picard classes along
  \(\Spec\varphi\) (\ref{pa-def:relPicAlgMap}): the degree of a relative Picard class
  does not depend on the field over which it is presented. This is the
  well-definedness core of the degree of an \'etale-locally-trivial class: a class
  defined over \(K_1\) and read over the larger field \(K_2\) has the same degree ---
  its points split, but order times residue degree rebalances exactly
  (\ref{pa-thm:sum_ordZ_residueDeg_baseFieldTransition}).
\end{theorem}
\begin{proof}
  \leanok
  Represent \(x = [L]\) with \(L \in \Pic(C_{K_1})\). Restriction along
  \(\Spec\varphi\) is induced on the quotients by pullback along the whiskered
  transition (\ref{pa-def:relPicAlgMap}), so \(\varphi^{*}[L] = [\pi^{*}L]\). By the
  anchor of \ref{pa-def:relPicDeg} and the invariance
  (\ref{pa-thm:classDeg_cechPicMap_baseFieldTransition}),
  \[
    \deg_{\mathrm{rel}}\bigl(\varphi^{*}[L]\bigr)
      = \deg_{\Pic}\bigl(\pi^{*}L\bigr)
      = \deg_{\Pic}(L)
      = \deg_{\mathrm{rel}}\bigl([L]\bigr). \qedhere
  \]
\end{proof}

\section{Invariance of the degree under isomorphisms}
\label{pa-sec:degree_iso_invariance}

The base-field invariance has a lighter sibling over a \emph{single} field: the degree
of a \v Cech Picard class is invariant under pullback along an isomorphism of curve
bundles over one field \(K\). The proof is the ladder of
Section~\ref{pa-sec:degree_base_field_invariance} with the transition replaced by the
isomorphism, and every leg degenerates: the fibre of a closed point is a single point,
the coefficient of the pulled point divisor there is \(+1\) --- the maximal-adic
multiplicity transports along the stalk isomorphism, with no colength computation ---
and the residue degree transports along the induced isomorphism of residue fields.

Throughout this section \(K\) is a field and \(X\) is a curve bundle over \(K\) ---
integral, smooth of relative dimension one and quasi-compact over \(\Spec K\) --- with
generic point \(\eta_X\); the hypotheses on the second scheme \(W\) are stated per
node.

\begin{lemma}[Naturality of the algebra map on sections]
  \label{pa-lem:overAlgebraMap_app}
  \lean{AlgebraicGeometry.overAlgebraMap_app}
  \leanok
  \dcref{custom}

  \uses{pa-def:overAlgebraMap}
  Let \(W, X\) be schemes over \(\Spec K\) and \(f : W \to X\) a morphism of schemes
  over \(\Spec K\). For every open \(U \subseteq X\) and every scalar \(a \in K\), the
  pullback of sections along \(f\) carries the canonical image of \(a\) in
  \(\Gamma(X, U)\) (\ref{pa-def:overAlgebraMap}) to the canonical image of \(a\) in
  \(\Gamma(W, f^{-1}U)\):
  \[
    f^{\sharp}\bigl(\iota^{X}_{U}(a)\bigr) \;=\; \iota^{W}_{f^{-1}U}(a),
  \]
  where \(\iota\) denotes the algebra map on sections.
\end{lemma}
\begin{proof}
  \leanok
  The algebra map into \(\Gamma(X, U)\) is the composite of the map induced on global
  sections by the structure morphism of \(X\) with restriction to \(U\)
  (\ref{pa-def:overAlgebraMap}). Pullback along \(f\) commutes with restriction
  (naturality of \(f^{\sharp}\)), so \(f^{\sharp}(\iota^{X}_{U}(a))\) is the
  restriction to \(f^{-1}U\) of \(f^{\sharp}(\iota^{X}_{X}(a))\). And \(f\) lies over
  \(\Spec K\): the structure morphism of \(X\) composed with \(f\) is the structure
  morphism of \(W\), so on global sections \(f^{\sharp} \circ \iota^{X}_{X} =
  \iota^{W}_{W}\). Restricting to \(f^{-1}U\) gives the claim.
\end{proof}

\begin{lemma}[Residue degrees transport along an isomorphism]
  \label{pa-lem:residueDeg_eq_of_isIso}
  \lean{AlgebraicGeometry.residueDeg_eq_of_isIso}
  \leanok
  \dcref{custom}

  \uses{pa-def:residueDeg, pa-def:residueOverAlgebraMap}
  Let \(W, X\) be schemes over \(\Spec K\) and \(f : W \to X\) an isomorphism of
  schemes over \(\Spec K\). Then for every point \(z\) of \(W\),
  \[
    [\kappa(z) : K] \;=\; [\kappa(f(z)) : K].
  \]
\end{lemma}
\begin{proof}
  \leanok
  \uses{pa-lem:overAlgebraMap_app, pa-def:overAlgebraMap}
  The stalk map of \(f\) at \(z\) is an isomorphism of local rings
  \(\struct{X, f(z)} \to \struct{W, z}\), and a ring isomorphism of local rings carries
  the maximal ideal onto the maximal ideal (units correspond to units in both
  directions, and the maximal ideal of a local ring is the set of non-units); it
  therefore induces a ring isomorphism of residue fields
  \(r : \kappa(f(z)) \to \kappa(z)\).

  The isomorphism \(r\) intertwines the two \(K\)-structures
  (\ref{pa-def:residueOverAlgebraMap}): the \(K\)-structure of a residue field is the
  composite of the algebra map on global sections, the germ, and the residue map. By
  naturality of the algebra map on sections (\ref{pa-lem:overAlgebraMap_app}, at the full
  open) followed by the defining compatibility of the stalk map with germs, the stalk
  map of \(f\) carries the germ at \(f(z)\) of the canonical image of \(a \in K\) to
  the germ at \(z\) of its canonical image; passing to the residues, \(r\) carries the
  canonical image of \(a\) in \(\kappa(f(z))\) to its canonical image in
  \(\kappa(z)\).

  An additive bijection between two \(K\)-vector spaces which is multiplicative and
  matches the two canonical copies of \(K\) is \(K\)-linear: \(r(a \cdot m) =
  r(\bar a\, m) = r(\bar a)\, r(m) = \bar a\, r(m) = a \cdot r(m)\). So \(r\) is a
  \(K\)-linear equivalence, and the \(K\)-dimensions of \(\kappa(f(z))\) and
  \(\kappa(z)\) agree (\ref{pa-def:residueDeg}).
\end{proof}

\begin{definition}[The pulled tracked point system along a generic-to-generic morphism]
  \label{pa-def:pointPullbackEquations}
  \lean{AlgebraicGeometry.pointPullbackEquations, AlgebraicGeometry.pointPullbackUnit,
    AlgebraicGeometry.coe_pointPullbackUnit}
  \leanok
  \dcref{custom}

  \uses{pa-def:pointEquations, pa-lem:pointUniformizerData, pa-def:localEquations_pullback,
    pa-lem:pullbackEqn_regular, pa-def:functionFieldMap, pa-lem:uniformizer}
  Let \(x'\) be a non-generic point of \(X\), with tracked point equations at \(x'\)
  (\ref{pa-def:pointEquations}, from the tracked uniformizer data of
  \ref{pa-lem:pointUniformizerData}), and let \(W\) be an integral scheme over \(\Spec K\)
  and \(f : W \to X\) a morphism with \(f(\eta_W) = \eta_X\). The \emph{pulled tracked
  point system} is the pullback of this local-equation system along \(f\)
  (\ref{pa-def:localEquations_pullback}), with regularity discharged by integrality of
  both schemes (\ref{pa-lem:pullbackEqn_regular}). Its \emph{trivialising unit} is the
  pullback \(f^{\sharp}(t) \in K(W)^{\times}\) of the chosen uniformizer \(t\) at
  \(x'\) (\ref{pa-lem:uniformizer}, \ref{pa-def:functionFieldMap}). This is the construction
  of the pulled tracked point system of the base-field transition
  (\ref{pa-def:pointTransitionEquations}), for an arbitrary morphism mapping generic
  point to generic point.
\end{definition}

\begin{lemma}[The trivialising elements of the generic pulled point presentation]
  \label{pa-lem:pointPullbackEquations_elem}
  \lean{AlgebraicGeometry.pointPullbackEquations_presentation_elem_of_ne,
    AlgebraicGeometry.pointPullbackEquations_presentation_elem_of_eq}
  \leanok
  \dcref{custom}

  \uses{pa-def:pointPullbackEquations, pa-lem:localEquations_presentation,
    pa-def:pointEquations, pa-lem:functionFieldMap_germ}
  In the situation of \ref{pa-def:pointPullbackEquations}, the meromorphic presentation
  (\ref{pa-lem:localEquations_presentation}) of the pulled tracked point system has
  trivialising elements
  \[
    f_z \;=\;
    \begin{cases}
      f^{\sharp}(t), & f(z) = x', \\
      1, & f(z) \neq x',
    \end{cases}
  \]
  at each point \(z\) of \(W\).
\end{lemma}
\begin{proof}
  \leanok
  The trivialising element at \(z\) is the germ at \(\eta_W\) of the pulled equation
  indexed by \(z\), which is the \(f\)-pullback of the equation of the point system at
  \(f(z)\) (\ref{pa-def:pointPullbackEquations}). If \(f(z) \neq x'\), that equation is
  the constant \(1\) (\ref{pa-def:pointEquations}), and \(1\) pulls back to \(1\), with
  germ \(1\). If \(f(z) = x'\), the equation is the tracked section \(s\), whose germ
  at \(\eta_X\) is the chosen uniformizer \(t\) (\ref{pa-def:pointEquations}); by germ
  naturality (\ref{pa-lem:functionFieldMap_germ}) the germ at \(\eta_W\) of
  \(f^{\sharp}s\) is \(f^{\sharp}(t)\).
\end{proof}

\begin{lemma}[The pulled uniformizer has order one at the preimage point]
  \label{pa-lem:ordZ_pointPullbackUnit}
  \lean{AlgebraicGeometry.ordZ_pointPullbackUnit}
  \leanok
  \dcref{custom}

  \uses{pa-def:pointPullbackEquations, pa-def:scheme_ordZ, pa-lem:uniformizer,
    pa-lem:pointUniformizerData}
  Let \(W\) be a further curve bundle over \(K\) --- integral and smooth of relative
  dimension one over \(\Spec K\) --- and let \(f : W \to X\) be an isomorphism of
  schemes with \(f(\eta_W) = \eta_X\). Let \(x'\) be a non-generic point of \(X\) and
  \(z\) a non-generic point of \(W\) with \(f(z) = x'\). Then the trivialising unit of
  the pulled tracked point system has order one at \(z\):
  \[
    \mathrm{ord}_z\bigl(f^{\sharp}(t)\bigr) \;=\; 1.
  \]
\end{lemma}
\begin{proof}
  \leanok
  \uses{pa-lem:functionFieldMap_germ, pa-lem:germ_generic_eq_algebraMap_germ,
    pa-thm:closedpoint_dvr_stalk, pa-thm:closedpoint_dedekind_stalk, pa-def:scheme_ord,
    pa-lem:adic_valuation_count, pa-lem:mem_pow_iff_le_count}
  Let \(s\) be the tracked section of the uniformizer data at \(x'\), on its tracked
  neighbourhood \(V \ni x'\), with \(\mathrm{germ}_{\eta_X}(s) = t\)
  (\ref{pa-lem:pointUniformizerData}); write \(c \in \struct{X,x'}\) for the germ of
  \(s\) at \(x'\), a nonzerodivisor of the stalk by regularity of the tracked data,
  hence nonzero.

  \emph{Both orders are adic multiplicities of matched stalk elements.} The order
  valuation at a closed point is the maximal-adic valuation of the stalk, a discrete
  valuation ring and a Dedekind domain (\ref{pa-def:scheme_ord},
  \ref{pa-thm:closedpoint_dvr_stalk}, \ref{pa-thm:closedpoint_dedekind_stalk}); its value on
  a rational function which is the image of a nonzero germ under the localization of
  the stalk into the function field is the multiplicity of the maximal ideal in the
  prime factorization of the principal ideal generated by that germ
  (\ref{pa-lem:adic_valuation_count}, applied in the local Dedekind stalk). The
  uniformizer \(t\) is such an image: \(t\) is the image of \(c\) under
  \(\struct{X,x'} \to K(X)\) (\ref{pa-lem:germ_generic_eq_algebraMap_germ}), so
  \(\mathrm{ord}_{x'}(t)\) is the multiplicity of \(\mathfrak m_{x'}\) in the
  factorization of \((c)\). On the other side, by germ naturality
  (\ref{pa-lem:functionFieldMap_germ}) the pullback \(f^{\sharp}(t)\) is the germ at
  \(\eta_W\) of the pulled section \(f^{\sharp}(s)\) on \(f^{-1}V \ni z\), which in
  turn is the image under \(\struct{W,z} \to K(W)\) of the germ of \(f^{\sharp}(s)\)
  at \(z\) (\ref{pa-lem:germ_generic_eq_algebraMap_germ} again); and that germ is the
  image of \(c\) under the stalk map \(\struct{X,x'} \to \struct{W,z}\) of \(f\) at
  \(z\) --- the defining compatibility of the stalk map with germs. So
  \(\mathrm{ord}_z(f^{\sharp}t)\) is the multiplicity of \(\mathfrak m_z\) in the
  factorization of the principal ideal generated by the image of \(c\).

  \emph{The multiplicity is an isomorphism invariant.} Since \(f\) is an isomorphism,
  its stalk map at \(z\) is a ring isomorphism \(e : \struct{X,x'} \to \struct{W,z}\)
  of local Dedekind domains. It carries the maximal ideal onto the maximal ideal
  (units correspond to units in both directions), hence carries
  \(\mathfrak m_{x'}^{\,n}\) onto \(\mathfrak m_z^{\,n}\) for every \(n\), so
  \(c \in \mathfrak m_{x'}^{\,n}\) if and only if \(e(c) \in \mathfrak m_z^{\,n}\).
  Membership in the powers of the maximal ideal characterizes the multiplicity
  (\ref{pa-lem:mem_pow_iff_le_count}, applied in each stalk with the maximal ideal as the
  prime), so the two multiplicities agree.

  Chaining, \(\mathrm{ord}_z(f^{\sharp}t) = \mathrm{ord}_{x'}(t) = 1\), the
  uniformizer having order one (\ref{pa-lem:uniformizer}).
\end{proof}

\begin{theorem}[The degree is invariant under isomorphisms: scheme form]
  \label{pa-thm:classDeg_cechPicMap_of_isIso}
  \lean{AlgebraicGeometry.classDeg_cechPicMap_of_isIso}
  \leanok
  \dcref{custom}

  \uses{pa-def:classDeg}
  Let \(W\) and \(X\) both be curve bundles over the field \(K\) --- integral, smooth
  of relative dimension one and quasi-compact over \(\Spec K\) --- with \(H^0\) and
  \(H^1\) of their structure sheaves finite-dimensional over \(K\), and let
  \(f : W \to X\) be an isomorphism of schemes over \(\Spec K\). Then for every \v Cech
  Picard class \(L\) on \(X\),
  \[
    \deg_{\Pic}\bigl(f^{*}L\bigr) \;=\; \deg_{\Pic}(L).
  \]
\end{theorem}
\begin{proof}
  \leanok
  \uses{pa-lem:genericPoint_eq_of_surjective, pa-def:curveDivisor_picClass,
    pa-thm:curveDivisor_picClass_hom, pa-thm:classDeg_mul, pa-lem:classDeg_zpow,
    pa-thm:curveDivisor_picClass_single, pa-thm:localEquations_picClass_pullback,
    pa-def:pointPullbackEquations, pa-lem:pointPullbackEquations_elem,
    pa-lem:ordZ_pointPullbackUnit, pa-lem:residueDeg_eq_of_isIso,
    pa-lem:localEquations_presentation, pa-def:presentationDivisor, pa-def:curveDivisor_single,
    pa-lem:curveDivisor_single_arith, pa-def:scheme_ordZ, pa-def:residueDeg}
  This is the ladder of the base-field invariance
  (\ref{pa-thm:classDeg_cechPicMap_baseFieldTransition}) with the transition replaced by
  the isomorphism \(f\), every leg degenerate. As an isomorphism, \(f\) is a
  homeomorphism on underlying spaces --- in particular injective and surjective on
  points --- and both spaces are irreducible, so it matches the generic points:
  \(f(\eta_W) = \eta_X\) (\ref{pa-lem:genericPoint_eq_of_surjective}).

  \emph{Reduction to a one-point divisor.} The class map
  \(\Div(X) \to \Pic(X)\) is surjective (\ref{pa-def:curveDivisor_picClass}), so
  \(L = \struct{X}(D)\) for a Weil divisor \(D\), a finite \(\Z\)-combination of
  one-point divisors. Both sides of the claimed identity are additive in \(L\): the
  class map and the pullback \(f^{*}\) are homomorphisms
  (\ref{pa-thm:curveDivisor_picClass_hom}) and \(\deg_{\Pic}\) is additive with
  \(\deg_{\Pic}(L^n) = n \deg_{\Pic}(L)\) (\ref{pa-thm:classDeg_mul},
  \ref{pa-lem:classDeg_zpow}). It therefore suffices to prove the identity for
  \(L = \struct{X}(1 \cdot x')\), \(x'\) a non-generic point.

  \emph{The pulled point class as a divisor class.} The one-point class is the class
  of the tracked point equations at \(x'\)
  (\ref{pa-thm:curveDivisor_picClass_single}). Pulling back a local-equation system
  computes the pullback of its class (\ref{pa-thm:localEquations_picClass_pullback}, with
  regularity discharged as in \ref{pa-def:pointPullbackEquations}), so \(f^{*}L\) is the
  Picard class of the pulled tracked point system; and the class of any local-equation
  system is the class of the divisor of its presentation
  (\ref{pa-lem:localEquations_presentation}, \ref{pa-def:curveDivisor_picClass}). Writing
  \(D''\) for the presentation divisor of the pulled system
  (\ref{pa-def:presentationDivisor}), the normalization of \ref{pa-def:classDeg} turns both
  sides into divisor degrees:
  \[
    \deg_{\Pic}\bigl(f^{*}L\bigr) = \deg_{K}\bigl(D''\bigr), \qquad
    \deg_{\Pic}(L) = \deg_{K}\bigl(1 \cdot x'\bigr) = \bigl[\kappa(x') : K\bigr]
  \]
  (\ref{pa-lem:curveDivisor_single_arith}).

  \emph{The pulled point divisor is a single point.} The coefficient of \(D''\) at a
  non-generic point \(z\) of \(W\) is the order at \(z\) of the trivialising element
  \(f_z\) (\ref{pa-def:presentationDivisor}). Since \(f\) is injective on points, the
  only point of \(W\) mapping to \(x'\) is \(z_0 := f^{-1}(x')\), which is
  non-generic: its image \(x'\) differs from \(\eta_X = f(\eta_W)\). At \(z \neq
  z_0\), \(f(z) \neq x'\) and the trivialising element is \(1\)
  (\ref{pa-lem:pointPullbackEquations_elem}), of order zero; at \(z_0\) the element is
  the pulled uniformizer \(f^{\sharp}(t)\) (\ref{pa-lem:pointPullbackEquations_elem}), of
  order \(1\) (\ref{pa-lem:ordZ_pointPullbackUnit}). Hence \(D'' = 1 \cdot z_0\)
  (\ref{pa-def:curveDivisor_single}), of degree
  \(\deg_{K}(D'') = [\kappa(z_0) : K]\) (\ref{pa-lem:curveDivisor_single_arith}).

  \emph{Closing with the residue transport.} By
  \ref{pa-lem:residueDeg_eq_of_isIso},
  \([\kappa(z_0) : K] = [\kappa(f(z_0)) : K] = [\kappa(x') : K]\). Chaining the
  displays proves the one-point case, and with it the theorem.
\end{proof}

\begin{theorem}[The degree is invariant under isomorphisms of curve bundles]
  \label{pa-thm:classDeg_map_iso}
  \lean{AlgebraicGeometry.classDeg_map_iso}
  \leanok
  \dcref{custom}

  \uses{pa-def:classDeg, pa-thm:smooth_geometrically_reduced, pa-prop:curve_integral,
    pa-thm:h0_moduleKSheaf_curve, pa-thm:moduleFinite_hModule_one}
  Let \(K\) be a field and \(X, Y\) objects of \(\Over(\Spec K)\) which are proper,
  smooth of relative dimension one and geometrically irreducible over \(K\). Then for
  every isomorphism \(e : X \cong Y\) of \(\Over(\Spec K)\) and every \v Cech Picard
  class \(L\) on the underlying scheme of \(Y\),
  \[
    \deg_{\Pic}\bigl(e^{*}L\bigr) \;=\; \deg_{\Pic}(L),
  \]
  where \(e^{*}\) is pullback of \v Cech Picard classes along the underlying
  isomorphism of schemes.
\end{theorem}
\begin{proof}
  \leanok
  \uses{pa-thm:classDeg_cechPicMap_of_isIso, pa-thm:smooth_geometrically_reduced,
    pa-prop:curve_integral, pa-thm:h0_moduleKSheaf_curve, pa-thm:moduleFinite_hModule_one}
  The hypotheses of the scheme form are discharged on both bundles by the curve pack:
  smoothness gives geometric reducedness (\ref{pa-thm:smooth_geometrically_reduced}),
  which together with geometric irreducibility gives integrality
  (\ref{pa-prop:curve_integral}); properness includes quasi-compactness; and \(H^0\) and
  \(H^1\) of the structure sheaf are finite-dimensional over \(K\)
  (\ref{pa-thm:h0_moduleKSheaf_curve}, \ref{pa-thm:moduleFinite_hModule_one}). The
  underlying morphism of \(e\) is an isomorphism of schemes over \(\Spec K\): the
  forgetful functor preserves isomorphisms, and every morphism of \(\Over(\Spec K)\)
  commutes with the structure morphisms. Apply
  \ref{pa-thm:classDeg_cechPicMap_of_isIso} to it and to \(L\).
\end{proof}

\section{The degree seam of the same-carrier comparison}
\label{pa-sec:degree_crossBase_seam}

The base-field shuffle of the relative Picard group at an affine test
(Section~\ref{pa-sec:relPicCrossBase}) reads one class on two presentations of one fiber:
at a \emph{field} test \(K\) in a scalar tower \(k \to L \to K\), the fiber of the
\(k\)-curve at \(K\) and the fiber of the base-changed \(L\)-curve at \(K\) are
identified by the same-carrier comparison (\ref{pa-def:crossBaseAffineIso}), an
isomorphism of schemes over \(\Spec K\). This section shows that the degree over
\(K\) reads the same integer on the two presentations --- the invariance under
isomorphisms of Section~\ref{pa-sec:degree_iso_invariance} fired at the comparison.

Throughout, \(C\) is the challenge curve over \(k\) --- proper, smooth of relative
dimension one and geometrically irreducible --- and \(k \to L \to K\) is a scalar
tower of fields, with \(\sigma : \Spec L \to \Spec k\) the induced morphism and
\(C_L := \sigma^{*}C\) the base-changed curve over \(L\) (\ref{pa-def:baseChange}). We
write \(\Spec K/\Spec F\) for the affine test of \(K\) over a field \(F\), and
\(\mathrm{relPic}_L(C_L, K)\) for the relative Picard group of \(C_L\) at the affine
test of \(K\) over \(L\), as in Section~\ref{pa-sec:relPicCrossBase}.

\begin{theorem}[Degree invariance along the same-carrier comparison]
  \label{pa-thm:classDeg_cechPicMap_crossBaseAffineIso}
  \lean{AlgebraicGeometry.classDeg_cechPicMap_crossBaseAffineIso}
  \leanok
  \dcref{custom}

  \uses{pa-def:classDeg, pa-thm:classDeg_cechPicMap_of_isIso, pa-def:crossBaseAffineIso,
    pa-lem:crossBaseAffineIso_snd, pa-thm:baseChangeCurve_bundle,
    pa-thm:baseChangeCurve_integral, pa-thm:baseChangeCurve_cohomFinite, pa-def:baseChange}
  Let \(c : \bigl(C_L \otimes (\Spec K/\Spec L)\bigr)_{\mathrm{left}} \to
  \bigl(C \otimes (\Spec K/\Spec k)\bigr)_{\mathrm{left}}\) be the same-carrier
  comparison at the field test \(K\) (\ref{pa-def:crossBaseAffineIso}). Then for every
  \v Cech Picard class \(\Lambda\) on the fiber
  \(C_K = \bigl(C \otimes (\Spec K/\Spec k)\bigr)_{\mathrm{left}}\),
  \[
    \deg_{\Pic}\bigl(c^{*}\Lambda\bigr) \;=\; \deg_{\Pic}(\Lambda),
  \]
  both degrees taken over \(K\).
\end{theorem}
\begin{proof}
  \leanok
  The two fibers satisfy every hypothesis of the scheme-form invariance engine. For
  the \(k\)-side fiber \(C_K\): it is smooth of relative dimension one and
  quasi-compact over \(K\) through the second projection, and integral, with \(H^0\)
  and \(H^1\) of its structure sheaf finite-dimensional over \(K\)
  (\ref{pa-thm:baseChangeCurve_bundle}, \ref{pa-thm:baseChangeCurve_integral},
  \ref{pa-thm:baseChangeCurve_cohomFinite}, at the extension \(K/k\)). For the
  \(L\)-side fiber: the base-changed curve \(C_L\) is itself proper, smooth of
  relative dimension one and geometrically irreducible over \(L\) --- its underlying
  data is the fiber of \(C\) at \(L\) with the second projection
  (\ref{pa-def:baseChange}, \ref{pa-thm:baseChangeCurve_bundle} at the extension
  \(L/k\)) --- so the same three theorems, applied over the base field \(L\) to the
  curve \(C_L\) at the extension \(K/L\), give the identical package for
  \(\bigl(C_L \otimes (\Spec K/\Spec L)\bigr)_{\mathrm{left}}\) over \(K\).

  The comparison \(c\) is an isomorphism of schemes commuting with the two second
  projections (\ref{pa-lem:crossBaseAffineIso_snd}), which are the structure morphisms
  of the two fibers over \(\Spec K\); so \(c\) is an isomorphism of curve bundles
  over \(\Spec K\). Apply the scheme-form invariance
  (\ref{pa-thm:classDeg_cechPicMap_of_isIso}) to \(c\) and \(\Lambda\).
\end{proof}

\begin{theorem}[The relative degree is invariant under the base-field shuffle]
  \label{pa-thm:relPicDeg_relPicCrossBase}
  \lean{AlgebraicGeometry.relPicDeg_relPicCrossBase}
  \leanok
  \dcref{custom}

  \uses{pa-def:relPicDeg, pa-def:relPicCrossBase,
    pa-thm:classDeg_cechPicMap_crossBaseAffineIso}
  Let \(\mathrm{sh}_K : \mathrm{relPic}(C, \Spec K) \to \mathrm{relPic}_L(C_L, K)\)
  be the base-field shuffle at the field test \(K\) (\ref{pa-def:relPicCrossBase}), and
  let the relative degree on the target be that of Definition~\ref{pa-def:relPicDeg}
  applied over the base field \(L\) to the curve \(C_L\) at the extension \(K/L\)
  (its hypotheses discharged as in
  \ref{pa-thm:classDeg_cechPicMap_crossBaseAffineIso}). Then for every relative Picard
  class \(x \in \mathrm{relPic}(C, \Spec K)\),
  \[
    \deg_{\mathrm{rel}}\bigl(\mathrm{sh}_K(x)\bigr) \;=\; \deg_{\mathrm{rel}}(x) :
  \]
  a relative Picard class of the \(k\)-curve at the field test \(K\) and its shuffled
  presentation on the base-changed curve read the same integer.
\end{theorem}
\begin{proof}
  \leanok
  Represent \(x = \pi(\Lambda)\) with \(\Lambda \in \Pic(C_K)\), the projection being
  surjective. The shuffle computes on cosets:
  \(\mathrm{sh}_K(x) = \pi(c^{*}\Lambda)\) (\ref{pa-def:relPicCrossBase}). The anchor of
  Definition~\ref{pa-def:relPicDeg}, on each side, reads the relative degree off the
  representative:
  \(\deg_{\mathrm{rel}}\bigl(\mathrm{sh}_K(x)\bigr) = \deg_{\Pic}(c^{*}\Lambda)\) and
  \(\deg_{\mathrm{rel}}(x) = \deg_{\Pic}(\Lambda)\). These agree by
  \ref{pa-thm:classDeg_cechPicMap_crossBaseAffineIso}.
\end{proof}

\section{Multiplicativity of the degree under a finite morphism}
\label{pa-sec:degree_pullback}

The invariance ladders of the two preceding sections generalize to the classical
multiplicativity of the degree: pulling back a \v Cech Picard class along a finite
morphism \(f : X \to Y\) of curve bundles over one field \(K\) multiplies its degree by
\(n = [K(X) : K(Y)]\), the degree of the induced extension of function fields. This is
Hartshorne's \(\deg f^{*}D = \deg f \cdot \deg D\) for complete nonsingular curves, in
bundle form. The proof is again the ladder of
Section~\ref{pa-sec:degree_base_field_invariance}, with two new algebraic inputs on a
Dedekind chart pair: the \emph{generic-rank dictionary} identifying \(n\) with the rank
of the chart extension, and the \emph{fiber-degree identity}
\(\sum_z e_z f_z = n\) over each closed point, which replaces the colength bridge of the
base-field transition.

Throughout, \(X\) and \(Y\) carry the hypotheses stated per node; \(\eta_X, \eta_Y\)
denote the generic points.

\begin{definition}[The chart square of a dominant morphism]
  \label{pa-def:pullbackSectionsAlgebra}
  \lean{AlgebraicGeometry.Scheme.Hom.pullbackSectionsAlgebra,
    AlgebraicGeometry.Scheme.Hom.algebraMap_pullbackSectionsAlgebra,
    AlgebraicGeometry.Scheme.Hom.functionFieldAlgebra,
    AlgebraicGeometry.Scheme.Hom.algebraMap_functionFieldAlgebra,
    AlgebraicGeometry.Scheme.Hom.pullbackSectionsFunctionFieldAlgebra,
    AlgebraicGeometry.Scheme.Hom.isScalarTower_pullbackSections_functionField,
    AlgebraicGeometry.Scheme.Hom.isScalarTower_functionFieldAlgebra}
  \leanok
  \dcref{custom}

  \uses{pa-def:functionFieldMap}
  Let \(f : X \to Y\) be a morphism of schemes and \(V \subseteq Y\) an open. Pullback
  of sections along \(f\) makes \(\Gamma(X, f^{-1}V)\) an algebra over
  \(\Gamma(Y, V)\), the \emph{chart extension} of \(f\) at \(V\). If moreover \(X\) and
  \(Y\) are integral with \(f(\eta_X) = \eta_Y \in V\), the function-field map
  (\ref{pa-def:functionFieldMap}) makes \(K(X)\) an algebra over \(K(Y)\), and the chart
  extension followed by the germ at \(\eta_X\) makes \(K(X)\) an algebra over
  \(\Gamma(Y, V)\); these structures fit in a commuting square
  \[
    \begin{array}{ccc}
      \Gamma(Y, V) & \longrightarrow & \Gamma(X, f^{-1}V) \\
      \downarrow & & \downarrow \\
      K(Y) & \longrightarrow & K(X)
    \end{array}
  \]
  whose vertical maps are the germs at the generic points.
\end{definition}
\begin{proof}
  \leanok
  \uses{pa-lem:functionFieldMap_germ}
  Note first \(\eta_X \in f^{-1}V\), since \(f(\eta_X) = \eta_Y \in V\). The
  right-then-down composite of the square is the chart extension followed by the germ
  at \(\eta_X\); the down-then-right composite is the germ at \(\eta_Y\) followed by
  the function-field map. They agree by germ naturality
  (\ref{pa-lem:functionFieldMap_germ}): the germ at \(\eta_X\) of a pulled-back section is
  the function-field image of the germ of the section at \(\eta_Y\).
\end{proof}

\begin{lemma}[Module finiteness of the chart extension]
  \label{pa-lem:moduleFinite_pullbackSections}
  \lean{AlgebraicGeometry.Scheme.Hom.moduleFinite_pullbackSections}
  \leanok
  \dcref{custom}

  \uses{pa-def:pullbackSectionsAlgebra}
  Let \(f : X \to Y\) be a finite morphism of schemes and \(V \subseteq Y\) an affine
  open. Then \(\Gamma(X, f^{-1}V)\) is a finite module over \(\Gamma(Y, V)\).
\end{lemma}
\begin{proof}
  \leanok
  Finiteness of a morphism is affine-local on the target, so the chart extension
  \(\Gamma(Y, V) \to \Gamma(X, f^{-1}V)\) is an integral ring map and of finite type.
  An integral ring map of finite type is module-finite: each of the finitely many
  algebra generators is integral, so the image algebra is generated as a module by the
  finitely many monomials in the generators of degree below the degrees of their
  integral equations.
\end{proof}

\begin{lemma}[Injectivity of the chart extension]
  \label{pa-lem:faithfulSMul_pullbackSections}
  \lean{AlgebraicGeometry.Scheme.Hom.faithfulSMul_pullbackSections}
  \leanok
  \dcref{custom}

  \uses{pa-def:pullbackSectionsAlgebra}
  Let \(f : X \to Y\) be a morphism of integral schemes with \(f(\eta_X) = \eta_Y\),
  and let \(V \subseteq Y\) be an open containing \(\eta_Y\). Then the chart extension
  \(\Gamma(Y, V) \to \Gamma(X, f^{-1}V)\) is injective.
\end{lemma}
\begin{proof}
  \leanok
  \uses{pa-lem:functionFieldMap_germ}
  Let \(a, b \in \Gamma(Y, V)\) have equal pullbacks. Taking germs at \(\eta_X\) and
  applying germ naturality (\ref{pa-lem:functionFieldMap_germ}), the function-field map
  \(K(Y) \to K(X)\) takes the germs of \(a\) and \(b\) at \(\eta_Y\) to equal values;
  the function-field map is injective (\ref{pa-lem:functionFieldMap_germ}), so the germs
  of \(a\) and \(b\) at \(\eta_Y\) agree; and on an integral scheme the germ map at the
  generic point is injective on sections of a nonempty open, so \(a = b\).
\end{proof}

\begin{theorem}[The generic-rank dictionary]
  \label{pa-thm:finrank_functionField_eq_finrank_pullbackSections}
  \lean{AlgebraicGeometry.Scheme.Hom.finrank_functionField_eq_finrank_pullbackSections}
  \leanok
  \source{hartshorne-algebraic-geometry:page-0154, hartshorne-algebraic-geometry:page-0155}

  \uses{pa-def:pullbackSectionsAlgebra}
  \dcref{custom}
  Let \(f : X \to Y\) be a finite morphism of integral schemes with
  \(f(\eta_X) = \eta_Y\), and let \(V \subseteq Y\) be an affine open containing
  \(\eta_Y\). Then the degree of the function-field extension equals the rank of the
  chart extension:
  \[
    \bigl[K(X) : K(Y)\bigr]
      \;=\; \operatorname{rank}_{\Gamma(Y, V)} \Gamma(X, f^{-1}V),
  \]
  the rank of the finite torsion-free module \(\Gamma(X, f^{-1}V)\) over the domain
  \(\Gamma(Y, V)\) being the maximal number of linearly independent elements.
\end{theorem}
\begin{proof}
  \leanok
  \uses{pa-lem:moduleFinite_pullbackSections, pa-lem:faithfulSMul_pullbackSections}
  Write \(B = \Gamma(Y, V)\) and \(A = \Gamma(X, f^{-1}V)\), domains, since nonempty
  opens of integral schemes have integral section rings; \(A\) is a finite \(B\)-module
  (\ref{pa-lem:moduleFinite_pullbackSections}) and \(B \to A\) is injective
  (\ref{pa-lem:faithfulSMul_pullbackSections}). The germ maps at the generic points
  identify \(K(Y)\) with the fraction field of \(B\) and \(K(X)\) with the fraction
  field of \(A\) --- on an integral scheme, the function field is the fraction field of
  the sections on any nonempty affine open --- and under these identifications the
  square of \ref{pa-def:pullbackSectionsAlgebra} commutes, so the extension
  \(K(Y) \subseteq K(X)\) is obtained from \(B \subseteq A\) by passing to fractions.

  Let \(S = B \setminus \{0\}\), a multiplicative set of nonzerodivisors of \(A\) (its
  image avoids \(0\) by injectivity, and \(A\) is a domain). The localization
  \(S^{-1}A = K(Y) \otimes_B A\) is a finite-dimensional \(K(Y)\)-algebra, and it is a
  domain, being a subring of \(K(X)\); a finite-dimensional algebra over a field which
  is a domain is a field (multiplication by a nonzero element is injective, hence
  surjective, by dimension). A field lying between \(A\) and its fraction field
  contains the inverse of every nonzero element of \(A\), so \(S^{-1}A = K(X)\).

  Finally, rank passes to fractions: a family of elements of \(A\) is linearly
  independent over \(B\) if and only if it is linearly independent over \(K(Y)\) in
  \(S^{-1}A\) --- one direction clears denominators, the other restricts scalars ---
  so the rank of \(A\) over \(B\) is the dimension of \(S^{-1}A = K(X)\) over
  \(K(Y)\).
\end{proof}

\begin{theorem}[The fiber-degree identity for a finite morphism]
  \label{pa-thm:sum_ordZ_residueDeg_of_isFinite}
  \lean{AlgebraicGeometry.sum_ordZ_residueDeg_of_isFinite}
  \leanok
  \source{hartshorne-algebraic-geometry:page-0155}

  \uses{pa-def:scheme_ordZ, pa-def:residueDeg, pa-def:functionFieldMap,
    pa-def:pullbackSectionsAlgebra}
  \dcref{custom}
  Let \(K\) be a field, \(X\) and \(Y\) integral schemes smooth of relative dimension
  one over \(\Spec K\), and \(f : X \to Y\) a finite morphism over \(\Spec K\) with
  \(f(\eta_X) = \eta_Y\). Let \(V \ni \eta_Y\) be an affine chart of \(Y\), let
  \(t \in \Gamma(Y, V)\) be a section whose germ at \(\eta_Y\) is the value of a unit
  \(g \in K(Y)^{\times}\), let \(y \in V\) be a non-generic point with
  \(\mathrm{ord}_y(g) = 1\) at which the tracked section is normalized, \(t\) being a
  unit at every point of \(V\) other than \(y\); and let \(S\) be a finite set of
  non-generic points of \(X\), contained in \(f^{-1}V\), outside which the pulled
  section \(f^{\sharp}t\) is a unit on \(f^{-1}V\). Then, in \(\Z\),
  \[
    \sum_{z \in S} \mathrm{ord}_z\bigl(f^{\sharp}g\bigr)
      \cdot \bigl[\kappa(z) : K\bigr]
    \;=\; \bigl[K(X) : K(Y)\bigr] \cdot \bigl[\kappa(y) : K\bigr].
  \]
  (The hypotheses force \(S\) to exhaust the non-unit locus of \(f^{\sharp}t\), which
  is contained in the fiber \(f^{-1}(y)\); members of \(S\) outside that locus
  contribute order zero. The statement is parameterised by such a finite \(S\) so that
  the consumer may take for \(S\) the support of a divisor already in hand.)
\end{theorem}
\begin{proof}
  \leanok
  \uses{pa-thm:isDedekindDomain_section, pa-lem:moduleFinite_pullbackSections,
    pa-lem:faithfulSMul_pullbackSections,
    pa-thm:finrank_functionField_eq_finrank_pullbackSections,
    pa-thm:toAdd_ordZ_eq_count_factors, pa-thm:finrank_quotient_primeIdealOf,
    pa-lem:prime_factor_span, pa-lem:primeIdealOf_ne_bot, pa-lem:primeIdealOf_fromSpec_base,
    pa-lem:fromSpec_base_mem, pa-lem:fromSpec_base_ne_genericPoint, pa-lem:appLE_overAlgebraMap,
    pa-thm:residueDeg_pos, pa-lem:functionFieldMap_germ}
  Write \(B = \Gamma(Y, V)\) and \(A = \Gamma(X, f^{-1}V)\). The chart \(f^{-1}V\) is
  affine (finite morphisms are affine) and contains \(\eta_X\), so both \(B\) and \(A\)
  are Dedekind domains (\ref{pa-thm:isDedekindDomain_section}); \(A\) is a finite
  \(B\)-module (\ref{pa-lem:moduleFinite_pullbackSections}) along the injective chart
  extension (\ref{pa-lem:faithfulSMul_pullbackSections}), hence torsion-free --- \(A\) is
  a domain on which the nonzero elements of \(B\) act injectively. The chart extension
  commutes with the two \(K\)-algebra structures, since \(f\) lies over \(\Spec K\)
  (\ref{pa-lem:appLE_overAlgebraMap}). By germ naturality
  (\ref{pa-lem:functionFieldMap_germ}), the germ of \(f^{\sharp}t\) at \(\eta_X\) is the
  value of the unit \(f^{\sharp}g\). Under \(V \cong \Spec B\), a point \(z\) of the
  chart corresponds to a prime with the germ of a section at \(z\) becoming its image
  in the localization at that prime; so \emph{a section lies in the prime of \(z\) if
  and only if its germ at \(z\) is a non-unit}. We use this dictionary on both charts
  (\ref{pa-lem:fromSpec_base_mem}, \ref{pa-lem:primeIdealOf_fromSpec_base}).

  \emph{The tracked section cuts out the tracked prime.} The point \(y\) corresponds
  to a nonzero prime \(\mathfrak p\) of \(B\) (\ref{pa-lem:primeIdealOf_ne_bot}), maximal
  since nonzero primes of a Dedekind domain are maximal. In the Dedekind domain \(B\)
  the nonzero ideal \((t)\) factors into primes. Every prime factor \(\mathfrak P\) of
  \((t)\) contains \(t\) (\ref{pa-lem:prime_factor_span}), so at the point of \(V\)
  corresponding to \(\mathfrak P\) the germ of \(t\) is a non-unit; by the unit
  hypothesis that point is \(y\), so \(\mathfrak P = \mathfrak p\). The multiplicity
  of \(\mathfrak p\) in \((t)\) is \(\mathrm{ord}_y(g) = 1\)
  (\ref{pa-thm:toAdd_ordZ_eq_count_factors}), so \((t) = \mathfrak p\). Extending along
  \(B \to A\), which carries \(t\) to \(f^{\sharp}t\),
  \[
    \bigl(f^{\sharp}t\bigr) \;=\; \mathfrak p A \;\neq\; 0 .
  \]

  \emph{Each term is \(e \cdot f\) times \([\kappa(y):K]\).} Let \(z \in S\) contribute
  a nonzero term; then the germ of \(f^{\sharp}t\) at \(z\) is a non-unit --- a unit
  germ at \(z\) forces \(\mathrm{ord}_z(f^{\sharp}g) = 0\)
  (\ref{pa-thm:toAdd_ordZ_eq_count_factors}: the order is the multiplicity of the prime of
  \(z\) in \((f^{\sharp}t)\), and a unit generates an ideal not contained in that
  prime; concretely, \(f^{\sharp}t \notin \mathfrak q_z\)). So the nonzero prime
  \(\mathfrak q = \mathfrak q_z\) of \(A\) at \(z\) (\ref{pa-lem:primeIdealOf_ne_bot})
  contains \(f^{\sharp}t\), hence \(\mathfrak q \supseteq \mathfrak p A\), hence
  \(\mathfrak q \cap B \supseteq \mathfrak p\); maximality of \(\mathfrak p\) gives
  \(\mathfrak q \cap B = \mathfrak p\): the prime \(\mathfrak q\) \emph{lies over}
  \(\mathfrak p\). Write \(e_{\mathfrak q}\) for the multiplicity of \(\mathfrak q\) in
  the factorization of \(\mathfrak p A\) (the ramification index) and
  \(f_{\mathfrak q} = \dim_{B/\mathfrak p} A/\mathfrak q\) (the inertia degree). Then:
  \[
    \mathrm{ord}_z\bigl(f^{\sharp}g\bigr)
      = \text{multiplicity of } \mathfrak q \text{ in } \bigl(f^{\sharp}t\bigr)
      = e_{\mathfrak q}
    \quad (\ref{pa-thm:toAdd_ordZ_eq_count_factors}),
  \]
  and, using the residue leg \([\kappa(z) : K] = \dim_K A/\mathfrak q\) on both charts
  (\ref{pa-thm:finrank_quotient_primeIdealOf}) and the tower
  \(K \to B/\mathfrak p \to A/\mathfrak q\) --- whose \(K\)-structures match by
  \ref{pa-lem:appLE_overAlgebraMap} ---
  \[
    \bigl[\kappa(z) : K\bigr] = \dim_K A/\mathfrak q
      = f_{\mathfrak q} \cdot \dim_K B/\mathfrak p
      = f_{\mathfrak q} \cdot \bigl[\kappa(y) : K\bigr].
  \]

  \emph{Reindexing the sum over the primes over \(\mathfrak p\).} The assignment
  \(z \mapsto \mathfrak q_z\) is a bijection from the members of \(S\) with nonzero
  term onto the primes of \(A\) over \(\mathfrak p\): it is injective because the
  points--primes correspondence on the chart is
  (\ref{pa-lem:primeIdealOf_fromSpec_base}); and every prime \(\mathfrak q\) over
  \(\mathfrak p\) arises --- it is nonzero (its contraction is \(\mathfrak p \neq 0\)),
  so it corresponds to a non-generic point \(z\) of \(f^{-1}V\)
  (\ref{pa-lem:fromSpec_base_mem}, \ref{pa-lem:fromSpec_base_ne_genericPoint}) at which
  \(f^{\sharp}t \in \mathfrak p A \subseteq \mathfrak q\) is a non-unit germ, so
  \(z \in S\) by the exhaustion hypothesis, with term
  \(e_{\mathfrak q} \cdot f_{\mathfrak q} \cdot [\kappa(y):K] \neq 0\), the
  ramification index being at least \(1\) and the residue degree positive
  (\ref{pa-thm:residueDeg_pos}). Members of \(S\) with zero term contribute nothing.
  Hence
  \[
    \sum_{z \in S} \mathrm{ord}_z\bigl(f^{\sharp}g\bigr) \cdot
      \bigl[\kappa(z) : K\bigr]
    \;=\; \Bigl(\sum_{\mathfrak q \mid \mathfrak p}
      e_{\mathfrak q} f_{\mathfrak q}\Bigr) \cdot \bigl[\kappa(y) : K\bigr].
  \]

  \emph{The fundamental identity \(\sum e f = n\).} Localize at \(\mathfrak p\): over
  the discrete valuation ring \(B_{\mathfrak p}\) the module \(A_{\mathfrak p}\) is
  finite and torsion-free, hence free, of rank
  \(r = \operatorname{rank}_B A\) (rank is preserved by localization). Thus
  \(A_{\mathfrak p}/\mathfrak p A_{\mathfrak p}\) is an \(r\)-dimensional vector space
  over \(B/\mathfrak p\). It is an Artinian ring whose primes are the (finitely many)
  primes of \(A\) over \(\mathfrak p\), so it is the product of its localizations at
  them; the localization at \(\mathfrak q\) is
  \(A_{\mathfrak q}/\mathfrak p A_{\mathfrak q}\), and in the discrete valuation ring
  \(A_{\mathfrak q}\) we have \(\mathfrak p A_{\mathfrak q} =
  \mathfrak q^{e_{\mathfrak q}} A_{\mathfrak q}\), whose quotient filtration
  \(A_{\mathfrak q} \supseteq \mathfrak q A_{\mathfrak q} \supseteq \cdots \supseteq
  \mathfrak q^{e_{\mathfrak q}} A_{\mathfrak q}\) has \(e_{\mathfrak q}\) successive
  quotients, each one-dimensional over \(\kappa(\mathfrak q) = A/\mathfrak q\), hence
  of dimension \(f_{\mathfrak q}\) over \(B/\mathfrak p\). Summing dimensions,
  \(\sum_{\mathfrak q \mid \mathfrak p} e_{\mathfrak q} f_{\mathfrak q} = r\). By the
  generic-rank dictionary
  (\ref{pa-thm:finrank_functionField_eq_finrank_pullbackSections}),
  \(r = [K(X) : K(Y)]\). Chaining the displays gives the identity.
\end{proof}

\begin{theorem}[Degree multiplicativity under pullback: scheme form]
  \label{pa-thm:classDeg_cechPicMap_eq_finrank_mul}
  \lean{AlgebraicGeometry.classDeg_cechPicMap_eq_finrank_mul}
  \leanok
  \source{hartshorne-algebraic-geometry:page-0155}

  \uses{pa-def:classDeg, pa-def:pullbackSectionsAlgebra}
  \dcref{custom}
  Let \(X\) and \(Y\) be curve bundles over the field \(K\) --- integral, smooth of
  relative dimension one and quasi-compact over \(\Spec K\) --- with \(H^0\) and
  \(H^1\) of their structure sheaves finite-dimensional over \(K\), and let
  \(f : X \to Y\) be a finite morphism over \(\Spec K\) with \(f(\eta_X) = \eta_Y\).
  Then for every \v Cech Picard class \(L\) on \(Y\),
  \[
    \deg_{\Pic}\bigl(f^{*}L\bigr)
      \;=\; \bigl[K(X) : K(Y)\bigr] \cdot \deg_{\Pic}(L).
  \]
\end{theorem}
\begin{proof}
  \leanok
  \uses{pa-def:curveDivisor_picClass, pa-thm:curveDivisor_picClass_hom, pa-thm:classDeg_mul,
    pa-lem:classDeg_zpow, pa-thm:curveDivisor_picClass_single,
    pa-thm:localEquations_picClass_pullback, pa-def:pointPullbackEquations,
    pa-lem:pointPullbackEquations_elem, pa-lem:localEquations_presentation,
    pa-def:presentationDivisor, pa-def:curveDivisor_single, pa-lem:curveDivisor_single_arith,
    pa-thm:sum_ordZ_residueDeg_of_isFinite, pa-lem:uniformizer, pa-lem:pointUniformizerData,
    pa-thm:ordZ_germGenericUnits, pa-def:scheme_ordZ, pa-def:residueDeg}
  Write \(n = [K(X) : K(Y)]\).

  \emph{Reduction to a one-point divisor.} The class map
  \(\Div(Y) \to \Pic(Y)\) is surjective (\ref{pa-def:curveDivisor_picClass}), so
  \(L = \struct{Y}(D)\) for a Weil divisor \(D\), a finite \(\Z\)-combination of
  one-point divisors. Both sides of the claimed identity are additive in \(L\): the
  class map and the pullback \(f^{*}\) are homomorphisms
  (\ref{pa-thm:curveDivisor_picClass_hom}) and \(\deg_{\Pic}\) is additive with
  \(\deg_{\Pic}(L^m) = m \deg_{\Pic}(L)\) (\ref{pa-thm:classDeg_mul},
  \ref{pa-lem:classDeg_zpow}). It therefore suffices to prove the identity for
  \(L = \struct{Y}(1 \cdot x')\), \(x'\) a non-generic point of \(Y\).

  \emph{The pulled point class as a divisor class.} The one-point class is the class
  of the tracked point equations at \(x'\) (\ref{pa-thm:curveDivisor_picClass_single}).
  Pulling back a local-equation system computes the pullback of its class
  (\ref{pa-thm:localEquations_picClass_pullback}, with regularity discharged as in
  \ref{pa-def:pointPullbackEquations}), so \(f^{*}L\) is the Picard class of the pulled
  tracked point system along \(f\) (\ref{pa-def:pointPullbackEquations}, for the
  generic-to-generic morphism \(f\)); and the class of any local-equation system is
  the class of the divisor of its presentation
  (\ref{pa-lem:localEquations_presentation}, \ref{pa-def:curveDivisor_picClass}). Writing
  \(D''\) for the presentation divisor of the pulled system
  (\ref{pa-def:presentationDivisor}), the normalization of \ref{pa-def:classDeg} turns both
  sides into divisor degrees:
  \[
    \deg_{\Pic}\bigl(f^{*}L\bigr) = \deg_{K}\bigl(D''\bigr), \qquad
    \deg_{\Pic}(L) = \deg_{K}\bigl(1 \cdot x'\bigr) = \bigl[\kappa(x') : K\bigr]
  \]
  (\ref{pa-lem:curveDivisor_single_arith}).

  \emph{The degree of the pulled point divisor.} The coefficient of \(D''\) at a
  non-generic point \(z\) of \(X\) is the order at \(z\) of the trivialising element
  \(f_z\) (\ref{pa-def:presentationDivisor}). Off the fiber of \(x'\) the element is
  \(1\) (\ref{pa-lem:pointPullbackEquations_elem}), of order zero; so \(D''\) is
  supported in the finite set \(S := \mathrm{supp}(D'') \subseteq f^{-1}(x')\), and on
  the fiber the element is the pulled uniformizer
  (\ref{pa-lem:pointPullbackEquations_elem}), so
  \[
    \deg_{K}(D'') \;=\; \sum_{z \in S}
      \mathrm{ord}_z\bigl(f^{\sharp}t\bigr) \cdot \bigl[\kappa(z) : K\bigr].
  \]

  \emph{Closing with the fiber-degree identity.} Choose an affine chart
  \(V \ni x'\) inside the tracked neighbourhood of the uniformizer data (affine opens
  form a basis), containing the generic point (a nonempty open of an irreducible
  space), and restrict the tracked section \(s\) to it; the restriction still has germ
  \(t\) at \(\eta_Y\) and is still a unit at every point of \(V\) other than \(x'\)
  (\ref{pa-lem:pointUniformizerData}), and \(\mathrm{ord}_{x'}(t) = 1\)
  (\ref{pa-lem:uniformizer}). The hypotheses of the fiber-degree identity
  (\ref{pa-thm:sum_ordZ_residueDeg_of_isFinite}) hold with this chart, the unit
  \(g = t\), and the finite set \(S\): its members lie in \(f^{-1}V\) (they lie in the
  fiber of \(x' \in V\)), and at a non-generic point \(z\) of \(f^{-1}V\) outside
  \(S\) the pulled section is a unit --- off the fiber because \(s\) is a unit at
  \(f(z) \neq x'\) and unit germs pull back to unit germs; on the fiber because a
  non-unit germ at \(z\) would give \(\mathrm{ord}_z(f^{\sharp}t) \neq 0\) (a rational
  function which is the generic germ of a section non-invertible at \(z\) has
  nontrivial order there --- the contrapositive of \ref{pa-thm:ordZ_germGenericUnits},
  through the discrete valuation of the stalk at \(z\)), putting \(z\) in the support
  \(S\). The identity gives
  \[
    \sum_{z \in S} \mathrm{ord}_z\bigl(f^{\sharp}t\bigr) \cdot
      \bigl[\kappa(z) : K\bigr]
    \;=\; n \cdot \bigl[\kappa(x') : K\bigr].
  \]
  Chaining the displays proves the one-point case, and with it the theorem.
\end{proof}

\begin{theorem}[Degree multiplicativity under pullback of curve bundles]
  \label{pa-thm:classDeg_cechPicMap_of_isFinite}
  \lean{AlgebraicGeometry.classDeg_cechPicMap_of_isFinite}
  \leanok
  \source{hartshorne-algebraic-geometry:page-0155}

  \uses{pa-def:classDeg, pa-def:pullbackSectionsAlgebra}
  \dcref{custom}
  Let \(K\) be a field and \(D, E\) objects of \(\Over(\Spec K)\) which are proper,
  smooth of relative dimension one and geometrically irreducible over \(K\), and let
  \(h : D \to E\) be a morphism of \(\Over(\Spec K)\) whose underlying morphism is
  surjective and finite. Then for every \v Cech Picard class \(L\) on the underlying
  scheme of \(E\),
  \[
    \deg_{\Pic}\bigl(h^{*}L\bigr)
      \;=\; \bigl[K(D) : K(E)\bigr] \cdot \deg_{\Pic}(L).
  \]
\end{theorem}
\begin{proof}
  \leanok
  \uses{pa-thm:classDeg_cechPicMap_eq_finrank_mul, pa-thm:smooth_geometrically_reduced,
    pa-prop:curve_integral, pa-thm:h0_moduleKSheaf_curve, pa-thm:moduleFinite_hModule_one,
    pa-lem:genericPoint_eq_of_surjective}
  The hypotheses of the scheme form are discharged on both bundles by the curve pack:
  smoothness gives geometric reducedness (\ref{pa-thm:smooth_geometrically_reduced}),
  which together with geometric irreducibility gives integrality
  (\ref{pa-prop:curve_integral}); properness includes quasi-compactness; and \(H^0\) and
  \(H^1\) of the structure sheaves are finite-dimensional over \(K\)
  (\ref{pa-thm:h0_moduleKSheaf_curve}, \ref{pa-thm:moduleFinite_hModule_one}). The
  underlying morphism of \(h\) lies over \(\Spec K\) (every morphism of
  \(\Over(\Spec K)\) commutes with the structure morphisms) and, being surjective,
  matches the generic points (\ref{pa-lem:genericPoint_eq_of_surjective}). Apply
  \ref{pa-thm:classDeg_cechPicMap_eq_finrank_mul}.
\end{proof}

\begin{theorem}[The constant leg: pullbacks along one-point morphisms have degree zero]
  \label{pa-thm:classDeg_cechPicMap_eq_zero_of_range_eq_singleton}
  \lean{AlgebraicGeometry.classDeg_cechPicMap_eq_zero_of_range_eq_singleton}
  \leanok
  \dcref{custom}

  \uses{pa-def:classDeg}
  Let \(K\) be a field, \(D\) an object of \(\Over(\Spec K)\) which is proper, smooth
  of relative dimension one and geometrically irreducible over \(K\), and \(E\) an
  object of \(\Over(\Spec K)\) with proper structure morphism. Let \(h : D \to E\) be
  a morphism of \(\Over(\Spec K)\) whose set-range is a single closed point \(x\) of
  \(E\). Then for every \v Cech Picard class \(L\) on the underlying scheme of \(E\),
  \[
    \deg_{\Pic}\bigl(h^{*}L\bigr) \;=\; 0 .
  \]
\end{theorem}
\begin{proof}
  \leanok
  \uses{pa-cor:cechPicMap_eq_one_of_range_eq_singleton, pa-thm:overHom_isProper,
    pa-thm:smooth_geometrically_reduced, pa-prop:curve_integral, pa-thm:h0_moduleKSheaf_curve,
    pa-thm:moduleFinite_hModule_one}
  The bundle \(D\) is integral (\ref{pa-thm:smooth_geometrically_reduced},
  \ref{pa-prop:curve_integral}), in particular reduced, and its \(H^0\) and \(H^1\) are
  finite-dimensional over \(K\) (\ref{pa-thm:h0_moduleKSheaf_curve},
  \ref{pa-thm:moduleFinite_hModule_one}), so the degree homomorphism of
  \ref{pa-def:classDeg} is available on \(D\). The underlying morphism of \(h\) is proper
  (\ref{pa-thm:overHom_isProper}: the structure morphism of \(D\) is proper and that of
  \(E\) is separated, properness including separatedness), in particular
  quasi-compact. By \ref{pa-cor:cechPicMap_eq_one_of_range_eq_singleton},
  \(h^{*}L = 1\); and the degree of the trivial class is zero, \(\deg_{\Pic}\) being a
  homomorphism (\ref{pa-def:classDeg}).
\end{proof}

\section{The fibre twist of a map to the projective line}

Throughout this section \(Y\) is an integral scheme, smooth of relative dimension one and
quasi-compact over a field \(K\), equipped with a \emph{dominant} morphism \(\pi \colon Y \to
\PP^1\). Pulling the two standard charts of \(\PP^1\) back along \(\pi\) gives an affine two-cover
of \(Y\); the pullback of the chart coordinate cuts out, on this cover, the divisor of a fibre of
\(\pi\). The class of \(n\) copies of that fibre is the \emph{fibre twist} \(\Theta_n\), and the
purpose of the section is to construct it as a class trivialised on the two-cover by construction
and to record that its degree scales linearly in \(n\).

\begin{definition}[The pinned two-cover and its coordinate]
  \label{pa-def:fiberCover}
  \lean{AlgebraicGeometry.fiberCoord, AlgebraicGeometry.preimage_inf_eq_basicOpen_fiberCoord,
    AlgebraicGeometry.isUnit_fiberCoord_res_inf, AlgebraicGeometry.fiberCover}
  \leanok
  \dcref{custom}

  \uses{pa-def:P1_chartOpen, pa-def:P1_chartCoord, pa-lem:preimage_chart_sup, pa-lem:P1_chartOpen_inf}
  Write \(V_0 = \pi^{-1}(U_0)\) and \(V_1 = \pi^{-1}(U_1)\) for the preimages of the two standard
  charts \(U_i = D_+(X_i)\) of \(\PP^1\) (\ref{pa-def:P1_chartOpen}); they cover \(Y\)
  (\ref{pa-lem:preimage_chart_sup}). The pinned two-cover assigns to a point of \(Y\) the chart \(V_0\)
  if the point lies in \(V_0\), and \(V_1\) otherwise. Let
  \[
    t_0 \;=\; \pi^{\sharp}(t_{01}) \;\in\; \Gamma(Y, V_0)
  \]
  be the pullback along \(\pi\) of the chart-\(0\) coordinate \(t_{01} = X_1/X_0\)
  (\ref{pa-def:P1_chartCoord}). Then the overlap of the two-cover is exactly the locus where \(t_0\) is
  invertible:
  \[
    V_0 \cap V_1 \;=\; D(t_0),
  \]
  the basic open of \(t_0\); consequently the restriction of \(t_0\) to \(V_0 \cap V_1\) is a unit of
  \(\Gamma(Y, V_0 \cap V_1)\).
\end{definition}
\begin{proof}
  \leanok
  Since \(U_0 \cup U_1 = \PP^1\), the preimages cover \(Y\) (\ref{pa-lem:preimage_chart_sup}), so a
  point outside \(V_0\) lies in \(V_1\); the assignment is therefore a pointed cover. On \(\PP^1\)
  the coordinate \(t_{01} = X_1/X_0\) is a regular function on \(U_0\) that is invertible precisely
  where \(X_1\) also fails to vanish, that is, on \(U_0 \cap U_1 = D_+(X_0 X_1)\)
  (\ref{pa-lem:P1_chartOpen_inf}); in other words the basic open of \(t_{01}\) in \(U_0\) is
  \(U_0 \cap U_1\). Forming a basic open commutes with pullback along a morphism, so the locus where
  \(t_0 = \pi^{\sharp}(t_{01})\) is invertible is \(\pi^{-1}(U_0 \cap U_1) = V_0 \cap V_1\); that is,
  \(V_0 \cap V_1 = D(t_0)\). By definition of the basic open a section is invertible at every point
  of its basic open, so \(t_0\) restricts to a unit on \(V_0 \cap V_1 = D(t_0)\).
\end{proof}

\begin{definition}[The fibre divisor of \(n\cdot(\text{fibre})\)]
  \label{pa-def:fiberDivisor}
  \lean{AlgebraicGeometry.fiberDivisor, AlgebraicGeometry.germ_fiberCoord_mem_nonZeroDivisors,
    AlgebraicGeometry.genericPoint_mem_preimage_inf}
  \leanok
  \dcref{custom}

  \uses{pa-def:fiberCover, pa-def:localEquations}
  For \(n \in \N\), the \emph{fibre divisor of \(n\cdot(\text{fibre of }\pi)\)} is the
  local-equation system on the pinned two-cover (\ref{pa-def:localEquations}) whose equation is
  \(t_0^{\,n}\) on the chart \(V_0\) and the constant \(1\) on \(V_1\). It is a genuine system of
  local equations: each equation is regular, and on the overlap the two equations differ by the unit
  factor \(t_0^{\,n}\).
\end{definition}
\begin{proof}
  \leanok
  \uses{pa-lem:germ_generic_eq_algebraMap_germ}
  We check regularity and the overlap condition. Regularity on \(V_1\) is trivial, the equation
  being \(1\). On \(V_0\) it suffices to show that the germ of \(t_0\) at every point \(y \in V_0\)
  is a nonzerodivisor of the stalk \(\struct{Y,y}\), for then so is the germ of its power
  \(t_0^{\,n}\). Since \(Y\) is integral, each stalk embeds in the function field \(K(Y)\) through
  the localisation map, so a germ is a nonzerodivisor as soon as it is nonzero. The generic point
  \(\eta\) lies in the overlap \(V_0 \cap V_1\): as \(\pi\) is dominant its image is dense, and the
  \(\PP^1\)-overlap \(U_0 \cap U_1\) is a nonempty open (it is the range of an affine chart), so the
  preimage \(V_0 \cap V_1\) is a nonempty open of the irreducible space \(Y\) and hence contains
  \(\eta\). On \(V_0 \cap V_1 = D(t_0)\) the germ of \(t_0\) at \(\eta\) is a unit
  (\ref{pa-def:fiberCover}), in particular nonzero. For any \(y \in V_0\), the germ of \(t_0\) at
  \(\eta\) is the image of the germ at \(y\) under the localisation map \(\struct{Y,y} \to K(Y)\)
  (\ref{pa-lem:germ_generic_eq_algebraMap_germ}); were the germ at \(y\) zero, its image --- the
  nonzero unit germ at \(\eta\) --- would be zero, a contradiction. Hence the germ of \(t_0\) at each
  \(y \in V_0\) is nonzero, thus a nonzerodivisor. For the overlap condition, whenever a pair of
  cover members overlap the ratio of their equations is a unit: two \(V_0\)-members or two
  \(V_1\)-members give ratio \(1\), while a \(V_0\)--\(V_1\) overlap lies in \(V_0 \cap V_1\), where
  the ratio is \(t_0^{\,n}/1 = t_0^{\,n}\) (or its inverse), a unit since \(t_0\) is
  (\ref{pa-def:fiberCover}).
\end{proof}

\begin{definition}[The fibre twist and its power law]
  \label{pa-def:fiberTwist}
  \lean{AlgebraicGeometry.fiberTwist, AlgebraicGeometry.fiberTwist_succ,
    AlgebraicGeometry.fiberTwist_pow}
  \leanok
  \dcref{custom}

  \uses{pa-def:fiberDivisor, pa-def:localEquations_picClass}
  The \emph{fibre twist} \(\Theta_n \in \Pic(Y)\) is the Picard class of the fibre divisor of
  \(n\cdot(\text{fibre})\):
  \[
    \Theta_n \;=\; \struct Y(n\cdot\text{fibre}) \;=\; \mathrm{picClass}(\text{fibre divisor } n)
  \]
  (\ref{pa-def:localEquations_picClass}). It obeys the power law \(\Theta_{n+1} = \Theta_n \cdot
  \Theta_1\), with \(\Theta_0 = 1\); hence \(\Theta_n = \Theta_1^{\,n}\).
\end{definition}
\begin{proof}
  \leanok
  \uses{pa-lem:localEquations_ratioUnit}
  By construction \(\Theta_n\) is the class of the unit \v Cech cocycle of transition units of the
  fibre divisor (\ref{pa-def:localEquations_picClass}). Write \(g^{(n)}_{xy}\) for its transition unit
  between the members at points \(x, y\), the unique unit with \(f^{(n)}_x|_{\cap} = g^{(n)}_{xy}
  \cdot f^{(n)}_y|_{\cap}\) on the overlap (\ref{pa-lem:localEquations_ratioUnit}), where \(f^{(n)}\)
  denotes the fibre equation of exponent \(n\). Pointwise \(f^{(n+1)} = f^{(n)} \cdot f^{(1)}\), the
  equation being \(t_0^{\,n+1} = t_0^{\,n}\,t_0\) on \(V_0\) and \(1 = 1\cdot 1\) on \(V_1\).
  Restricting to an overlap and using the defining equations,
  \[
    g^{(n+1)}_{xy}\, f^{(n+1)}_y|_{\cap}
      = f^{(n+1)}_x|_{\cap}
      = f^{(n)}_x|_{\cap}\, f^{(1)}_x|_{\cap}
      = g^{(n)}_{xy}\, g^{(1)}_{xy}\, f^{(n)}_y|_{\cap}\, f^{(1)}_y|_{\cap}
      = g^{(n)}_{xy}\, g^{(1)}_{xy}\, f^{(n+1)}_y|_{\cap};
  \]
  the restriction of \(f^{(n+1)}_y\) to the overlap is a nonzerodivisor (regularity), so cancelling
  it gives \(g^{(n+1)}_{xy} = g^{(n)}_{xy}\, g^{(1)}_{xy}\). Thus the cocycle of exponent \(n+1\) is
  the product of those of exponents \(n\) and \(1\); passing to classes, and using that the class map
  from unit cocycles to \(\Pic(Y)\) is multiplicative, \(\Theta_{n+1} = \Theta_n \cdot \Theta_1\).
  The exponent-\(0\) equation is the constant \(1\), with trivial cocycle, so \(\Theta_0 = 1\).
  Induction on \(n\) yields \(\Theta_n = \Theta_1^{\,n}\).
\end{proof}

\begin{theorem}[The degree of the fibre twist scales linearly]
  \label{pa-thm:classDeg_fiberTwist}
  \lean{AlgebraicGeometry.classDeg_fiberTwist}
  \leanok
  \dcref{custom}

  \uses{pa-def:fiberTwist, pa-def:classDeg, pa-thm:classDeg_mul}
  Assume in addition that \(Y\) is smooth of relative dimension one and quasi-compact over \(K\)
  with finite-dimensional \(H^0\) and \(H^1\), so that the degree homomorphism \(\deg_{\Pic}\) of
  \ref{pa-def:classDeg} is available. Then the degree of the fibre twist scales linearly in \(n\):
  \[
    \deg_{\Pic} \Theta_n \;=\; n \cdot \deg_{\Pic} \Theta_1 .
  \]
\end{theorem}
\begin{proof}
  \leanok
  By the power law \(\Theta_n = \Theta_1^{\,n}\) (\ref{pa-def:fiberTwist}) and additivity of the class
  degree (\ref{pa-thm:classDeg_mul}), we argue by induction on \(n\). The base case is
  \(\deg_{\Pic}\Theta_0 = \deg_{\Pic} 1 = 0\). For the step,
  \[
    \deg_{\Pic}\Theta_{n+1} = \deg_{\Pic}(\Theta_n \cdot \Theta_1)
      = \deg_{\Pic}\Theta_n + \deg_{\Pic}\Theta_1
      = n\cdot\deg_{\Pic}\Theta_1 + \deg_{\Pic}\Theta_1
      = (n+1)\cdot\deg_{\Pic}\Theta_1 . \qedhere
  \]
\end{proof}

\section{The fibre unit and the fibre Weil divisor}

We now build the counterpart of the fibre twist on the side of the function field \(K(Y)\) and of
Weil divisors. The standing hypotheses are as before, with \(Y\) in addition locally of finite
type over \(K\); the coordinate \(t_0\), read at the generic point, becomes a unit \(u \in
K(Y)^{\times}\), and the effective divisor of its zeros is the fibre Weil divisor \(F\). Its zeros
lie on the fibre of \(\pi\) over the point \([1 : 0]\), the locus \(V_0 \setminus V_1\); this fibre
is nonempty whenever \(\pi\) is surjective, in particular for a finite dominant \(\pi\).

\begin{definition}[The fibre unit]
  \label{pa-def:fiberCoordUnit}
  \lean{AlgebraicGeometry.fiberCoord₁, AlgebraicGeometry.fiberCoord_mul_fiberCoord₁_res,
    AlgebraicGeometry.fiberCoordUnit}
  \leanok
  \dcref{custom}

  \uses{pa-def:fiberCover, pa-def:fiberDivisor, pa-def:P1_chartCoord, pa-lem:P1_chartOpen_inf}
  Let \(t_1 = \pi^{\sharp}(t_{10}) \in \Gamma(Y, V_1)\) be the pullback of the chart-\(1\) coordinate
  \(t_{10} = X_0/X_1\) (\ref{pa-def:P1_chartCoord}), a regular section on \(V_1\). On the overlap the
  two coordinates are mutually inverse,
  \[
    t_0|_{\cap} \cdot t_1|_{\cap} \;=\; 1 \qquad \text{in } \Gamma(Y, V_0 \cap V_1).
  \]
  Consequently the germ \(u := \mathrm{germ}_{\eta}(t_0) \in K(Y)\) is a unit of the function field,
  the \emph{fibre unit}; its inverse is \(\mathrm{germ}_{\eta}(t_1)\).
\end{definition}
\begin{proof}
  \leanok
  On \(\PP^1\) the two chart coordinates satisfy \(t_{01}\,t_{10} = (X_1/X_0)(X_0/X_1) = 1\) on the
  overlap \(U_0 \cap U_1 = D_+(X_0 X_1)\) (\ref{pa-def:P1_chartCoord}, \ref{pa-lem:P1_chartOpen_inf}).
  Pulling this identity back along \(\pi\)'s section map over the overlap gives \(t_0|_{\cap} \cdot
  t_1|_{\cap} = 1\) in \(\Gamma(Y, V_0 \cap V_1)\). The generic point \(\eta\) lies in the overlap
  (\ref{pa-def:fiberDivisor}), so taking germs at \(\eta\) --- a ring homomorphism, compatible with
  restriction --- yields \(\mathrm{germ}_{\eta}(t_0)\cdot \mathrm{germ}_{\eta}(t_1) = 1\) in
  \(K(Y)\). Hence \(u = \mathrm{germ}_{\eta}(t_0)\) is invertible with inverse
  \(\mathrm{germ}_{\eta}(t_1)\), a unit of \(K(Y)\).
\end{proof}

\begin{theorem}[The order table of the fibre unit]
  \label{pa-thm:fiberCoordUnit_order}
  \lean{AlgebraicGeometry.fiberCoordUnit_coeffAt_divOf_nonneg_of_mem_chart₀,
    AlgebraicGeometry.fiberCoordUnit_coeffAt_divOf_nonpos_of_mem_chart₁,
    AlgebraicGeometry.one_le_fiberCoordUnit_coeffAt_divOf_of_notMem_chart₁}
  \leanok
  \dcref{custom}

  \uses{pa-def:fiberCoordUnit, pa-def:fiberCover, pa-def:divOf, pa-def:scheme_ord, pa-def:scheme_ordZ,
    pa-lem:germ_generic_eq_algebraMap_germ, pa-lem:ord_algebraMap_stalk_le_one}
  The principal divisor \(\mathrm{div}(u)\) of the fibre unit (\ref{pa-def:divOf}) has, at a closed
  point \(x\), the following orders:
  \begin{itemize}
    \item if \(x \in V_0\), then \(\mathrm{ord}_x(u) \ge 0\) (no pole);
    \item if \(x \in V_1\), then \(\mathrm{ord}_x(u) \le 0\) (no zero);
    \item if \(x \in V_0\) but \(x \notin V_1\), then \(\mathrm{ord}_x(u) \ge 1\) (a genuine zero).
  \end{itemize}
  In particular \(u\) is a unit of the stalk \(\struct{Y,x}\) exactly on the overlap, and its zeros
  lie on \(V_0 \setminus V_1\), the fibre of \(\pi\) over \([1 : 0]\).
\end{theorem}
\begin{proof}
  \leanok
  For \(x \in V_0\), the value \(u = \mathrm{germ}_{\eta}(t_0)\) is the image, under the localisation
  map \(\struct{Y,x} \to K(Y)\), of the germ of \(t_0\) at \(x\)
  (\ref{pa-lem:germ_generic_eq_algebraMap_germ}). The image in \(K(Y)\) of a stalk element has no pole
  at \(x\) (\ref{pa-lem:ord_algebraMap_stalk_le_one}), so \(\mathrm{ord}_x(u) \ge 0\).

  For \(x \in V_1\), apply the same to \(u^{-1} = \mathrm{germ}_{\eta}(t_1)\) and the regular section
  \(t_1\) on \(V_1\): \(\mathrm{ord}_x(u^{-1}) \ge 0\). As \(\mathrm{ord}_x\) is a homomorphism
  \(K(Y)^{\times} \to \Z\) (\ref{pa-def:scheme_ordZ}), \(\mathrm{ord}_x(u) = -\mathrm{ord}_x(u^{-1}) \le
  0\).

  For \(x \in V_0\) with \(x \notin V_1\), the germ of \(t_0\) at \(x\) is not a unit, since \(x\)
  lies outside \(D(t_0) = V_0 \cap V_1\) (\ref{pa-def:fiberCover}); it is nonzero, its image in \(K(Y)\)
  being the unit \(u\). The stalk \(\struct{Y,x}\) is a discrete valuation ring (\ref{pa-def:scheme_ord}),
  in which a nonzero non-unit has strictly positive order; hence \(\mathrm{ord}_x(u) \ge 1\).
\end{proof}

\begin{definition}[The fibre Weil divisor]
  \label{pa-def:fiberWeilDivisor}
  \lean{AlgebraicGeometry.fiberWeilDivisor, AlgebraicGeometry.fiberWeilDivisor_nonneg,
    AlgebraicGeometry.fiberWeilDivisor_coeffAt_of_mem_chart₁,
    AlgebraicGeometry.fiberWeilDivisor_coeffAt_of_mem_chart₀,
    AlgebraicGeometry.fiberWeilDivisor_deg_nonneg,
    AlgebraicGeometry.fiberWeilDivisor_ne_zero_of_notMem_chart₁,
    AlgebraicGeometry.fiberWeilDivisor_deg_pos_of_notMem_chart₁}
  \leanok
  \dcref{custom}

  \uses{pa-thm:fiberCoordUnit_order, pa-def:divOf, pa-def:curveDivisor, pa-def:curveDivisor_deg,
    pa-thm:residueDeg_pos}
  The \emph{fibre Weil divisor} \(F\) is the effective part of \(\mathrm{div}(u)\): the Weil divisor
  (\ref{pa-def:curveDivisor}) with coefficient \(F_x = \max(\mathrm{ord}_x(u), 0)\) at each closed point
  \(x\). It is effective (\(F \ge 0\)); it vanishes on \(V_1\) (\(F_x = 0\) for \(x \in V_1\)); it
  agrees with \(\mathrm{div}(u)\) on \(V_0\) (\(F_x = \mathrm{ord}_x(u)\) for \(x \in V_0\)); and its
  degree is nonnegative. If moreover some closed point \(x_0 \in V_0\) lies outside \(V_1\) --- a
  point of the fibre over \([1 : 0]\) --- then \(F\) is nonzero and has strictly positive degree.
\end{definition}
\begin{proof}
  \leanok
  Effectivity is immediate from \(F_x = \max(\mathrm{ord}_x(u), 0) \ge 0\). On \(V_1\) one has
  \(\mathrm{ord}_x(u) \le 0\) (\ref{pa-thm:fiberCoordUnit_order}), so \(F_x = \max(\mathrm{ord}_x(u), 0)
  = 0\); on \(V_0\) one has \(\mathrm{ord}_x(u) \ge 0\), so \(F_x = \mathrm{ord}_x(u)\). The degree
  \(\deg_K F = \sum_x F_x\,[\kappa(x):K]\) (\ref{pa-def:curveDivisor_deg}) is a finite sum of
  nonnegative terms, each \(F_x \ge 0\) and each residue degree nonnegative, hence \(\deg_K F \ge 0\).
  Given a witness \(x_0 \in V_0 \setminus V_1\): there \(F_{x_0} = \mathrm{ord}_{x_0}(u) \ge 1\)
  (\ref{pa-thm:fiberCoordUnit_order}), so \(F \ne 0\); and the degree, again a sum of nonnegative terms,
  is bounded below by the single term \(F_{x_0}\,[\kappa(x_0):K] \ge 1\cdot[\kappa(x_0):K] > 0\),
  using positivity of the residue degree (\ref{pa-thm:residueDeg_pos}). Hence \(\deg_K F > 0\).
\end{proof}

\section{The twist lattices on the two-cover}

Fix a Weil divisor \(D\) on \(Y\). The twisted divisors \(D + nF\) have divisor sheaves whose
sections over the two charts and their overlap are \(K\)-submodules of the function field \(K(Y)\).
We record how these section lattices vary with \(n\): fixed over \(V_1\) and over the overlap,
growing over \(V_0\); from this the two-cover denominator lattices form an increasing chain that
exhausts the fixed overlap module \(N = \struct Y(D)(V_0 \cap V_1)\), by clearing poles at the
finite fibre over \([1 : 0]\).

\begin{theorem}[The twist leaves the lattice fixed off \(V_0\)]
  \label{pa-thm:divisorSections_twist_fixed}
  \lean{AlgebraicGeometry.divisorSections_add_nsmul_fiberWeilDivisor_chart₁,
    AlgebraicGeometry.divisorSections_add_nsmul_fiberWeilDivisor_overlap}
  \leanok
  \dcref{custom}

  \uses{pa-def:fiberWeilDivisor, pa-def:divisorSections}
  For every Weil divisor \(D\) and every \(n \in \N\), the sections of \(\struct Y(D + nF)\)
  coincide with those of \(\struct Y(D)\) over \(V_1\) and over the overlap \(V_0 \cap V_1\):
  \[
    \struct Y(D + nF)(V_1) = \struct Y(D)(V_1), \qquad
    \struct Y(D + nF)(V_0 \cap V_1) = \struct Y(D)(V_0 \cap V_1).
  \]
\end{theorem}
\begin{proof}
  \leanok
  The fibre divisor \(F\) vanishes at every point of \(V_1\) (\ref{pa-def:fiberWeilDivisor}), so at each
  closed point \(x \in V_1\) the coefficient of the twist is \((D + nF)_x = D_x + n F_x = D_x\). The
  overlap \(V_0 \cap V_1\) is contained in \(V_1\), so the same equality holds there. The sections of
  \(\struct Y(A)\) over an open \(U\) are the rational functions whose pole order at each closed point
  of \(U\) is bounded by \(A\) (\ref{pa-def:divisorSections}); since the bounds imposed by \(D + nF\) and
  by \(D\) agree at every closed point of \(V_1\) (respectively of the overlap), the two section
  modules coincide.
\end{proof}

\begin{theorem}[The twist grows the lattice over \(V_0\)]
  \label{pa-thm:divisorSections_twist_chart0}
  \lean{AlgebraicGeometry.divisorSections_add_nsmul_fiberWeilDivisor_chart₀}
  \leanok
  \dcref{custom}

  \uses{pa-def:fiberWeilDivisor, pa-def:fiberCoordUnit, pa-def:divisorSections, pa-thm:mem_boundedSections_mul_iff,
    pa-def:divOf}
  For every Weil divisor \(D\) and every \(n \in \N\), multiplication by \(u^{-n}\) carries
  \(\struct Y(D)(V_0)\) isomorphically onto \(\struct Y(D + nF)(V_0)\); that is,
  \[
    \struct Y(D + nF)(V_0) \;=\; u^{-n}\cdot \struct Y(D)(V_0)
  \]
  as \(K\)-submodules of \(K(Y)\).
\end{theorem}
\begin{proof}
  \leanok
  On \(V_0\) the fibre divisor agrees with \(\mathrm{div}(u)\) (\ref{pa-def:fiberWeilDivisor}), so at
  each closed point \(x \in V_0\) one has \((D + nF)_x = D_x + n\,\mathrm{ord}_x(u)\). Since
  \(\mathrm{div}(u^{-n}) = -n\,\mathrm{div}(u)\) (the principal divisor is a homomorphism,
  \ref{pa-def:divOf}), this reads \((D + nF)_x = D_x - \mathrm{ord}_x(u^{-n}) = \bigl(D -
  \mathrm{div}(u^{-n})\bigr)_x\). Hence \(\struct Y(D + nF)(V_0) = \struct Y\bigl(D -
  \mathrm{div}(u^{-n})\bigr)(V_0)\). By the bound-shift (\ref{pa-thm:mem_boundedSections_mul_iff}) with
  \(g = u^{-n}\), a rational function \(h\) has poles bounded by \(D\) on \(V_0\) if and only if
  \(u^{-n} h\) has poles bounded by \(D - \mathrm{div}(u^{-n})\) on \(V_0\); equivalently,
  multiplication by \(u^{-n}\) is a \(K\)-linear bijection from \(\struct Y(D)(V_0)\) onto
  \(\struct Y\bigl(D - \mathrm{div}(u^{-n})\bigr)(V_0) = \struct Y(D + nF)(V_0)\). As \(F\) is
  effective, these submodules increase with \(n\).
\end{proof}

\begin{definition}[The denominator lattice and its exhaustion]
  \label{pa-def:fiberLattice}
  \lean{AlgebraicGeometry.fiberLattice, AlgebraicGeometry.fiberLattice_le_overlap,
    AlgebraicGeometry.fiberLattice_mono, AlgebraicGeometry.iSup_fiberLattice}
  \leanok
  \dcref{custom}

  \uses{pa-thm:divisorSections_twist_fixed, pa-thm:divisorSections_twist_chart0, pa-def:fiberWeilDivisor,
    pa-thm:fiberCoordUnit_order, pa-def:divisorSections, pa-def:divOf}
  The \emph{denominator lattice} \(A_n\) of \(\struct Y(D + nF)\) is the join, inside \(K(Y)\), of the
  section lattices over the two charts:
  \[
    A_n \;=\; \struct Y(D + nF)(V_0) \;+\; \struct Y(D + nF)(V_1)
    \qquad (\text{as } K\text{-submodules of } K(Y)).
  \]
  It lies in the fixed ambient \(N := \struct Y(D)(V_0 \cap V_1)\), the lattices form an increasing
  chain, and they exhaust \(N\):
  \[
    A_n \le N, \qquad A_n \le A_{n+1}, \qquad \bigsqcup_{n} A_n = N .
  \]
\end{definition}
\begin{proof}
  \leanok
  \emph{Containment in \(N\).} For an inclusion of nonempty opens \(W \subseteq W'\), the section
  module \(\struct Y(A)(W')\) is contained in \(\struct Y(A)(W)\), a bound over more points being more
  restrictive (\ref{pa-def:divisorSections}). Applying this to \(V_0 \cap V_1 \subseteq V_0\) and \(V_0
  \cap V_1 \subseteq V_1\), and using \(\struct Y(D + nF)(V_0 \cap V_1) = N\)
  (\ref{pa-thm:divisorSections_twist_fixed}), both \(\struct Y(D + nF)(V_0)\) and \(\struct Y(D +
  nF)(V_1)\) are contained in \(N\); hence so is their join \(A_n\).

  \emph{The chain is increasing.} \(F\) is effective (\ref{pa-def:fiberWeilDivisor}), so \(D + nF \le D +
  (n+1)F\); the section lattice of a larger divisor contains that of a smaller one over any fixed open
  (\ref{pa-def:divisorSections}), over both \(V_0\) and \(V_1\), whence \(A_n \le A_{n+1}\).

  \emph{Exhaustion.} The inclusion \(\bigsqcup_n A_n \le N\) is the containment just proved. For the
  reverse, let \(s \in N\); we find \(n\) with \(s \in \struct Y(D + nF)(V_0) \subseteq A_n\). If
  \(s = 0\) it lies in every \(A_n\). Otherwise \(s\) is a unit of \(K(Y)\), and both \(D\) and
  \(\mathrm{div}(s)\) are supported on a finite set of closed points (\ref{pa-def:curveDivisor},
  \ref{pa-def:divOf}). Consider a closed point \(x \in V_0\). If also \(x \in V_1\), then \(F_x = 0\)
  (\ref{pa-def:fiberWeilDivisor}) and membership \(s \in N\) gives \(-\mathrm{ord}_x(s) \le D_x = (D +
  nF)_x\). If \(x \notin V_1\), then \(F_x \ge 1\) (\ref{pa-thm:fiberCoordUnit_order}, via
  \ref{pa-def:fiberWeilDivisor}); choosing \(n\) at least the maximum, over the finitely many exceptional
  points, of the pole deficit \(-\mathrm{ord}_x(s) - D_x\) makes \(n F_x \ge -\mathrm{ord}_x(s) -
  D_x\), that is \(-\mathrm{ord}_x(s) \le D_x + n F_x = (D + nF)_x\); at the remaining points both
  \(D_x\) and \(\mathrm{ord}_x(s)\) vanish and the bound is automatic. Thus for this \(n\), \(s\) has
  poles bounded by \(D + nF\) at every closed point of \(V_0\), so \(s \in \struct Y(D + nF)(V_0)
  \subseteq A_n \subseteq \bigsqcup_m A_m\). Hence \(N \le \bigsqcup_n A_n\), and equality holds.
\end{proof}

\section{Quasi-coherence of the divisor sheaves}

The two-cover computation of \(H^1\) with general coefficients
(Theorem~\ref{pa-thm:twoCoverH1LinearEquiv}) runs on any sheaf of \(K\)-modules whose degree-one
cohomology vanishes over each affine chart of the cover. For the divisor sheaves \(\struct X(D)\)
this vanishing follows, exactly as for the structure sheaf, from a quasi-coherence packaging
(Definition~\ref{pa-def:qcohOn}). This section supplies that packaging on every open, so the twisted
affine vanishing fires on every \(\struct X(D)\).

\begin{theorem}[The divisor sheaves are quasi-coherently packaged]
  \label{pa-thm:divisorSheaf_qcohOn}
  \lean{AlgebraicGeometry.Scheme.divisorSheaf.instQcohOn}
  \leanok
  \dcref{custom}

  \uses{pa-def:qcohOn, pa-def:divisorSheaf, pa-def:divisorSections, pa-lem:ord_algebraMap_stalk_le_one,
    pa-lem:ord_eq_one_of_mem_basicOpen, pa-def:divOf, pa-def:curveDivisor}
  On the curve bundle \(X\), for every Weil divisor \(D\) the divisor sheaf \(\struct X(D)\),
  viewed as a sheaf of \(K\)-modules, carries a quasi-coherence packaging on every open \(U\): an
  ambient section \(r \in \Gamma(X, U)\) acts on \(\struct X(D)(W)\), \(W \le U\), by multiplication
  with its germ at the generic point \(\eta\).
\end{theorem}
\begin{proof}
  \leanok
  When \(\eta \notin U\) all section modules over \(W \le U\) are zero (the sheaf is guarded by the
  generic point) and every packaging axiom is trivial. Assume \(\eta \in U\).

  \emph{(P1) action.} For \(r \in \Gamma(X, U)\) the germ \(\mathrm{germ}_\eta(r)\) is integral at
  every closed point \(x \in U\) --- its order is \(\le 1\), no pole (Lemma~\ref{pa-lem:ord_algebraMap_stalk_le_one})
  --- so multiplying a section of \(\struct X(D)\) by it can only improve the pole bounds: the
  product again lies in \(\struct X(D)(W)\) (Definition~\ref{pa-def:divisorSections}). The module laws
  over \(\Gamma(X, U)\) and naturality under restriction hold because restriction of \(\struct
  X(D)\)-sections is the inclusion of submodules of the function field \(K(X)\), on which the action
  is ordinary multiplication.

  \emph{(P2a) denominator clearing.} Let \(g \in \Gamma(X, U)\) and \(c \in \struct X(D)(V \cap
  D(g))\) with \(V \cap D(g)\) nonempty. On \(D(g)\) the germ \(\mathrm{germ}_\eta(g)\) has order
  \(0\) (Lemma~\ref{pa-lem:ord_eq_one_of_mem_basicOpen}), while off \(D(g)\), at a zero of \(g\), it
  has order \(\ge 1\); both \(D\) and \(\mathrm{div}(c)\) are finitely supported
  (Definitions~\ref{pa-def:curveDivisor}, \ref{pa-def:divOf}). Choosing \(N\) at least the finite
  pole-deficit budget over the exceptional points, the section \(\mathrm{germ}_\eta(g)^N \cdot c\)
  has poles bounded by \(D\) on all of \(V\): at points of \(D(g)\) the factor is a unit and the
  bound is inherited from the overlap, at zeros of \(g\) the \(N\)-fold vanishing clears the pole,
  elsewhere all orders vanish. Thus \(c = (g^N \cdot e)|_{V \cap D(g)}\) with \(e =
  \mathrm{germ}_\eta(g)^N \cdot c \in \struct X(D)(V)\).

  \emph{(P2b) defect annihilation.} Restrictions between nonempty opens are injective on the
  underlying rational functions, so a section killed on \(V \cap D(g)\) is already zero when \(V
  \cap D(g)\) is nonempty; and when \(V \cap D(g) = \emptyset\) while \(V\) is nonempty, the germ
  \(\mathrm{germ}_\eta(g)\) is a non-unit of the function \emph{field}, hence \(0\), so \(g\) itself
  annihilates. Either way a power of \(g\) kills the defect.
\end{proof}

\begin{theorem}[Twisted affine vanishing for divisor sheaves]
  \label{pa-thm:subsingleton_hModule_divisorSheaf_one}
  \lean{AlgebraicGeometry.IsAffineOpen.subsingleton_hModule'_divisorSheaf_one}
  \leanok
  \dcref{custom}

  \uses{pa-thm:divisorSheaf_qcohOn, pa-thm:hModule1_vanishing_qcoh}
  For every affine open \(U\) of the curve bundle \(X\) and every Weil divisor \(D\), the
  degree-one cohomology of \(U\) with coefficients in \(\struct X(D)\) vanishes: \(H'^1(U, \struct
  X(D)) = 0\).
\end{theorem}
\begin{proof}
  \leanok
  Apply the twisted affine vanishing \ref{pa-thm:hModule1_vanishing_qcoh} to the quasi-coherence
  packaging of \(\struct X(D)\) on \(U\) supplied by Theorem~\ref{pa-thm:divisorSheaf_qcohOn}.
\end{proof}

\section{Fibrewise large-twist vanishing}

We now fix a dominant affine morphism \(\pi \colon Y \to \mathbb P^1\) from the curve bundle \(Y\)
and its fibre divisor \(F = \mathrm{fiberWeilDivisor}(\pi)\) (Definition~\ref{pa-def:fiberWeilDivisor}),
with the affine two-cover \(V_0 = \pi^{-1}D_+(X_0)\), \(V_1 = \pi^{-1}D_+(X_1)\). The main theorem
computes the twisted cohomology \(H^1(\struct Y(D + nF))\) as a quotient of the fixed overlap
lattice by the fibre lattice, then runs the Noetherian stabilisation of the exhausting chain to
conclude vanishing for all large \(n\).

\begin{definition}[The fibre lattice inside the overlap, and the twisted \(H^1\)]
  \label{pa-def:fiberLatticeH1Equiv}
  \lean{AlgebraicGeometry.fiberLatticeOverlap, AlgebraicGeometry.mem_fiberLatticeOverlap,
    AlgebraicGeometry.fiberLatticeOverlap_mono, AlgebraicGeometry.iSup_fiberLatticeOverlap,
    AlgebraicGeometry.fiberLatticeH1Equiv}
  \leanok
  \dcref{custom}

  \uses{pa-def:fiberLattice, pa-thm:twoCoverH1LinearEquiv, pa-thm:subsingleton_hModule_divisorSheaf_one,
    pa-thm:divisorSections_twist_fixed, pa-def:divisorSheaf, pa-def:TwoCover_diff}
  Write \(N = \struct Y(D)(V_0 \cap V_1)\) for the fixed ambient overlap module and \(A_n\) for the
  denominator lattice of \(\struct Y(D + nF)\) (Definition~\ref{pa-def:fiberLattice}). Viewed inside
  \(N\) through the inclusion \(N \hookrightarrow K(Y)\), the lattices \(A_n\) form an increasing
  chain of submodules \(\overline A_n \le N\) exhausting \(N\), \(\bigsqcup_n \overline A_n = N\).
  Moreover the degree-one cohomology of the twisted sheaf is the quotient by this chain:
  \[
    H^1\bigl(\struct Y(D + nF)\bigr) \;\cong\; N \big/ \overline A_n \qquad (K\text{-linearly}).
  \]
\end{definition}
\begin{proof}
  \leanok
  Containment, monotonicity and exhaustion of \(\overline A_n\) inside \(N\) are those of \(A_n\)
  (Definition~\ref{pa-def:fiberLattice}) read through the inclusion of \(N\), the comap of a submodule
  along the subtype inclusion. For the isomorphism, the two-cover computation
  (Theorem~\ref{pa-thm:twoCoverH1LinearEquiv}) applies to \(\struct Y(D + nF)\) on the cover \(V_0,
  V_1\): its two affine vanishing hypotheses \(H'^1(V_i, \struct Y(D + nF)) = 0\) are supplied by
  Theorem~\ref{pa-thm:subsingleton_hModule_divisorSheaf_one}, giving
  \[
    H^1\bigl(\struct Y(D + nF)\bigr) \cong \struct Y(D + nF)(V_0 \cap V_1) \big/
      \mathrm{range}\,(\text{restriction-difference}) .
  \]
  The restriction-difference map (Definition~\ref{pa-def:TwoCover_diff}) reads on underlying rational
  functions as the literal difference in \(K(Y)\), so its range is exactly the join \(A_n = \struct
  Y(D + nF)(V_0) + \struct Y(D + nF)(V_1)\) inside the overlap module. Finally the overlap module is
  constant in \(n\): \(\struct Y(D + nF)(V_0 \cap V_1) = \struct Y(D)(V_0 \cap V_1) = N\)
  (Theorem~\ref{pa-thm:divisorSections_twist_fixed}). Substituting identifies the cokernel with \(N /
  \overline A_n\).
\end{proof}

\begin{theorem}[Noetherian stabilisation of an exhausting chain]
  \label{pa-thm:eventually_eq_top_of_monotone_of_iSup_eq_top}
  \lean{Submodule.eventually_eq_top_of_monotone_of_iSup_eq_top}
  \leanok
  \dcref{custom}

  Let \(M\) be a module over a ring \(R\) and \((A_n)_{n \in \N}\) an increasing chain of submodules
  with \(\bigsqcup_n A_n = M\). If the quotient \(M / A_0\) is Noetherian, then the chain is
  eventually everything: there is \(n_0\) with \(A_n = M\) for all \(n \ge n_0\).
\end{theorem}
\begin{proof}
  \leanok
  Push the chain into the Noetherian quotient \(M / A_0\): the images \(\overline{A_n} = A_n / A_0\)
  form an increasing chain of submodules of a Noetherian module, hence stabilise at some \(n_0\),
  \(\overline{A_n} = \overline{A_{n_0}}\) for \(n \ge n_0\). Their join is the image of \(\bigsqcup_n
  A_n = M\), namely all of \(M / A_0\); since the stabilised value dominates every term, it equals
  the whole quotient, \(\overline{A_{n_0}} = M / A_0\). Comapping back along the quotient map
  \(M \to M / A_0\) and using \(A_0 \le A_n\) gives \(A_n = M\) for all \(n \ge n_0\).
\end{proof}

\begin{theorem}[Fibrewise large-twist vanishing]
  \label{pa-thm:subsingleton_hModule_divisorSheaf_one_of_isDominant_toP1}
  \lean{AlgebraicGeometry.subsingleton_hModule_divisorSheaf_one_of_isDominant_toP1}
  \leanok
  \dcref{custom}

  \uses{pa-def:fiberLatticeH1Equiv, pa-thm:eventually_eq_top_of_monotone_of_iSup_eq_top, pa-def:fiberLattice,
    pa-thm:moduleFinite_divisorSheaf, pa-def:fiberWeilDivisor}
  Assume \(H^0(Y, \struct{Y/K})\) and \(H^1(Y, \struct{Y/K})\) are finite-dimensional over \(K\).
  Then for every Weil divisor \(D\) on \(Y\) there is \(n_0\) with
  \[
    H^1\bigl(\struct Y(D + nF)\bigr) = 0 \qquad \text{for all } n \ge n_0 ,
  \]
  where \(F\) is the fibre divisor of the dominant affine \(\pi \colon Y \to \mathbb P^1\). The
  bound \(n_0\) is per-class; the argument uses no duality and no higher direct image.
\end{theorem}
\begin{proof}
  \leanok
  Through the identification \(H^1(\struct Y(D + nF)) \cong N / \overline A_n\)
  (Definition~\ref{pa-def:fiberLatticeH1Equiv}) it suffices to show \(\overline A_n = N\) for all large
  \(n\). The chain \(\overline A_n\) is increasing and exhausts \(N\) (the exhaustion of
  Definition~\ref{pa-def:fiberLattice}, by pole clearing at the finite fibre over \([1 : 0]\)). Its
  base quotient \(N / \overline A_0 \cong H^1(\struct Y(D))\) is finite-dimensional: every divisor
  sheaf has finite \(H^1\) under the standing structure-sheaf finiteness hypotheses, by dévissage
  (Theorem~\ref{pa-thm:moduleFinite_divisorSheaf}). A finite-dimensional \(K\)-module is Noetherian, so
  the stabilisation Theorem~\ref{pa-thm:eventually_eq_top_of_monotone_of_iSup_eq_top} gives \(n_0\) with
  \(\overline A_n = N\) for \(n \ge n_0\); then \(N / \overline A_n = 0\) and the twisted \(H^1\)
  vanishes.
\end{proof}

\begin{theorem}[Fibrewise large-twist vanishing along a finite map]
  \label{pa-thm:subsingleton_hModule_divisorSheaf_one_of_isFinite_toP1}
  \lean{AlgebraicGeometry.subsingleton_hModule_divisorSheaf_one_of_isFinite_toP1}
  \leanok
  \dcref{custom}

  \uses{pa-thm:subsingleton_hModule_divisorSheaf_one_of_isDominant_toP1,
    pa-thm:moduleFinite_hModule_one_of_isFinite_toP1}
  Let \(\pi \colon Y \to \mathbb P^1\) be a finite dominant morphism compatible with the structure
  morphisms, and assume \(H^0(Y, \struct{Y/K})\) is finite-dimensional. Then for every Weil divisor
  \(D\) there is \(n_0\) with \(H^1(\struct Y(D + nF)) = 0\) for all \(n \ge n_0\).
\end{theorem}
\begin{proof}
  \leanok
  Finiteness of \(\pi\) makes \(H^1(Y, \struct{Y/K})\) finite-dimensional
  (Theorem~\ref{pa-thm:moduleFinite_hModule_one_of_isFinite_toP1}), discharging the remaining
  finiteness hypothesis of Theorem~\ref{pa-thm:subsingleton_hModule_divisorSheaf_one_of_isDominant_toP1}.
\end{proof}

\section{Effective witnesses and peeling}
\label{pa-sec:effective_peeling}

This section supplies the three remaining inputs of the uniform vanishing theorems: the
strict positivity of the fiber divisor's degree, the realisation of every class of large
degree by an \emph{effective} divisor, and the transport of \(H^1\)-vanishing up an
effective divisor. Throughout, \(Y\) is integral, smooth of relative dimension one, locally
of finite type and quasi-compact over a field \(K\), with finite-dimensional \(H^0\) and
\(H^1\) where degrees are used.

\begin{lemma}[Class arithmetic]
  \label{pa-lem:picClass_nsmul}
  \lean{AlgebraicGeometry.divisorBound_le_iff, AlgebraicGeometry.picClass_nsmul,
    AlgebraicGeometry.classDeg_pow}
  \leanok
  \dcref{custom}

  \uses{pa-def:curveDivisor_picClass, pa-thm:curveDivisor_picClass_hom, pa-def:classDeg,
    pa-thm:classDeg_mul, pa-lem:classDeg_zpow}
  For a Weil divisor \(D\) and \(n \in \N\): \(\struct Y(n D) = \struct Y(D)^{n}\), and for
  a Picard class \(L\): \(\deg_{\Pic}(L^{n}) = n \cdot \deg_{\Pic}(L)\). Moreover the
  divisor-bound comparison at a point is the pointwise coefficient comparison in \(\Z\).
\end{lemma}
\begin{proof}
  \leanok
  Induction on \(n\) from additivity of the class map
  (Theorem~\ref{pa-thm:curveDivisor_picClass_hom}) and of the degree
  (Theorem~\ref{pa-thm:classDeg_mul}); the bound comparison unfolds the definition of the
  valuation bound to its integer coefficient.
\end{proof}

\begin{lemma}[A point of \(\PP^1\) off the second chart]
  \label{pa-lem:P1_point_off_chart1}
  \lean{AlgebraicGeometry.exists_mem_chartOpen_zero_notMem_one}
  \leanok
  \dcref{custom}

  \uses{pa-def:P1_chartOpen, pa-def:P1_chartIota, pa-lem:P1_chartIota_range, pa-def:P1_chartCoord,
    pa-def:P1_awayAlgEquiv}
  For any field \(k\) there is a point of \(\PP^1_k\) lying in the chart \(D_+(X_0)\) but
  not in \(D_+(X_1)\) — the origin \([1 : 0]\).
\end{lemma}
\begin{proof}
  \leanok
  On the chart \(D_+(X_0) = \Spec k[X_1/X_0]\) (Definitions~\ref{pa-def:P1_chartOpen},
  \ref{pa-def:P1_awayAlgEquiv}), the preimage of \(D_+(X_1)\) is the basic open of the chart
  coordinate \(X_1/X_0\) (Definition~\ref{pa-def:P1_chartCoord}). The chart coordinate is not
  a unit (it corresponds to the polynomial variable), so the ideal it generates is proper
  and is contained in a maximal ideal \(M\). The image of \(M\) under the chart embedding
  (Definition~\ref{pa-def:P1_chartIota}, with range \(D_+(X_0)\) by
  Lemma~\ref{pa-lem:P1_chartIota_range}) lies in \(D_+(X_0)\) but not in the basic open of the
  coordinate, hence not in \(D_+(X_1)\).
\end{proof}

\begin{lemma}[A finite dominant morphism is surjective]
  \label{pa-lem:surjective_of_isFinite_of_isDominant}
  \lean{AlgebraicGeometry.surjective_of_isFinite_of_isDominant}
  \leanok
  \dcref{custom}

  A finite dominant morphism of schemes is surjective.
\end{lemma}
\begin{proof}
  \leanok
  Finite morphisms are closed maps, so the range is closed; dominance makes it dense; a
  closed dense subset is everything.
\end{proof}

\begin{theorem}[Strict positivity of the fiber divisor's degree]
  \label{pa-thm:zero_lt_deg_fiberWeilDivisor}
  \lean{AlgebraicGeometry.zero_lt_deg_fiberWeilDivisor}
  \leanok
  \dcref{custom}

  \uses{pa-lem:P1_point_off_chart1, pa-lem:surjective_of_isFinite_of_isDominant,
    pa-def:fiberWeilDivisor, pa-def:fiberCover}
  For a finite dominant \(\pi \colon Y \to \PP^1\), the fiber Weil divisor
  \(F = \mathrm{fiberWeilDivisor}(\pi)\) has strictly positive degree: \(0 < \deg F\).
\end{theorem}
\begin{proof}
  \leanok
  Choose a point \(z \in D_+(X_0) \setminus D_+(X_1)\) of \(\PP^1\)
  (Lemma~\ref{pa-lem:P1_point_off_chart1}). By surjectivity of the finite dominant \(\pi\)
  (Lemma~\ref{pa-lem:surjective_of_isFinite_of_isDominant}) there is \(x_0 \in Y\) over \(z\);
  it lies in the fiber chart \(V_0\) but not in \(V_1\), and it is not the generic point
  (the generic point maps into the overlap). Such a point is a pole witness: the fiber
  divisor has strictly positive coefficient at some point off \(V_1\), which forces
  \(0 < \deg F\) (the pole-existence estimate packaged with
  Definition~\ref{pa-def:fiberWeilDivisor}).
\end{proof}

\begin{lemma}[An effective witness of a class of large degree]
  \label{pa-lem:exists_effective_of_picClass}
  \lean{AlgebraicGeometry.exists_effective_of_picClass}
  \leanok
  \dcref{custom}

  \uses{pa-thm:riemann_inequality, pa-def:h0_h1_chi, pa-def:divisorSheaf, pa-def:HModule_linearEquiv0,
    pa-def:divisorSections, pa-def:scheme_ord, pa-def:divOf, pa-thm:curveDivisor_picClass_hom,
    pa-def:curveDivisor}
  If a Weil divisor \(W\) satisfies \(1 \le \deg W + \chi(\struct Y)\), then its Picard
  class is realised by an \emph{effective} divisor: there is \(E \ge 0\) with
  \(\struct Y(E) = \struct Y(W)\).
\end{lemma}
\begin{proof}
  \leanok
  By the Riemann inequality (Theorem~\ref{pa-thm:riemann_inequality}),
  \(h^0(\struct Y(W)) \ge \deg W + \chi(\struct Y) \ge 1\), so \(H^0(\struct Y(W))\) is
  nontrivial. A nonzero global section is, through the degree-zero identification
  (Definition~\ref{pa-def:HModule_linearEquiv0}) and the section description of the divisor
  sheaf (Definitions~\ref{pa-def:divisorSheaf}, \ref{pa-def:divisorSections}), a nonzero rational
  function \(f \in K(Y)^{\times}\) with \(\mathrm{ord}_x(f) \ge -\mathrm{coeff}_x(W)\) at
  every point (Definition~\ref{pa-def:scheme_ord}). Set \(E := W + \mathrm{div}(f)\)
  (Definition~\ref{pa-def:divOf}): the order bound says exactly \(E \ge 0\) coefficientwise,
  and principal divisors have trivial class
  (Theorem~\ref{pa-thm:curveDivisor_picClass_hom}), so \(\struct Y(E) = \struct Y(W)\).
\end{proof}

\begin{lemma}[Peeling one point]
  \label{pa-lem:peel_single}
  \lean{AlgebraicGeometry.peel_single, AlgebraicGeometry.peel_nsmul_single}
  \leanok
  \dcref{custom}

  \uses{pa-def:devissageSES, pa-thm:devissageSES_shortExact, pa-thm:skyModule_h1_vanishing,
    pa-thm:HModule_surjective_map_f, pa-lem:devissageDivisor_sub, pa-def:curveDivisor_single,
    pa-def:jumpModule}
  Let \(x\) be a closed point and \(B\) a Weil divisor with
  \(H^1(\struct Y(B)) = 0\) (i.e.\ a subsingleton). Then
  \(H^1(\struct Y(B + x)) = 0\), and by iteration
  \(H^1(\struct Y(B + n\,x)) = 0\) for all \(n \in \N\).
\end{lemma}
\begin{proof}
  \leanok
  Set \(D := B + x\); the d\'evissage divisor of \(D\) at \(x\) is \(B\)
  (Lemma~\ref{pa-lem:devissageDivisor_sub}). The d\'evissage short exact sequence
  (Definition~\ref{pa-def:devissageSES}, Theorem~\ref{pa-thm:devissageSES_shortExact})
  \(0 \to \struct Y(B) \to \struct Y(D) \to \mathrm{sky}_x \to 0\) has skyscraper cokernel,
  whose \(H^1\) vanishes (Theorem~\ref{pa-thm:skyModule_h1_vanishing}); the long exact
  sequence therefore gives a surjection \(H^1(\struct Y(B)) \twoheadrightarrow
  H^1(\struct Y(D))\) (Theorem~\ref{pa-thm:HModule_surjective_map_f}), and a quotient of a
  subsingleton is a subsingleton. Iterate for \(n\,x\).
\end{proof}

\begin{lemma}[Effective peeling]
  \label{pa-lem:peel_effective}
  \lean{AlgebraicGeometry.peel_effective}
  \leanok
  \dcref{custom}

  \uses{pa-lem:peel_single, pa-def:curveDivisor_single, pa-def:curveDivisor}
  Let \(A\) be a Weil divisor with \(H^1(\struct Y(A)) = 0\) and \(E \ge 0\) an effective
  divisor. Then \(H^1(\struct Y(A + E)) = 0\).
\end{lemma}
\begin{proof}
  \leanok
  Induction on the size of the (finite) support of \(E\). If the support is empty, \(E = 0\).
  Otherwise pick \(p\) in the support, with coefficient \(b > 0\); write
  \(E = b\,p + E'\) where \(E' = E - b\,p \ge 0\) has support one point smaller. By the
  induction hypothesis \(H^1(\struct Y(A + E')) = 0\), and peeling the multiplicity block
  \(b \, p\) (Lemma~\ref{pa-lem:peel_single}) gives \(H^1(\struct Y(A + E' + b\,p)) =
  H^1(\struct Y(A + E)) = 0\).
\end{proof}

\begin{theorem}[The class form of the fibrewise large-twist vanishing]
  \label{pa-thm:flv_class_form}
  \lean{AlgebraicGeometry.exists_subsingleton_hModule_one_of_one_le_classDeg_of_isFinite_toP1}
  \leanok
  \dcref{custom}

  \uses{pa-thm:subsingleton_hModule_divisorSheaf_one_of_isFinite_toP1,
    pa-thm:zero_lt_deg_fiberWeilDivisor, pa-lem:exists_effective_of_picClass, pa-lem:peel_effective,
    pa-thm:subsingleton_hModule_one_class, pa-lem:picClass_nsmul, pa-def:classDeg,
    pa-def:curveDivisor_picClass, pa-thm:classDeg_mul, pa-lem:classDeg_zpow, pa-def:fiberWeilDivisor}
  Let \(\pi \colon Y \to \PP^1\) be finite dominant, compatible with the structure
  morphisms, and \(l, \theta\) Picard classes with \(1 \le \deg_{\Pic}\theta\). Then there
  is \(n_0 \in \N\) such that for all \(n \ge n_0\) and \emph{every} Weil divisor \(D\) of
  class \(\struct Y(D) = l \cdot \theta^{n}\),
  \[
    H^1\bigl(\struct Y(D)\bigr) = 0 .
  \]
\end{theorem}
\begin{proof}
  \leanok
  Realise \(l\) by a divisor \(D_0\) (Definition~\ref{pa-def:curveDivisor_picClass}). Run the
  fibrewise vanishing at \(D_0\)
  (Theorem~\ref{pa-thm:subsingleton_hModule_divisorSheaf_one_of_isFinite_toP1}): there is
  \(m_0\) with \(H^1(\struct Y(D_0 + m\,F)) = 0\) for all \(m \ge m_0\), \(F\) the fiber
  divisor; fix \(m = m_0\). Let \(d' := \deg F > 0\)
  (Theorem~\ref{pa-thm:zero_lt_deg_fiberWeilDivisor}) and \(\varphi := \struct Y(F)\). Take
  \(n_0 := \max(0,\, m_0 d' + 1 - \chi(\struct Y))\).

  For \(n \ge n_0\), the residual class \(c := \theta^{n}\varphi^{-m_0}\) has degree
  \(n \deg\theta - m_0 d' \ge n - m_0 d' \ge 1 - \chi(\struct Y)\)
  (Lemma~\ref{pa-lem:picClass_nsmul}, Theorem~\ref{pa-thm:classDeg_mul},
  Lemma~\ref{pa-lem:classDeg_zpow}, using \(\deg\theta \ge 1\)). Realise \(c\) by a divisor
  \(W\); then \(1 \le \deg W + \chi(\struct Y)\), so \(c\) has an effective witness \(E\)
  (Lemma~\ref{pa-lem:exists_effective_of_picClass}). The divisor \(D_0 + m_0 F + E\) has class
  \(l \cdot \varphi^{m_0} \cdot \theta^{n}\varphi^{-m_0} = l\,\theta^{n}\); its \(H^1\)
  vanishes by effective peeling from the base vanishing at \(D_0 + m_0 F\)
  (Lemma~\ref{pa-lem:peel_effective}). Finally, witness-independence
  (Theorem~\ref{pa-thm:subsingleton_hModule_one_class}) transports the vanishing from this one
  witness to \emph{every} divisor \(D\) of the class \(l\,\theta^{n}\).
\end{proof}

\begin{theorem}[The rank anchor]
  \label{pa-thm:h0_eq_deg_add_chi}
  \lean{AlgebraicGeometry.h0_eq_deg_add_chi_of_subsingleton_hModule_one}
  \leanok
  \dcref{custom}

  \uses{pa-thm:chi_eq_h0, pa-thm:chi_divisorSheaf, pa-def:h0_h1_chi}
  For a Weil divisor \(D\) with \(H^1(\struct Y(D)) = 0\),
  \[
    h^0\bigl(\struct Y(D)\bigr) \;=\; \deg D + \chi(\struct Y) .
  \]
  (Riemann--Roch with vanishing \(h^1\); with \(\chi(\struct Y) = 1 - g\) this is the
  classical \(h^0(D) = \deg D + 1 - g\).)
\end{theorem}
\begin{proof}
  \leanok
  Under the vanishing, \(\chi(\struct Y(D)) = h^0(\struct Y(D))\)
  (Theorem~\ref{pa-thm:chi_eq_h0}), and the \(\chi\)-ledger reads
  \(\chi(\struct Y(D)) = \deg D + \chi(\struct Y)\)
  (Theorem~\ref{pa-thm:chi_divisorSheaf}).
\end{proof}

\section{The uniform vanishing bound}
\label{pa-sec:uniformVanishing}

\begin{theorem}[The uniform \(H^1\)-vanishing bound]
  \label{pa-thm:exists_bound_subsingleton_hModule_one}
  \lean{AlgebraicGeometry.exists_bound_subsingleton_hModule_one_of_isFinite_toP1}
  \leanok
  \dcref{custom}

  \uses{pa-thm:subsingleton_hModule_divisorSheaf_one_of_isFinite_toP1,
    pa-thm:zero_lt_deg_fiberWeilDivisor, pa-lem:exists_effective_of_picClass, pa-lem:peel_effective,
    pa-thm:subsingleton_hModule_one_class, pa-lem:picClass_nsmul, pa-def:classDeg,
    pa-def:curveDivisor_picClass, pa-thm:classDeg_mul, pa-lem:classDeg_zpow, pa-def:fiberWeilDivisor,
    pa-def:h0_h1_chi}
  Let \(\pi \colon Y \to \PP^1\) be finite dominant, compatible with the structure
  morphisms. There is a single threshold \(b \in \Z\), depending only on \((Y, \pi)\), such
  that for \emph{every} Weil divisor \(D\),
  \[
    b \le \deg D \implies H^1\bigl(\struct Y(D)\bigr) = 0 .
  \]
  Explicitly \(b = n_1 \deg F + 1 - \chi(\struct Y)\), where \(F\) is the fiber Weil
  divisor and \(n_1\) is the fibrewise-vanishing threshold at the divisor \(0\).
\end{theorem}
\begin{proof}
  \leanok
  Run the fibrewise vanishing at \(D_0 = 0\)
  (Theorem~\ref{pa-thm:subsingleton_hModule_divisorSheaf_one_of_isFinite_toP1}): there is
  \(n_1\) with \(H^1(\struct Y(n\,F)) = 0\) for all \(n \ge n_1\); fix the base twist
  \(n_1 F\). The threshold \(b := n_1 \deg F + 1 - \chi(\struct Y)\) is legal since
  \(0 < \deg F\) (Theorem~\ref{pa-thm:zero_lt_deg_fiberWeilDivisor}).

  Let \(D\) have \(\deg D \ge b\). The residual class
  \(c := \struct Y(D) \cdot \struct Y(F)^{-n_1}\) has degree
  \(\deg D - n_1\deg F \ge 1 - \chi(\struct Y)\) (Lemma~\ref{pa-lem:picClass_nsmul},
  Theorem~\ref{pa-thm:classDeg_mul}, Lemma~\ref{pa-lem:classDeg_zpow}); realise it by a divisor
  \(W\) (Definition~\ref{pa-def:curveDivisor_picClass}) and take an effective witness
  \(E \ge 0\) of its class (Lemma~\ref{pa-lem:exists_effective_of_picClass}). Then
  \(n_1 F + E\) is in the class of \(D\). Vanishing holds at \(n_1 F\) by the base choice;
  effective peeling (Lemma~\ref{pa-lem:peel_effective}) transports it to \(n_1 F + E\); and
  witness-independence (Theorem~\ref{pa-thm:subsingleton_hModule_one_class}) upgrades it to
  \(\struct Y(D)\) itself. Because the base twist is fixed once, the effective witness
  absorbs the whole degree slack and \(b\) depends only on \((Y, \pi)\).
\end{proof}

\section{The degree of the fiber twist}
\label{pa-sec:thetaDegree}

\begin{theorem}[The presentation divisor of the unit fiber divisor]
  \label{pa-thm:presentationDivisor_fiberDivisor_one}
  \lean{AlgebraicGeometry.presentationDivisor_presentation_fiberDivisor_one}
  \leanok
  \dcref{custom}

  \uses{pa-def:fiberDivisor, pa-def:fiberCoordUnit, pa-thm:fiberCoordUnit_order,
    pa-def:fiberWeilDivisor, pa-def:presentationDivisor, pa-lem:localEquations_presentation,
    pa-def:curveDivisor, pa-def:fiberCover}
  For a dominant \(\pi \colon Y \to \PP^1\), the presentation divisor of the meromorphic
  presentation of the unit fiber divisor system \(\mathrm{fiberDivisor}(\pi,1)\) is the
  fiber Weil divisor:
  \[
    \mathrm{presentationDivisor}\bigl(\mathrm{fiberDivisor}(\pi,1).\mathrm{presentation}\bigr)
      \;=\; F = \mathrm{fiberWeilDivisor}(\pi) .
  \]
\end{theorem}
\begin{proof}
  \leanok
  Compare coefficients point by point. The presentation
  (Lemma~\ref{pa-lem:localEquations_presentation}) of the pinned two-cover system \((t_0\)
  on \(V_0\), \(1\) on \(V_1)\) (Definition~\ref{pa-def:fiberDivisor}) has trivialising
  element the fiber coordinate unit \(u = t_0\) at points of \(V_0\)
  (Definition~\ref{pa-def:fiberCoordUnit}) and the constant \(1\) at points off \(V_0\)
  (which lie in \(V_1\), the two charts covering \(Y\)). At \(x \in V_0\) the coefficient
  of the presentation divisor (Definition~\ref{pa-def:presentationDivisor}) is
  \(\mathrm{ord}_x(u)\), which the fiber order table identifies with the coefficient of
  \(F = (\mathrm{div}\,u)^{+}\) at \(x\) (Theorem~\ref{pa-thm:fiberCoordUnit_order},
  Definition~\ref{pa-def:fiberWeilDivisor}); at \(x \notin V_0\) both coefficients are zero —
  the element is \(1\), and \(F\) is supported in \(V_0\)
  (Definition~\ref{pa-def:fiberWeilDivisor}).
\end{proof}

\begin{theorem}[\(\Theta\)-positivity of the unit fiber twist]
  \label{pa-thm:classDeg_fiberTwist_one}
  \lean{AlgebraicGeometry.fiberTwist_one_eq_picClass_fiberWeilDivisor,
    AlgebraicGeometry.classDeg_fiberTwist_one,
    AlgebraicGeometry.zero_lt_classDeg_fiberTwist_one,
    AlgebraicGeometry.one_le_classDeg_fiberTwist_one}
  \leanok
  \dcref{custom}

  \uses{pa-thm:presentationDivisor_fiberDivisor_one, pa-lem:localEquations_presentation,
    pa-def:curveDivisor_picClass, pa-def:classDeg, pa-def:fiberTwist,
    pa-thm:zero_lt_deg_fiberWeilDivisor, pa-thm:picClass_eq_of_presentationDivisor_eq}
  For a dominant \(\pi \colon Y \to \PP^1\), the unit fiber twist is the class of the fiber
  Weil divisor, \(\Theta_1 = \struct Y(F)\), and
  \[
    \deg_{\Pic}(\Theta_1) = \deg F .
  \]
  When \(\pi\) is moreover finite, the degree is strictly positive:
  \(1 \le \deg_{\Pic}(\Theta_1)\). (Only positivity is needed downstream; the identity
  \(\deg \Theta_1 = \deg \pi\) is never used.)
\end{theorem}
\begin{proof}
  \leanok
  \(\Theta_1\) is the Picard class of the unit fiber divisor system
  (Definition~\ref{pa-def:fiberTwist}); the class of a local-equation system is the class of
  the presentation divisor of its meromorphic presentation
  (Lemma~\ref{pa-lem:localEquations_presentation},
  Theorem~\ref{pa-thm:picClass_eq_of_presentationDivisor_eq} via the anchored class map of
  Definition~\ref{pa-def:curveDivisor_picClass}), which is \(F\)
  (Theorem~\ref{pa-thm:presentationDivisor_fiberDivisor_one}). Hence
  \(\Theta_1 = \struct Y(F)\), and the degree of an anchored class reads the divisor degree
  (Definition~\ref{pa-def:classDeg}). Positivity is
  Theorem~\ref{pa-thm:zero_lt_deg_fiberWeilDivisor}.
\end{proof}

\begin{theorem}[The positivity constant transports across base-field transitions]
  \label{pa-thm:one_le_classDeg_baseFieldTransition}
  \lean{AlgebraicGeometry.one_le_classDeg_cechPicMap_baseFieldTransition_of_one_le}
  \leanok
  \dcref{custom}

  \uses{pa-thm:classDeg_cechPicMap_baseFieldTransition}
  For the challenge curve \(C\) over \(k\), a base-field transition
  \(\varphi \colon K_1 \to K_2\) of field extensions of \(k\), and a Picard class \(L\) on
  \(C_{K_1}\) with \(1 \le \deg_{\Pic} L\): the pulled-back class on \(C_{K_2}\) retains the
  threshold, \(1 \le \deg_{\Pic}\bigl(\pi^{*}L\bigr)\), where \(\pi\) is the base-field
  transition of the curve.
\end{theorem}
\begin{proof}
  \leanok
  The degree is invariant under base-field transitions
  (Theorem~\ref{pa-thm:classDeg_cechPicMap_baseFieldTransition}); the threshold transfers
  verbatim.
\end{proof}

\section{Evaluation at the graph section}
\label{pa-sec:graphSectionEval}

For a field point \(t \colon \Spec K \to C\) of the challenge curve, the section
\(g_t\) of the projection \(C_K := (C \otimes \Spec K)_{\mathrm{left}} \to \Spec K\)
evaluates sections of any open containing the graph point into a copy of \(K\). This is
the residue half of the degree-one certificate.

\begin{definition}[The evaluation at the graph section]
  \label{pa-def:graphSectionEval}
  \lean{AlgebraicGeometry.Over.top_le_sectionOfPoint_preimage,
    AlgebraicGeometry.graphSectionEval, AlgebraicGeometry.graphSectionEval_res}
  \leanok
  \dcref{custom}

  \uses{pa-def:sectionOfPoint, pa-def:graphPoint}
  Let \(W\) be an open of \(C_K\) containing the graph point \(x_t\). The preimage of
  \(W\) under the section \(g_t\) is all of \(\Spec K\) (a one-point space whose point maps
  to \(x_t \in W\)), so pullback of sections along \(g_t\) defines a ring homomorphism
  \[
    \varepsilon_t \colon \Gamma(C_K, W) \longrightarrow \Gamma(\Spec K, \top),
  \]
  compatible with restriction of the open.
\end{definition}

\begin{theorem}[The evaluation splits the \(K\)-structure]
  \label{pa-thm:graphSectionEval_algebraMap}
  \lean{AlgebraicGeometry.graphSectionEval_algebraMap,
    AlgebraicGeometry.bijective_graphSectionEval_algebraMap,
    AlgebraicGeometry.isField_sections_top_overSpec,
    AlgebraicGeometry.surjective_graphSectionEval}
  \leanok
  \dcref{custom}

  \uses{pa-def:graphSectionEval, pa-def:overSpecGammaTop, pa-def:sectionOfPoint}
  The composite of the structure map \(K \to \Gamma(C_K, W)\) (sections carry the
  \(K\)-structure through the second projection) with the evaluation \(\varepsilon_t\) is
  the canonical identification \(K \cong \Gamma(\Spec K, \top)\)
  (Definition~\ref{pa-def:overSpecGammaTop}) — in particular it is bijective, the target is a
  field, and \(\varepsilon_t\) is surjective.
\end{theorem}
\begin{proof}
  \leanok
  The structure map of the sections is the pullback along the second projection
  \(\mathrm{snd}\), pre-composed with the identification; and
  \(g_t \circ \mathrm{snd} = \id\) (Definition~\ref{pa-def:sectionOfPoint}), so the composite
  pullback \(\varepsilon_t \circ \mathrm{snd}^{\sharp}\) is the pullback along the
  identity — the canonical identification, a ring isomorphism onto
  \(\Gamma(\Spec K, \top) \cong K\). Surjectivity of \(\varepsilon_t\) follows since
  already the composite is surjective.
\end{proof}

\begin{theorem}[The point--evaluation dictionary]
  \label{pa-thm:ker_graphSectionEval_eq_primeIdealOf}
  \lean{AlgebraicGeometry.ker_graphSectionEval_eq_primeIdealOf}
  \leanok
  \dcref{custom}

  \uses{pa-def:graphSectionEval, pa-thm:graphSectionEval_algebraMap, pa-def:graphPoint}
  On an \emph{affine} open \(W \ni x_t\), the kernel of the evaluation at the graph section
  is the prime ideal of the graph point:
  \[
    \ker \varepsilon_t \;=\; \mathfrak p_{x_t} \subseteq \Gamma(C_K, W).
  \]
\end{theorem}
\begin{proof}
  \leanok
  The stalk of \(C_K\) at \(x_t\) is the localisation of \(\Gamma(C_K, W)\) at
  \(\mathfrak p_{x_t}\), so a section \(f\) lies in \(\mathfrak p_{x_t}\) iff its germ at
  \(x_t\) is a non-unit. Chase germs through the section: the germ of
  \(\varepsilon_t(f)\) at the unique point of \(\Spec K\) is the stalk map of \(g_t\)
  applied to the germ of \(f\) at \(x_t\).

  If \(\varepsilon_t(f) = 0\) but the germ of \(f\) at \(x_t\) were a unit, its image under
  the stalk map would be a unit and also the germ of \(0\) — impossible. Conversely if
  \(f \in \mathfrak p_{x_t}\) but \(\varepsilon_t(f) \neq 0\), then \(\varepsilon_t(f)\) is
  a unit (the target is a field, Theorem~\ref{pa-thm:graphSectionEval_algebraMap}), so its
  germ is a unit; stalk maps are local ring homomorphisms, so reflecting units along the
  stalk map of \(g_t\) makes the germ of \(f\) at \(x_t\) a unit — contradicting
  \(f \in \mathfrak p_{x_t}\).
\end{proof}

\begin{theorem}[The quotient by the kernel is a copy of \(K\)]
  \label{pa-thm:finrank_quotient_ker_graphSectionEval}
  \lean{AlgebraicGeometry.finrank_quotient_ker_graphSectionEval,
    AlgebraicGeometry.isMaximal_ker_graphSectionEval}
  \leanok
  \dcref{custom}

  \uses{pa-def:graphSectionEval, pa-thm:graphSectionEval_algebraMap}
  For any open \(W \ni x_t\), the kernel of \(\varepsilon_t\) is a maximal ideal and
  \[
    \dim_K \bigl(\Gamma(C_K, W) \,/\, \ker\varepsilon_t\bigr) \;=\; 1 .
  \]
\end{theorem}
\begin{proof}
  \leanok
  \(\varepsilon_t\) is surjective onto a field
  (Theorem~\ref{pa-thm:graphSectionEval_algebraMap}), so the quotient by its kernel is
  isomorphic to that field — whence maximality — and as a \(K\)-vector space the field
  \(\Gamma(\Spec K, \top)\) is one-dimensional: the \(K\)-structure map onto it is
  bijective (Theorem~\ref{pa-thm:graphSectionEval_algebraMap} again).
\end{proof}

\section{The graph chart}
\label{pa-sec:graphChart}

The multiplicity half of the degree-one certificate: an affine chart of the fibre curve on
which the ideal of the graph section is principal on the transported point generator.

\begin{definition}[The graph chart and the graph equation]
  \label{pa-def:graphChart}
  \lean{AlgebraicGeometry.graphBaseChart, AlgebraicGeometry.isAffineOpen_graphBaseChart,
    AlgebraicGeometry.basePoint_mem_graphBaseChart,
    AlgebraicGeometry.top_le_preimage_graphBaseChart, AlgebraicGeometry.graphCoordEval,
    AlgebraicGeometry.graphGen, AlgebraicGeometry.graphElift, AlgebraicGeometry.graphChart,
    AlgebraicGeometry.isAffineOpen_graphChart, AlgebraicGeometry.graphChart_le_productChart,
    AlgebraicGeometry.graphChartEqn}
  \leanok
  \dcref{custom}

  \uses{pa-def:diagonalChartData, pa-def:productChart, pa-def:productChartSections,
    pa-thm:isLocalization_away_basicOpen_productChartSections, pa-def:pointGen_pointEv}
  Let \(t \colon \Spec K \to C\) be a field point with image point \(p_0\), and let
  \(U := \mathrm{chart}(p_0)\) be the frozen \'etale-coordinate chart of the curve at
  \(p_0\) (Definition~\ref{pa-def:diagonalChartData}), with coordinate \(u\), and
  \(F := \Gamma(\Spec K, \top)\). Define:
  \begin{itemize}
  \item the \emph{point pullback} \(c := t^{\sharp} \colon \Gamma(C, U) \to F\) (the whole
    of \(\Spec K\) maps into \(U\));
  \item the \emph{point generator} \(g := u \otimes 1 - 1 \otimes c(u) \in
    \Gamma(C,U) \otimes_k F\) (Definition~\ref{pa-def:pointGen_pointEv});
  \item the \emph{pushed idempotent lift} \(\tilde e_F\), the image of the frozen diagonal
    idempotent lift under \(1 \otimes c\);
  \item the \emph{graph chart} \(\mathfrak G := D(1 - \tilde e_F) \subseteq
    \mathfrak W(U, \top)\), the basic open of the product chart
    (Definitions~\ref{pa-def:productChart}, \ref{pa-def:productChartSections}) at the pushed
    complementary idempotent — an affine open of \(C_K\), and a localisation of
    \(\Gamma(C,U) \otimes_k F\) away from \(1 - \tilde e_F\)
    (Theorem~\ref{pa-thm:isLocalization_away_basicOpen_productChartSections});
  \item the \emph{graph equation} \(\gamma := g|_{\mathfrak G} \in
    \Gamma(C_K, \mathfrak G)\), the restriction of the identified point generator.
  \end{itemize}
\end{definition}

\begin{lemma}[The tensor-level evaluation and the graph point]
  \label{pa-lem:graphTensorEval}
  \lean{AlgebraicGeometry.graphTensorEval, AlgebraicGeometry.graphTensorEval_tmul,
    AlgebraicGeometry.graphTensorEval_map, AlgebraicGeometry.graphTensorEval_graphElift,
    AlgebraicGeometry.graphPoint_mem_productChart,
    AlgebraicGeometry.graphSectionEval_productChartSections}
  \leanok
  \dcref{custom}

  \uses{pa-def:graphChart, pa-thm:appLE_productChartSections, pa-def:graphSectionEval,
    pa-def:graphPoint, pa-def:diagonalChartData, pa-def:pointGen_pointEv}
  Write \(\mathrm{ev} \colon \Gamma(C,U) \otimes_k F \to F\),
  \(a \otimes \lambda \mapsto c(a)\lambda\), for the tensor-level point evaluation
  (Definition~\ref{pa-def:pointGen_pointEv}). Then:
  the graph point lies in the product chart \(\mathfrak W(U, \top)\);
  the evaluation of an identified tensor at the graph section is the tensor-level
  evaluation,
  \(\varepsilon_t\bigl(\mathrm{pcs}(y)\bigr) = \mathrm{ev}(y)\) for
  \(y \in \Gamma(C,U) \otimes_k F\); and the pushed idempotent lift evaluates to zero,
  \(\mathrm{ev}(\tilde e_F) = 0\).
\end{lemma}
\begin{proof}
  \leanok
  \emph{Membership}: the first projection of the graph point is the image point
  \(p_0 \in U\) (Definition~\ref{pa-def:graphPoint}), and the second-factor condition is
  vacuous at \(\top\). \emph{Evaluation}: on a pure tensor \(a \otimes b\), the lift-app
  rule for the section \(g_t = \langle t, \id\rangle\)
  (Theorem~\ref{pa-thm:appLE_productChartSections}) computes
  \(\varepsilon_t(\mathrm{pcs}(a \otimes b)) = t^{\sharp}(a) \cdot b = c(a)\,b =
  \mathrm{ev}(a \otimes b)\); extend additively. \emph{Idempotent}: the tensor evaluation
  of a second-factor push is the point pullback of the multiplication, and the frozen
  package satisfies \(\mathrm{mul}(\tilde e) = 0\)
  (Definition~\ref{pa-def:diagonalChartData}), so
  \(\mathrm{ev}(\tilde e_F) = c\bigl(\mathrm{mul}(\tilde e)\bigr) = 0\).
\end{proof}

\begin{theorem}[The graph point lies in the graph chart]
  \label{pa-thm:graphPoint_mem_graphChart}
  \lean{AlgebraicGeometry.graphPoint_mem_graphChart}
  \leanok
  \dcref{custom}

  \uses{pa-lem:graphTensorEval, pa-thm:ker_graphSectionEval_eq_primeIdealOf, pa-def:graphChart,
    pa-thm:graphSectionEval_algebraMap}
  \(x_t \in \mathfrak G\): the pushed complementary idempotent \(1 - \tilde e_F\) does not
  vanish at the graph point.
\end{theorem}
\begin{proof}
  \leanok
  Suppose it vanished at \(x_t\), i.e.\ its section were a non-unit in the stalk. On the
  affine product chart the point--evaluation dictionary
  (Theorem~\ref{pa-thm:ker_graphSectionEval_eq_primeIdealOf}) reads: the section lies in
  \(\ker\varepsilon_t\). But its evaluation is
  \(\mathrm{ev}(1 - \tilde e_F) = 1 - 0 = 1 \neq 0\) in the field \(F\)
  (Lemma~\ref{pa-lem:graphTensorEval}, Theorem~\ref{pa-thm:graphSectionEval_algebraMap}) — a
  contradiction.
\end{proof}

\begin{theorem}[The kernel presentation of the graph ideal]
  \label{pa-thm:ker_graphSectionEval_eq_span_graphChartEqn}
  \lean{AlgebraicGeometry.ker_graphSectionEval_eq_span_graphChartEqn}
  \leanok
  \dcref{custom}

  \uses{pa-def:graphChart, pa-lem:graphTensorEval, pa-thm:graphPoint_mem_graphChart,
    pa-thm:ker_pointEv_map_localization_eq, pa-def:diagonalChartData,
    pa-thm:isLocalization_away_basicOpen_productChartSections, pa-def:graphSectionEval}
  On the graph chart, the kernel of the evaluation at the graph section is the principal
  ideal on the graph equation:
  \[
    \ker\bigl(\varepsilon_t \colon \Gamma(C_K, \mathfrak G) \to F\bigr)
      \;=\; (\gamma).
  \]
\end{theorem}
\begin{proof}
  \leanok
  The chart sections are a localisation of the tensor ring \(\Gamma(C,U) \otimes_k F\)
  away from \(1 - \tilde e_F\)
  (Theorem~\ref{pa-thm:isLocalization_away_basicOpen_productChartSections}), and on the
  localisation the evaluation \(\varepsilon_t\) restricts, along the structure map, to the
  tensor-level evaluation \(\mathrm{ev}\) (Lemma~\ref{pa-lem:graphTensorEval}, compatible with
  the restriction to the basic open). The frozen diagonal package supplies the idempotent
  lift and the kernel presentation of the diagonal ideal
  (Definition~\ref{pa-def:diagonalChartData}), so the localised point-fiber keystone
  (Theorem~\ref{pa-thm:ker_pointEv_map_localization_eq}) computes the extension of
  \(\ker(\mathrm{ev})\) to the chart as the principal ideal on the image of the point
  generator — which is \(\gamma\) (Definition~\ref{pa-def:graphChart}).

  It remains to see \(\ker \varepsilon_t\) equals that extension. One inclusion is
  \(\varepsilon_t(\gamma) = \mathrm{ev}(g) = 0\). For the other, an element
  \(z \in \ker\varepsilon_t\) clears denominators: \(z \cdot s = \beta(b)\) with \(s\) a
  power of \(1 - \tilde e_F\) and \(b\) in the tensor ring; \(s\) evaluates to \(1\) and
  \(z\) to \(0\), so \(\mathrm{ev}(b) = 0\), and \(b\)'s image lies in the extension of
  \(\ker(\mathrm{ev})\); since \(s\) is invertible in the localisation, \(z\) does too.
\end{proof}

\begin{theorem}[The colength of the graph equation is one]
  \label{pa-thm:finrank_quotient_span_graphChartEqn}
  \lean{AlgebraicGeometry.finrank_quotient_span_graphChartEqn}
  \leanok
  \dcref{custom}

  \uses{pa-thm:ker_graphSectionEval_eq_span_graphChartEqn,
    pa-thm:finrank_quotient_ker_graphSectionEval}
  \[
    \dim_K \bigl(\Gamma(C_K, \mathfrak G) \,/\, (\gamma)\bigr) \;=\; 1 .
  \]
\end{theorem}
\begin{proof}
  \leanok
  \((\gamma) = \ker\varepsilon_t\)
  (Theorem~\ref{pa-thm:ker_graphSectionEval_eq_span_graphChartEqn}), and the quotient by the
  kernel is a copy of \(K\)
  (Theorem~\ref{pa-thm:finrank_quotient_ker_graphSectionEval}).
\end{proof}

\section{The degree-one certificate of the graph class}
\label{pa-sec:graphDegree}

\begin{lemma}[The graph equation is nonzero]
  \label{pa-lem:graphChartEqn_ne_zero}
  \lean{AlgebraicGeometry.graphGen_mem_nonZeroDivisors,
    AlgebraicGeometry.graphChartEqn_ne_zero}
  \leanok
  \dcref{custom}

  \uses{pa-thm:diagonal_regularity_engine, pa-def:graphChart, pa-thm:graphSectionEval_algebraMap,
    pa-thm:graphPoint_mem_graphChart, pa-def:diagonalChartData,
    pa-thm:isLocalization_away_basicOpen_productChartSections}
  The point generator \(g\) is a nonzerodivisor of the tensor ring
  \(\Gamma(C,U) \otimes_k F\), and the graph equation \(\gamma \neq 0\) in
  \(\Gamma(C_K, \mathfrak G)\).
\end{lemma}
\begin{proof}
  \leanok
  The regularity engine (Theorem~\ref{pa-thm:diagonal_regularity_engine}) makes
  \(u \otimes 1 - 1 \otimes b\) a nonzerodivisor of \(\Gamma(C,U) \otimes_k F\), uniformly
  in the second tensor factor and the element \(b\) — here \(b = c(u)\). If \(\gamma = 0\)
  in the localisation, then some power \(s\) of \(1 - \tilde e_F\) satisfies \(g \cdot s = 0\)
  in the tensor ring, so \(s = 0\) (nonzerodivisor); but then the localisation itself is
  the zero ring, contradicting that it evaluates onto the field \(F\)
  (Lemma~\ref{pa-lem:graphTensorEval}, Theorem~\ref{pa-thm:graphSectionEval_algebraMap},
  Theorem~\ref{pa-thm:graphPoint_mem_graphChart}), where \(1 \neq 0\).
\end{proof}

\begin{theorem}[The graph point is closed]
  \label{pa-thm:graphPoint_ne_genericPoint}
  \lean{AlgebraicGeometry.graphPoint_ne_genericPoint}
  \leanok
  \dcref{custom}

  \uses{pa-lem:graphChartEqn_ne_zero, pa-thm:ker_graphSectionEval_eq_span_graphChartEqn,
    pa-thm:ker_graphSectionEval_eq_primeIdealOf, pa-lem:fromSpec_base_ne_genericPoint,
    pa-lem:primeIdealOf_fromSpec_base}
  The graph point is not the generic point of the fibre curve:
  \(x_t \neq \eta\).
\end{theorem}
\begin{proof}
  \leanok
  On the affine graph chart, the prime of the graph point is
  \(\ker\varepsilon_t = (\gamma)\)
  (Theorems~\ref{pa-thm:ker_graphSectionEval_eq_primeIdealOf},
  \ref{pa-thm:ker_graphSectionEval_eq_span_graphChartEqn}), a \emph{nonzero} ideal since
  \(\gamma \neq 0\) (Lemma~\ref{pa-lem:graphChartEqn_ne_zero}). Points of nonzero primes of an
  affine open are not the generic point
  (Lemmas~\ref{pa-lem:fromSpec_base_ne_genericPoint},
  \ref{pa-lem:primeIdealOf_fromSpec_base}).
\end{proof}

\begin{lemma}[The graph equation is a unit away from the graph point]
  \label{pa-lem:isUnit_germ_graphChartEqn_of_ne}
  \lean{AlgebraicGeometry.isUnit_germ_graphChartEqn_of_ne}
  \leanok
  \dcref{custom}

  \uses{pa-thm:ker_graphSectionEval_eq_span_graphChartEqn,
    pa-thm:finrank_quotient_ker_graphSectionEval, pa-thm:ker_graphSectionEval_eq_primeIdealOf,
    pa-lem:primeIdealOf_fromSpec_base}
  For every point \(z \in \mathfrak G\) with \(z \neq x_t\), the germ of the graph equation
  \(\gamma\) at \(z\) is a unit.
\end{lemma}
\begin{proof}
  \leanok
  If the germ were a non-unit, \(\gamma\) would lie in the prime \(\mathfrak p_z\) of \(z\)
  on the chart, so \((\gamma) = \mathfrak p_{x_t} \subseteq \mathfrak p_z\)
  (Theorems~\ref{pa-thm:ker_graphSectionEval_eq_span_graphChartEqn},
  \ref{pa-thm:ker_graphSectionEval_eq_primeIdealOf}); but \(\mathfrak p_{x_t}\) is maximal
  (Theorem~\ref{pa-thm:finrank_quotient_ker_graphSectionEval}), forcing
  \(\mathfrak p_z = \mathfrak p_{x_t}\), hence \(z = x_t\) (points of an affine open are
  recovered from their primes, Lemma~\ref{pa-lem:primeIdealOf_fromSpec_base}) — contradiction.
\end{proof}

\begin{lemma}[The trivialising element away from the graph point is \(1\)]
  \label{pa-lem:presentationElem_graphLocalEquations_of_ne}
  \lean{AlgebraicGeometry.presentationElem_graphLocalEquations_of_ne}
  \leanok
  \dcref{custom}

  \uses{pa-def:graphLocalEquations, pa-def:graphPoint, pa-def:diagonalCover,
    pa-def:localEquations_pullback, pa-lem:localEquations_presentation}
  For \(x \neq x_t\), the trivialising element of the meromorphic presentation of the graph
  system at \(x\) is the constant \(1\).
\end{lemma}
\begin{proof}
  \leanok
  The graph system is the pullback of the diagonal system along the graph lift \(q\)
  (Definition~\ref{pa-def:graphLocalEquations}). At \(x \neq x_t\), the image \(q(x)\) misses
  the diagonal (Definition~\ref{pa-def:graphPoint}), so the diagonal equation at \(q(x)\) is
  the constant \(1\) (Definition~\ref{pa-def:diagonalCover}), whose pullback
  (Definition~\ref{pa-def:localEquations_pullback}) is \(q^{\sharp}(1) = 1\); the presentation
  element (Lemma~\ref{pa-lem:localEquations_presentation}) is its germ, i.e.\ \(1\).
\end{proof}

\begin{theorem}[The overlap identity at the graph point]
  \label{pa-thm:presentationElem_graphLocalEquations_val}
  \lean{AlgebraicGeometry.presentationElem_graphLocalEquations_val}
  \leanok
  \dcref{custom}

  \uses{pa-def:graphLocalEquations, pa-def:graphChart, pa-thm:appLE_productChartSections,
    pa-def:diagonalChartData, pa-def:diagonalCover, pa-def:diagonalChart,
    pa-def:localEquations_pullback, pa-lem:localEquations_presentation, pa-def:diagonalSetup,
    pa-def:graphPoint, pa-lem:pointedCover_genericPoint_mem, pa-thm:graphPoint_mem_graphChart,
    pa-def:productChartSections}
  The rational function trivialising the graph divisor at the graph point — the germ at the
  generic point \(\eta\) of the pulled diagonal equation — \emph{is} the germ at \(\eta\)
  of the graph equation:
  \[
    \mathrm{elem}_{x_t}\bigl(\Gamma_t\bigr)
      \;=\; \gamma_\eta \quad\text{in } K(C_K).
  \]
\end{theorem}
\begin{proof}
  \leanok
  Both sides restrict to the same section on the overlap
  \(W := (\text{pullback member at } x_t) \cap \mathfrak W(U, \top)\), an open containing
  \(\eta\); since germs at \(\eta\) are computed through any open containing it, the two
  germs agree.

  \emph{The left side on \(W\).} The pullback member at \(x_t\) carries the pulled
  diagonal equation \(q^{\sharp}(e_{p_0})\) (Definitions~\ref{pa-def:localEquations_pullback},
  \ref{pa-def:diagonalCover}), where the diagonal equation \(e_{p_0}\) on the diagonal chart
  is the identified diagonal generator \(u \otimes 1 - 1 \otimes u\)
  (Definitions~\ref{pa-def:diagonalChart}, \ref{pa-def:diagonalSetup}). The lift-app rule for
  \(q = \langle\mathrm{fst}, \mathrm{snd}\circ t\rangle\)
  (Theorem~\ref{pa-thm:appLE_productChartSections}) computes its restriction to \(W\) as
  \(\mathrm{fst}^{\sharp}(u) - (\mathrm{snd}\circ t)^{\sharp}(u)\). The argument is run at
  a generalised image point (\(\forall p,\ \mathrm{fst}(q(x_t)) = p \Rightarrow \dots\)),
  instantiated at \(p := t(\ast)\), so that the rewrite along
  \(\mathrm{fst}(q(x_t)) = p_0\) never crosses the statement.

  \emph{The right side on \(W\).} The graph equation is the identified point generator
  \(u \otimes 1 - 1 \otimes c(u)\) restricted to the graph chart
  (Definition~\ref{pa-def:graphChart}); by the product-chart section calculus
  (Definition~\ref{pa-def:productChartSections}), its restriction to \(W\) is
  \(\mathrm{fst}^{\sharp}(u) - \mathrm{snd}^{\sharp}(c(u))
    = \mathrm{fst}^{\sharp}(u) - (\mathrm{snd} \circ t)^{\sharp}(u)\),
  the same section. Finally \(\eta \in W\): \(W\) is a nonempty open of the irreducible
  space \(C_K\) (it contains \(x_t\):
  Lemma~\ref{pa-lem:pointedCover_genericPoint_mem} for the member,
  Theorem~\ref{pa-thm:graphPoint_mem_graphChart} for the chart side of the overlap).
\end{proof}

\begin{theorem}[The presentation divisor of the graph is a one-point divisor]
  \label{pa-thm:presentationDivisor_graphLocalEquations}
  \lean{AlgebraicGeometry.presentationDivisor_graphLocalEquations}
  \leanok
  \dcref{custom}

  \uses{pa-lem:presentationElem_graphLocalEquations_of_ne, pa-thm:graphPoint_ne_genericPoint,
    pa-def:presentationDivisor, pa-def:curveDivisor_single, pa-def:scheme_ordZ,
    pa-lem:localEquations_presentation, pa-def:graphLocalEquations}
  The presentation divisor of the graph system is supported at the graph point:
  \[
    \mathrm{presentationDivisor}\bigl(\Gamma_t\bigr)
      \;=\; \mathrm{ord}_{x_t}\bigl(\mathrm{elem}_{x_t}\bigr) \cdot x_t ,
  \]
  a single-point divisor (Definition~\ref{pa-def:curveDivisor_single}, legal since
  \(x_t \neq \eta\) by Theorem~\ref{pa-thm:graphPoint_ne_genericPoint}).
\end{theorem}
\begin{proof}
  \leanok
  Compare coefficients pointwise (Definition~\ref{pa-def:presentationDivisor}). At \(x_t\)
  both sides read \(\mathrm{ord}_{x_t}(\mathrm{elem}_{x_t})\)
  (Definition~\ref{pa-def:scheme_ordZ}). At \(x \neq x_t\) the trivialising element is \(1\)
  (Lemma~\ref{pa-lem:presentationElem_graphLocalEquations_of_ne}), of order zero, and the
  single-point divisor has coefficient zero there.
\end{proof}

\begin{theorem}[The graph class of a field point has degree one]
  \label{pa-thm:classDeg_graphPicClass}
  \lean{AlgebraicGeometry.classDeg_graphPicClass}
  \leanok
  \source{kleiman-picard}

  \uses{pa-thm:presentationDivisor_graphLocalEquations,
    pa-thm:presentationElem_graphLocalEquations_val,
    pa-thm:finrank_quotient_span_graphChartEqn, pa-lem:isUnit_germ_graphChartEqn_of_ne,
    pa-thm:finrank_quotient_span_section, pa-def:classDeg, pa-def:curveDivisor_picClass,
    pa-lem:curveDivisor_single_arith, pa-lem:localEquations_presentation,
    pa-thm:graphPoint_ne_genericPoint, pa-def:graphLocalEquations,
    pa-thm:picClass_eq_of_presentationDivisor_eq, pa-thm:graphPoint_mem_graphChart}
  \dcref{custom}
  For every field point \(t \colon \Spec K \to C\) of the challenge curve,
  \[
    \deg_{\Pic}\bigl(\struct{C_K}(\Gamma_t)\bigr) \;=\; 1 .
  \]
  This is the degree certificate of the Abel element: Kleiman's graph divisor
  \(\Gamma_g = (1 \times g)^{-1}\Delta\) is a relative effective divisor of fibre degree
  one.
\end{theorem}
\begin{proof}
  \leanok
  The class of the graph system is the anchored class of its presentation divisor
  (Lemma~\ref{pa-lem:localEquations_presentation},
  Theorem~\ref{pa-thm:picClass_eq_of_presentationDivisor_eq},
  Definition~\ref{pa-def:curveDivisor_picClass}), which is the one-point divisor
  \(m \cdot x_t\) with \(m := \mathrm{ord}_{x_t}(\mathrm{elem}_{x_t})\)
  (Theorem~\ref{pa-thm:presentationDivisor_graphLocalEquations}). The anchored degree of a
  one-point divisor is \(\deg(m \cdot x_t) = m \cdot [\kappa(x_t) : K]\)
  (Definition~\ref{pa-def:classDeg}, Lemma~\ref{pa-lem:curveDivisor_single_arith}).

  This product collapses in one shot through the chart colength dictionary
  (Theorem~\ref{pa-thm:finrank_quotient_span_section}) on the graph chart, applied to the
  rational function \(\mathrm{elem}_{x_t}\) with section representative \(\gamma\)
  (the overlap identity, Theorem~\ref{pa-thm:presentationElem_graphLocalEquations_val}) and
  the singleton point set \(\{x_t\}\): the hypotheses are that \(x_t\) lies in the chart
  (Theorem~\ref{pa-thm:graphPoint_mem_graphChart}) and that \(\gamma\) is a unit at every
  other closed point of the chart (Lemma~\ref{pa-lem:isUnit_germ_graphChartEqn_of_ne}); the
  dictionary then reads
  \[
    m \cdot [\kappa(x_t):K]
      \;=\; \dim_K\bigl(\Gamma(C_K, \mathfrak G)/(\gamma)\bigr)
      \;=\; 1
  \]
  by the colength computation
  (Theorem~\ref{pa-thm:finrank_quotient_span_graphChartEqn}).
\end{proof}
